% Leçon 18 — Grammaire : 被字句 — Chinois Tle A — Cours signé OnBuch+
% Programme officiel MINESEC. Compiler avec : tectonic ZH18-grammaire.tex
\newcommand{\DOCMATIERE}{Chinois}
\newcommand{\DOCNIVEAU}{Tle A}
\newcommand{\DOCMODULE}{Module 3 : Citoyenneté et ouverture au monde}
\newcommand{\DOCLECON}{ZH18}
\newcommand{\DOCTITRE}{Leçon 18 — Grammaire : 被字句}
% OnBuch+ — préambule « cours signés » ENRICHI (compilé avec tectonic / XeLaTeX)
% Système visuel « Pop » d'OnBuch+ : fond crème (jamais blanc), bordures encre
% épaisses, ombres « dures » décalées, titres Archivo Black en capitales.
% Version MATHÉMATIQUES : graphes (pgfplots/tikz), boîtes pédagogiques
% (définition, propriété, méthode, attention, le savais-tu, exercice résolu,
% exercice, TP, bilan), page de garde et signature OnBuch+ ancrée en bas.
%
% Macros attendues dans main.tex AVANT \input{preamble} :
%   \DOCMATIERE  \DOCNIVEAU  \DOCMODULE  \DOCLECON (numéro)  \DOCTITRE
\documentclass[11pt]{article}

\usepackage{fontspec}
\usepackage[french]{babel} % français = langue principale ; le chinois via \zh{..} / textebox (xeCJK)
\usepackage{xeCJK}
\usepackage{pifont} % \ding{51} \ding{55} utilisés par les modèles
\usepackage[a4paper,margin=2cm,top=3.4cm,bottom=2.8cm,headheight=1.4cm,footskip=1.3cm]{geometry}
\usepackage{amsmath,amssymb,amsfonts}
\usepackage{mathtools} % \xmapsto, \coloneqq... souvent utilisés par les modèles
\usepackage{cancel} % simplifications barrées (bilans de forces)
% \abs{z} (module, valeur absolue) et \norm{u} (norme), utilisés par les modèles
\providecommand{\abs}{}\let\abs\relax\DeclarePairedDelimiter{\abs}{\lvert}{\rvert}
\providecommand{\norm}{}\let\norm\relax\DeclarePairedDelimiter{\norm}{\lVert}{\rVert}
\usepackage[table]{xcolor} % [table] : fournit aussi \rowcolors (alternance de lignes)
\usepackage{pagecolor}
\usepackage{tikz}
\usetikzlibrary{arrows.meta,positioning,calc,shapes.geometric,shapes.misc,shapes.arrows,decorations.pathmorphing,decorations.pathreplacing,decorations.markings,patterns,shadows,mindmap,trees,fit,backgrounds,matrix,intersections,angles,quotes}
\usepackage{pgfplots}
\usepackage[version=4]{mhchem} % équations nucléaires \ce{^{235}_{92}U}
\usepackage[european,siunitx]{circuitikz} % schémas électriques
\pgfplotsset{compat=1.18}
\usepgfplotslibrary{fillbetween,groupplots} % aires entre courbes (calcul intégral)
\usepackage{enumitem}
\usepackage{fancyhdr}
% [most] charge toutes les bibliothèques tcolorbox (listings, minted, raster,
% documentation...) et gonfle inutilement la mémoire TeX sur un document long
% (~300 boîtes) : on ne charge que ce qui est réellement utilisé.
\usepackage[skins,breakable]{tcolorbox}
\usepackage{graphicx}
\usepackage{colortbl}
\usepackage{booktabs}
\usepackage{tabularx}
\usepackage{diagbox} % case d'en-tête diagonale des tableaux à double entrée
% Colonnes extensibles centrée / gauche / droite (C, L, R), souvent utilisées par les modèles
\newcolumntype{C}{>{\centering\arraybackslash}X}
\newcolumntype{L}{>{\raggedright\arraybackslash}X}
\newcolumntype{R}{>{\raggedleft\arraybackslash}X}
\usepackage{array}
\usepackage{multicol}
\usepackage{multirow}
\usepackage{siunitx}
% parse-numbers=false : accepte aussi \SI{2,5 \times 10^{-3}}{...}
\sisetup{locale=FR,output-decimal-marker={,},group-minimum-digits=4,parse-numbers=false}
\DeclareSIUnit\ohm{\ensuremath{\symup{\Omega}}} % U+2126 absent de la police
\DeclareSIUnit\division{div} % réglages d'oscilloscope (ms/div, V/div)
\DeclareSIUnit\tr{tr} % tours (vitesse de rotation, ex. \SI{1500}{\tr\per\minute})
\DeclareSIUnit\molar{M} % concentration molaire (ex. \SI{3}{\molar}, \SI{1}{\micro\molar})
\DeclareSIUnit\an{an} % année (le modèle confond parfois avec \year, un compteur TeX) — ex. \SI{5}{\kilo\meter\per\an}
\DeclareSIUnit\FCFA{FCFA} % franc CFA (ex. \SI{500}{\FCFA\per\kilo\gram})
\DeclareSIUnit\degreeBrix{\textdegree Brix} % teneur en sucre d'un sirop/jus (réfractométrie)
\DeclareSIUnit\month{mois} % le modèle utilise parfois \month, un compteur TeX, comme unité de durée
\usepackage{qrcode}
% \degree : commande fréquemment utilisée par les modèles pour un angle ou une
% température en dehors de \SI{}{} (non fournie par les paquets chargés).
\providecommand{\degree}{^{\circ}}
% \qed (fin de démonstration), utilisé par les modèles sans amsthm
\providecommand{\qed}{\hfill$\square$}
% \barre : double barre des valeurs interdites dans un tableau de variations (tkz-tab)
\providecommand{\barre}{\ensuremath{\|}}
% environnement proof (amsthm non chargé) utilisé par les modèles
\ifdefined\proof\else\newenvironment{proof}[1][Démonstration]{\par\noindent\textit{#1.}\ }{\qed\par}\fi
\usepackage{lastpage}
\usepackage{hyperref}
\hypersetup{hidelinks}

% ============================================================
% Palette « Pop » OnBuch+ — mêmes hex que la classe P (plus_ui.dart)
% ============================================================
\definecolor{poporange}{HTML}{F59321}
\definecolor{popdark}{HTML}{E07A0C}
\definecolor{popcream}{HTML}{FFF6E8}
\definecolor{popink}{HTML}{1C1714}
\definecolor{poppurple}{HTML}{7A5AE0}
\definecolor{popgreen}{HTML}{2E9E6A}
\definecolor{poppink}{HTML}{E0457B}
\definecolor{popblue}{HTML}{2F7DE1}
\definecolor{popgold}{HTML}{C98A12}
\definecolor{popmuted}{HTML}{8A7F76}
\definecolor{popline}{HTML}{EADFD2}
\definecolor{popgray}{HTML}{8A7F76}
\colorlet{popgrey}{popgray}
% Alias tolérants (le modèle nomme parfois les couleurs différemment)
\colorlet{popwhite}{white}
\colorlet{popblack}{popink}
\colorlet{popred}{poppink}
\colorlet{popyellow}{popgold}
% Teintes claires (fonds de tableaux/encadrés) — le modèle nomme parfois
% les couleurs avec un suffixe L (ex. popblueL) sans les avoir définies.
\colorlet{poporangeL}{poporange!20}
\colorlet{popdarkL}{popdark!20}
\colorlet{poppurpleL}{poppurple!20}
\colorlet{popgreenL}{popgreen!20}
\colorlet{poppinkL}{poppink!20}
\colorlet{popblueL}{popblue!20}
\colorlet{popgoldL}{popgold!20}
\colorlet{popredL}{popred!20}
\colorlet{popgrayL}{popgray!20}
% Teintes claires pour fonds de tableaux / zones
\colorlet{popblueL}{popblue!10!white}
\colorlet{poporangeL}{poporange!14!white}
\colorlet{popgreenL}{popgreen!12!white}
\colorlet{poppurpleL}{poppurple!12!white}
\colorlet{poppinkL}{poppink!12!white}
\colorlet{popgoldL}{popgold!14!white}

\pagecolor{popcream}

% ============================================================
% Typographie
% ============================================================
\defaultfontfeatures{Ligatures=TeX}
\setmainfont{PlusJakartaSans-Medium.ttf}[Path=../../../fonts/,BoldFont=PlusJakartaSans-Bold.ttf]
% Symboles, API et emojis absents de la police principale : repli sur DejaVu Sans (caractères actifs)
\newfontface\symfont{DejaVuSans.ttf}[Path=../../../fonts/]
\makeatletter
\newcommand\symactive[1]{\catcode#1=13 \begingroup\lccode`\~=#1 \lowercase{\endgroup\def~}{{\symfont\char#1}}}
\newcommand\symblank[1]{\catcode#1=13 \begingroup\lccode`\~=#1 \lowercase{\endgroup\def~}{}}
\newcommand\symhyph[1]{\catcode#1=13 \begingroup\lccode`\~=#1 \lowercase{\endgroup\def~}{\mbox{-}}}
\newcount\symc
\newcommand\symrange[2]{\symc=#1 \@whilenum\symc<#2 \do{\expandafter\symactive\expandafter{\the\symc}\advance\symc by 1 }}
\newcommand\symblankrange[2]{\symc=#1 \@whilenum\symc<#2 \do{\expandafter\symblank\expandafter{\the\symc}\advance\symc by 1 }}
\makeatother
\symrange{"0250}{"02FF}\symactive{"2713}\symactive{"2714}\symactive{"2717}\symactive{"2718}\symactive{"2610}\symactive{"2611}\symactive{"2612}
\symrange{"2600}{"27BF}
\catcode"2705=13 \begingroup\lccode`\~="2705 \lowercase{\endgroup\def~}{{\symfont\char"2713}}\symblank{"26BD}\symblank{"26A1}
\symrange{"2460}{"257F}\symactive{"223C}\symactive{"2248}\symactive{"2260}
\symhyph{"2011}\symhyph{"2E1A}\symblankrange{"1F300}{"1FAFF}\symblank{"FE0F}
% Caractères chinois : WenQuanYi Zen Hei (sous-ensemble GB2312 + pinyin), gras/italique simulés
\setCJKmainfont{WenQuanYiZenHei-subset.ttf}[Path=../../../fonts/,AutoFakeBold=3,AutoFakeSlant=0.2]
\xeCJKsetup{CJKspace=false}
% Pinyin : voyelles à caron (ǎ ǐ ǒ ǔ ǖ ǘ ǚ ǜ) absentes de la police principale -> caractères actifs
\newfontface\pyfont{WenQuanYiZenHei-subset.ttf}[Path=../../../fonts/]
\newcommand\pyactive[1]{\catcode#1=13 \begingroup\lccode`\~=#1 \lowercase{\endgroup\def~}{{\pyfont\char#1}}}
\pyactive{"01CD}\pyactive{"01CE}\pyactive{"01CF}\pyactive{"01D0}\pyactive{"01D1}\pyactive{"01D2}\pyactive{"01D3}\pyactive{"01D4}\pyactive{"01D5}\pyactive{"01D6}\pyactive{"01D7}\pyactive{"01D8}\pyactive{"01D9}\pyactive{"01DA}\pyactive{"01DB}\pyactive{"01DC}
% Maths en sans-serif (Fira Math) avec chiffres et lettres droites de Plus
% Jakarta, pour que les formules se fondent dans le texte.
\usepackage[math-style=ISO,bold-style=ISO]{unicode-math}
\setmathfont{FiraMath-Regular.otf}[Path=../../../fonts/]
\setmathfont{PlusJakartaSans-Medium.ttf}[Path=../../../fonts/,range={up/num,up/{Latin,latin},\mathup}]
\usepackage{icomma} % virgule décimale sans espace en mode math
% Caractères mathématiques Unicode absents des polices de texte (Plus Jakarta,
% Archivo Black) : rendus par leur équivalent mathématique, en texte comme en
% formule — sinon ils s'affichent en carré vide (ex. « Arithmétique dans ℤ »).
\usepackage{newunicodechar}
\newunicodechar{ℝ}{\ensuremath{\mathbb{R}}}
\newunicodechar{ℤ}{\ensuremath{\mathbb{Z}}}
\newunicodechar{ℂ}{\ensuremath{\mathbb{C}}}
\newunicodechar{ℕ}{\ensuremath{\mathbb{N}}}
\newunicodechar{ℚ}{\ensuremath{\mathbb{Q}}}
\newunicodechar{𝔻}{\ensuremath{\mathbb{D}}}
\newunicodechar{α}{\ensuremath{\alpha}}
\newunicodechar{β}{\ensuremath{\beta}}
\newunicodechar{γ}{\ensuremath{\gamma}}
\newunicodechar{δ}{\ensuremath{\delta}}
\newunicodechar{ε}{\ensuremath{\varepsilon}}
\newunicodechar{θ}{\ensuremath{\theta}}
\newunicodechar{λ}{\ensuremath{\lambda}}
\newunicodechar{μ}{\ensuremath{\mu}}
\newunicodechar{σ}{\ensuremath{\sigma}}
\newunicodechar{τ}{\ensuremath{\tau}}
\newunicodechar{φ}{\ensuremath{\varphi}}
\newunicodechar{ψ}{\ensuremath{\psi}}
\newunicodechar{ω}{\ensuremath{\omega}}
\newunicodechar{Σ}{\ensuremath{\Sigma}}
% majuscules produites par \MakeUppercase dans les titres (α-aminés -> Α)
\newunicodechar{Α}{\ensuremath{\alpha}}
\newunicodechar{Β}{\ensuremath{\beta}}
\newunicodechar{Γ}{\ensuremath{\Gamma}}
\newunicodechar{Δ}{\ensuremath{\Delta}}
\newunicodechar{Ε}{\ensuremath{\varepsilon}}
\newunicodechar{Θ}{\ensuremath{\Theta}}
\newunicodechar{Λ}{\ensuremath{\Lambda}}
\newunicodechar{Μ}{\ensuremath{\mu}}
\newunicodechar{Τ}{\ensuremath{\tau}}
\newunicodechar{Φ}{\ensuremath{\Phi}}
\newunicodechar{Ψ}{\ensuremath{\Psi}}
\newunicodechar{Ω}{\ensuremath{\Omega}}
\newunicodechar{↦}{\ensuremath{\mapsto}}
\newunicodechar{∈}{\ensuremath{\in}}
\newunicodechar{∉}{\ensuremath{\notin}}
\newunicodechar{∧}{\ensuremath{\wedge}}
\newunicodechar{⇔}{\ensuremath{\Leftrightarrow}}
\newunicodechar{⟺}{\ensuremath{\Longleftrightarrow}}
\newunicodechar{⇒}{\ensuremath{\Rightarrow}}
\newunicodechar{⟹}{\ensuremath{\Longrightarrow}}
\newunicodechar{∘}{\ensuremath{\circ}}
\newunicodechar{∪}{\ensuremath{\cup}}
\newunicodechar{∩}{\ensuremath{\cap}}
\newunicodechar{≡}{\ensuremath{\equiv}}
\newunicodechar{⊂}{\ensuremath{\subset}}
\newunicodechar{⊥}{\ensuremath{\perp}}
\newunicodechar{∃}{\ensuremath{\exists}}
\newunicodechar{∀}{\ensuremath{\forall}}
\newunicodechar{∛}{\ensuremath{\sqrt[3]{\,}}}
\newunicodechar{∜}{\ensuremath{\sqrt[4]{\,}}}
\newunicodechar{⃗}{\ensuremath{{}^{\to}}}
\newunicodechar{ⁿ}{\ifmmode^{n}\else\textsuperscript{\ensuremath{n}}\fi}
\newunicodechar{ˣ}{\ifmmode^{x}\else\textsuperscript{\ensuremath{x}}\fi}
\newunicodechar{ᵇ}{\ifmmode^{b}\else\textsuperscript{\ensuremath{b}}\fi}
\newunicodechar{ʳ}{\ifmmode^{r}\else\textsuperscript{\ensuremath{r}}\fi}
\newunicodechar{ᵘ}{\ifmmode^{u}\else\textsuperscript{\ensuremath{u}}\fi}
\newunicodechar{ʸ}{\ifmmode^{y}\else\textsuperscript{\ensuremath{y}}\fi}
\newunicodechar{ᵃ}{\ifmmode^{a}\else\textsuperscript{\ensuremath{a}}\fi}
\newunicodechar{ᵉ}{\ifmmode^{e}\else\textsuperscript{\ensuremath{e}}\fi}
\newunicodechar{ᵏ}{\ifmmode^{k}\else\textsuperscript{\ensuremath{k}}\fi}
\newunicodechar{ᵖ}{\ifmmode^{p}\else\textsuperscript{\ensuremath{p}}\fi}
\newunicodechar{ᴺ}{\ifmmode^{N}\else\textsuperscript{\ensuremath{N}}\fi}
\newunicodechar{ᶜ}{\ifmmode^{c}\else\textsuperscript{\ensuremath{c}}\fi}
\newunicodechar{ʰ}{\ifmmode^{h}\else\textsuperscript{\ensuremath{h}}\fi}
\newunicodechar{ᵠ}{\ifmmode^{q}\else\textsuperscript{\ensuremath{q}}\fi}
\newunicodechar{ₙ}{\ifmmode_{n}\else\textsubscript{\ensuremath{n}}\fi}
\newunicodechar{ₐ}{\ifmmode_{a}\else\textsubscript{\ensuremath{a}}\fi}
\newunicodechar{ᵢ}{\ifmmode_{i}\else\textsubscript{\ensuremath{i}}\fi}
\newunicodechar{ᵣ}{\ifmmode_{r}\else\textsubscript{\ensuremath{r}}\fi}
\newunicodechar{ₖ}{\ifmmode_{k}\else\textsubscript{\ensuremath{k}}\fi}
\newunicodechar{ᵦ}{\ifmmode_{\beta}\else\textsubscript{\ensuremath{\beta}}\fi}
\newunicodechar{ₓ}{\ifmmode_{x}\else\textsubscript{\ensuremath{x}}\fi}
\newfontfamily\popdisplay{ArchivoBlack-Regular.ttf}[Path=../../../fonts/]
\newfontfamily\poplogo{SpaceGrotesk-Bold.ttf}[Path=../../../fonts/]
\newfontfamily\popsemi{PlusJakartaSans-SemiBold.ttf}[Path=../../../fonts/]
\newfontfamily\popxbold{PlusJakartaSans-ExtraBold.ttf}[Path=../../../fonts/]
\newfontfamily\popxboldls{PlusJakartaSans-ExtraBold.ttf}[Path=../../../fonts/,LetterSpace=3.0]

\setlist[enumerate]{leftmargin=1.7em,itemsep=5pt,topsep=4pt}
\setlist[itemize]{leftmargin=1.7em,itemsep=4pt,topsep=4pt,label={\color{poporange}\textbullet}}
\setlength{\parindent}{0pt}
\setlength{\parskip}{5pt}
\renewcommand{\arraystretch}{1.25}

% pgfplots : style Pop par défaut
\pgfplotsset{
  every axis/.append style={axis line style={popink,line width=1pt},tick style={popink},
    grid style={popline,line width=0.5pt},label style={font=\small},tick label style={font=\footnotesize}},
  popcurve/.style={poporange,line width=1.8pt,no markers,smooth},
}
% Même style disponible dans un \draw TikZ simple (hors pgfplots)
\tikzset{popcurve/.style={poporange,line width=1.8pt}}
\tikzset{popgrid/.style={popline,very thin}}
\colorlet{popcurve}{poporange} % le nom du style est parfois utilisé comme couleur

% ============================================================
% Pill / squiggle
% ============================================================
\newcommand{\poppill}[3]{%
  \tikz[baseline=-0.55ex]\node[draw=popink,line width=1.2pt,rounded corners=2.6pt,%
    fill=#2,inner xsep=6.5pt,inner ysep=3pt]{%
    \popxboldls\fontsize{7}{7}\selectfont\color{#3}\MakeUppercase{#1}};%
}
\newcommand{\popsquiggle}[1]{%
  \tikz[baseline=-0.4ex]{
    \draw[#1,line width=2.6pt,line cap=round]
      plot [smooth] coordinates {(0,0.09) (0.34,0.01) (0.68,0.21) (1.02,0.01) (1.36,0.10)};
  }%
}
% Mot-clé surligné (marqueur orange sous le texte)
\newcommand{\cle}[1]{{\popxbold\color{popdark}#1}}

% ============================================================
% En-tête / pied
% ============================================================
\pagestyle{fancy}
\fancyhf{}
\renewcommand{\headrulewidth}{0pt}
\renewcommand{\footrulewidth}{0pt}
\lhead{\raisebox{-0.15ex}{\poplogo\fontsize{13}{13}\selectfont\color{popink}OnBuch\color{poporange}+}}
\rhead{\poppill{\DOCMATIERE\ \textbullet\ \DOCNIVEAU\ \textbullet\ Leçon \DOCLECON}{white}{popink}}
\lfoot{\popxbold\fontsize{7.6}{7.6}\selectfont\color{popmuted}\MakeUppercase{Cours signé OnBuch+}\ \textbullet\ {\popsemi\DOCTITRE}}
\rfoot{\tikz[baseline=-0.6ex]\node[rounded corners=8.5pt,fill=popink,minimum height=17pt,minimum width=17pt,inner xsep=5pt,inner sep=0]{\popxbold\fontsize{8}{8}\selectfont\color{popcream}\thepage\,/\,\pageref*{LastPage}};}

% ============================================================
% Titres
% ============================================================
\newcommand{\chaptitre}[2]{%
  \begin{tcolorbox}[enhanced,colback=poporange,colframe=popink,boxrule=2.6pt,arc=11pt,
      drop shadow=popink,left=15pt,right=15pt,top=12pt,bottom=11pt,before skip=3pt,after skip=18pt]
    \poppill{#2}{popink}{popcream}\par\vspace{9pt}
    {\raggedright\hyphenpenalty=10000\exhyphenpenalty=10000\popdisplay\fontsize{22}{24}\selectfont\color{popcream}\MakeUppercase{#1}\par}
  \end{tcolorbox}
}
% Section numérotée : pastille orange avec le numéro + titre Archivo
\newcounter{popsec}
\newcommand{\coursec}[1]{%
  \stepcounter{popsec}\setcounter{popsub}{0}%
  \par\vspace{16pt}\noindent
  \begin{minipage}{\linewidth}
  \tikz[baseline=-0.7ex]\node[draw=popink,line width=1.6pt,fill=poporange,rounded corners=4pt,minimum size=22pt,inner sep=0]{\popdisplay\fontsize{12}{12}\selectfont\color{popcream}\thepopsec};%
  \hspace{8pt}{\popdisplay\fontsize{15}{17}\selectfont\color{popink}\MakeUppercase{#1}}\par
  \vspace{4pt}\noindent{\color{poporange}\rule{36pt}{3.6pt}}
  \end{minipage}\par\nopagebreak\vspace{8pt}
}
\newcounter{popsub}
\newcommand{\courssub}[1]{%
  \stepcounter{popsub}%
  \par\vspace{9pt}\noindent{\popxbold\fontsize{11.5}{13}\selectfont\color{popink}{\color{poporange}\thepopsec.\thepopsub}\hspace{6pt}#1}\par\nopagebreak\vspace{4pt}
}

% ============================================================
% Boîtes pédagogiques — même recette : fond blanc, bordure épaisse colorée,
% ombre dure encre, badge plein. Titre optionnel [#1].
% ============================================================
\tcbset{popbase/.style={breakable,enhanced,colback=white,boxrule=2.3pt,arc=9pt,
  shadow={3pt}{-3pt}{0pt}{popink},left=14pt,right=14pt,top=11pt,bottom=11pt,before skip=11pt,after skip=15pt,
  fontupper=\popsemi\fontsize{10.6}{14.5}\selectfont\color{popink},
  fontlower=\popsemi\fontsize{10.6}{14.5}\selectfont\color{popink}}}
% Coche dessinée (le glyphe ✓ n'existe pas dans les polices du document)
\newcommand{\popcheck}[1]{\tikz[baseline=-0.5ex]\draw[#1,line width=1.6pt,line cap=round,line join=round](-0.13,0.0)--(-0.04,-0.09)--(0.13,0.1);}
\providecommand{\coche}{\popcheck{popgreen}}
\providecommand{\resultat}[1]{\par\smallskip\noindent\fcolorbox{popgreen}{popgreenL}{\parbox{\dimexpr\linewidth-2\fboxsep-2\fboxrule\relax}{\textbf{Résultat.} #1}}\par\smallskip}
\newcommand{\popboxhead}[3]{\poppill{#1}{#2}{white}\ifx\relax#3\relax\else\hspace{6pt}{\popxbold\fontsize{10.6}{12}\selectfont\color{popink}#3}\fi\par\vspace{7pt}}

\newtcolorbox{aretenir}[1][]{popbase,colframe=popgreen,before upper={\popboxhead{À retenir}{popgreen}{#1}}}
% Texte en chinois (document de lecture, dialogue, extrait) : cadre
% d'étude. Un titre optionnel (source, auteur) s'affiche dans l'en-tête.
\newtcolorbox{textebox}[1][]{popbase,colframe=popblue,before upper={\popboxhead{文本 · Texte}{popblue}{#1}},after upper={}}
% Marques ✓ / ✗ (forme correcte / fautive) souvent utilisées par les modèles
\providecommand{\cross}{\textcolor{poppink}{\ensuremath{\times}}}
\providecommand{\xmark}{\cross}\providecommand{\crossmark}{\cross}\providecommand{\cmark}{\ensuremath{\checkmark}}
% Ligne de réponse / justification à compléter (inventée par les modèles)
\providecommand{\justification}[1]{\par\noindent\textit{Justification~:} \dotfill\par}
% \textipa (package tipa non chargé) : la police ne couvre pas l'API -> simple passage
\providecommand{\textipa}[1]{#1}
% Soulignement (ulem non chargé) : \uline, \ul
\providecommand{\uline}[1]{\underline{#1}}\providecommand{\ul}[1]{\underline{#1}}
% \glossaire{\item ...} inventée par les modèles : liste de glossaire
\providecommand{\glossaire}[1]{\begin{itemize}#1\end{itemize}}
% \alert (beamer) inventée par les modèles : texte mis en évidence
\providecommand{\alert}[1]{\textbf{\textcolor{poppink}{#1}}}
% variantes ASCII d'ordinaux écrites en macro
\providecommand{\iere}{\textsuperscript{re}}\providecommand{\ieme}{\textsuperscript{e}}\providecommand{\ere}{\textsuperscript{re}}\providecommand{\eme}{\textsuperscript{e}}\providecommand{\er}{\textsuperscript{er}}\providecommand{\ie}{\textsuperscript{e}}
% Ordinaux français écrits en macro par les modèles : 1\ière, 3\ième
\newcommand{\ière}{\textsuperscript{re}}\newcommand{\ième}{\textsuperscript{e}}
% \exemple (inventée par les modèles) : étiquette « Exemple : »
\providecommand{\exemple}{\textbf{Exemple~:}\ }
% \square utilisable aussi hors mode math (cases à cocher)
\AtBeginDocument{\let\squareMath\square\renewcommand{\square}{\ensuremath{\squareMath}}}
% Mot ou phrase en chinois à l'intérieur d'un paragraphe français
\newcommand{\zh}[1]{\ifmmode\text{#1}\else{#1}\fi}
\let\eng\zh \let\esp\zh \let\all\zh % alias tolérés
% flèches d'intonation inventées par les modèles
\providecommand{\textsearrow}{}\renewcommand{\textsearrow}{\ensuremath{\searrow}}\providecommand{\textnearrow}{}\renewcommand{\textnearrow}{\ensuremath{\nearrow}}
% \fr{..} : mot français (inventé par les modèles)
\providecommand{\fr}[1]{\textit{#1}}
% zhbox : encadré de texte chinois inventé par les modèles
\newenvironment{zhbox}{\begin{quote}}{\end{quote}}
% \vs : « versus » inventé par les modèles
\providecommand{\vs}{\textit{vs}\ }
% \blank : case à compléter inventée par les modèles
\providecommand{\blank}[1][]{\underline{\hspace{1.6em}}}
\newtcolorbox{exemplebox}[1][]{popbase,colframe=poporange,before upper={\popboxhead{Exemple}{poporange}{#1}}}
\newtcolorbox{definition}[1][]{popbase,colframe=popblue,before upper={\popboxhead{Définition}{popblue}{#1}}}
\newtcolorbox{propriete}[1][]{popbase,colframe=poppurple,before upper={\popboxhead{Propriété}{poppurple}{#1}}}
\newtcolorbox{methode}[1][]{popbase,colframe=popgold,before upper={\popboxhead{Méthode}{popgold}{#1}}}
\newtcolorbox{attention}[1][]{popbase,colframe=poppink,before upper={\popboxhead{Attention}{poppink}{#1}}}
\newtcolorbox{experience}[1][]{popbase,colframe=popblue,colback=popblueL,before upper={\popboxhead{Expérience}{popblue}{#1}}}
% « Le savais-tu ? » : carte violette pleine (contexte, culture, Cameroun)
\newtcolorbox{savaistu}[1][]{popbase,colback=poppurple,colframe=popink,
  fontupper=\popsemi\fontsize{10.4}{14.2}\selectfont\color{white},
  before upper={\poppill{Le savais-tu\,?}{white}{poppurple}\ifx\relax#1\relax\else\hspace{6pt}{\popxbold\color{white}#1}\fi\par\vspace{7pt}}}
% Exercice résolu : énoncé en haut, \tcblower puis la solution sur fond crème
\newcounter{popexo}\newcounter{popexr}
\newtcolorbox{exoresolu}[1][]{popbase,colframe=poporange,colbacklower=popcream,
  segmentation style={popink,line width=1.2pt,dashed},
  before upper={\stepcounter{popexr}\popboxhead{Exercice résolu \thepopexr}{poporange}{#1}},
  before lower={\poppill{Solution}{popink}{popcream}\par\vspace{6pt}}}
% Exercice (non résolu en place) — la correction est dans la partie Corrigés
\newtcolorbox{exercice}[1][]{popbase,colframe=popink,
  before upper={\stepcounter{popexo}\popboxhead{Exercice \thepopexo}{popink}{#1}}}
% Corrigé
\newtcolorbox{corrige}[1][]{popbase,colframe=popgreen,colback=popgreenL,
  before upper={\popboxhead{Corrigé}{popgreen}{#1}}}
% Objectifs / prérequis / bilan
\newtcolorbox{objectifs}{popbase,colframe=popink,colback=poporangeL,
  before upper={\popboxhead{Objectifs de la leçon}{popink}{}}}
\newtcolorbox{prerequis}{popbase,colframe=popink,colback=popblueL,
  before upper={\popboxhead{Prérequis}{popblue}{}}}
\newtcolorbox{bilan}{popbase,colframe=popink,colback=white,boxrule=2.8pt,
  before upper={\popboxhead{Fiche bilan}{poporange}{L'essentiel à connaître pour l'examen}}}
% Figure encadrée avec légende
\newtcolorbox{popfigure}[1][]{enhanced,breakable=false,colback=white,colframe=popline,boxrule=1.4pt,arc=8pt,
  left=10pt,right=10pt,top=10pt,bottom=8pt,before skip=10pt,after skip=12pt,halign=center,
  lower separated=false,halign lower=center,
  fontlower=\popsemi\fontsize{9}{11.5}\selectfont\color{popmuted},
  #1}
\newcommand{\legende}[1]{\par\vspace{6pt}{\popsemi\fontsize{9}{11.5}\selectfont\color{popmuted}#1}}

% ============================================================
% Signature OnBuch+
% ============================================================
\newcommand{\poplogoinline}[1]{{\poplogo\fontsize{#1}{#1}\selectfont\color{popink}OnBuch\color{poporange}+}}
\newcommand{\signatureonbuch}{%
  \begin{tcolorbox}[enhanced,colback=popink,colframe=popink,boxrule=2.2pt,arc=10pt,
      drop shadow=poporange,left=14pt,right=14pt,top=10pt,bottom=10pt,before skip=0pt,after skip=0pt]
    \tikz[baseline=-0.7ex]\node[circle,fill=popgreen,draw=popcream,line width=1.3pt,minimum size=22pt,inner sep=0]{\popcheck{white}};%
    \hspace{9pt}\begin{minipage}[c]{0.62\linewidth}
      {\popxbold\fontsize{12}{13}\selectfont\color{popcream}Signé\ \textbullet\ {\poplogo OnBuch\color{poporange}+}}\par\vspace{2pt}
      {\raggedright\popsemi\fontsize{8}{10.5}\selectfont\color{popline}\MakeUppercase{Équipe pédagogique OnBuch+}\\\MakeUppercase{Conforme au programme officiel MINESEC}\par}
    \end{minipage}\hfill
    {\poplogo\fontsize{22}{22}\selectfont\color{popcream}OnBuch\color{poporange}+}
  \end{tcolorbox}%
}

% ============================================================
% Page de garde
% #1 titre · #2 matière · #3 niveau · #4 module · #5 n° leçon · #6 durée ·
% #7 description / objectifs (texte ou itemize)
% ============================================================
\newcommand{\pagegarde}[7]{%
  \thispagestyle{empty}
  \newgeometry{a4paper,margin=1.7cm,top=1.6cm,bottom=1.4cm}%
  \begin{tikzpicture}[remember picture,overlay]
    \fill[poporange!18] ([xshift=-3.2cm,yshift=-3.6cm]current page.north east) circle (4.6cm);
    \fill[poppurple!14] ([xshift=2.4cm,yshift=7.6cm]current page.south west) circle (3.8cm);
  \end{tikzpicture}%
  {\poplogo\fontsize{30}{30}\selectfont\color{popink}OnBuch\color{poporange}+}\hfill
  \raisebox{0.6ex}{\poppill{Cours signé \textbullet\ Premium}{popink}{popcream}}\par
  \vspace{1pt}\popsquiggle{poppurple}\par
  \vspace{8pt}
  \begin{tcolorbox}[enhanced,colback=poporange,colframe=popink,boxrule=2.8pt,arc=13pt,
      drop shadow=popink,left=18pt,right=18pt,top=12pt,bottom=13pt,after skip=10pt]
    \poppill{#2 \textbullet\ #3}{popink}{popcream}\hspace{5pt}\poppill{#4}{white}{popink}\par\vspace{8pt}
    {\popdisplay\fontsize{14}{14}\selectfont\color{popink}LEÇON #5}\par\vspace{4pt}
    {\raggedright\hyphenpenalty=10000\exhyphenpenalty=10000\popdisplay\fontsize{25}{27}\selectfont\color{popcream}\MakeUppercase{#1}\par}
    \vspace{8pt}
    {\popxbold\fontsize{9.5}{9.5}\selectfont\color{popink}\MakeUppercase{Durée conseillée : #6}}
  \end{tcolorbox}
  \begin{tcolorbox}[enhanced,colback=white,colframe=popink,boxrule=2.2pt,arc=10pt,
      drop shadow=popink,left=16pt,right=16pt,top=10pt,bottom=10pt,after skip=8pt]
    \poppill{Dans cette leçon}{poppurple}{white}\par\vspace{5pt}
    \popsemi\fontsize{9.3}{12.6}\selectfont\color{popink}#7
  \end{tcolorbox}
  \noindent\begin{minipage}[t]{0.487\linewidth}
    \begin{tcolorbox}[enhanced,colback=popgreen,colframe=popink,boxrule=2.2pt,arc=10pt,
        drop shadow=popink,left=11pt,right=11pt,top=10pt,bottom=10pt,height=3.1cm,valign=center]
      \tcbox[colback=white,colframe=popink,boxrule=1.3pt,arc=5pt,boxsep=0pt,nobeforeafter,
        left=3pt,right=3pt,top=3pt,bottom=3pt,valign=center]{\qrcode[height=1.9cm]{https://play.google.com/store/apps/details?id=com.luvvix.onbuch&hl=fr}}%
      \hspace{8pt}\begin{minipage}[c]{0.5\linewidth}
        {\popdisplay\fontsize{11}{12.5}\selectfont\color{white}\MakeUppercase{Télécharge OnBuch}}\par\vspace{4pt}
        {\raggedright\popsemi\fontsize{8.4}{11}\selectfont\color{white}Gratuit \textbullet\ Google Play\\Tuteur IA, annales, concours, résultats\par}
      \end{minipage}
    \end{tcolorbox}
  \end{minipage}\hfill
  \begin{minipage}[t]{0.487\linewidth}
    \begin{tcolorbox}[enhanced,colback=poppurple,colframe=popink,boxrule=2.2pt,arc=10pt,
        drop shadow=popink,left=11pt,right=11pt,top=10pt,bottom=10pt,height=3.1cm,valign=center]
      \tikz[baseline=-0.7ex]\node[circle,fill=white,draw=popink,line width=1.3pt,minimum size=1.6cm,inner sep=0]{\popdisplay\fontsize{24}{24}\selectfont\color{poppurple}+};%
      \hspace{8pt}\begin{minipage}[c]{0.6\linewidth}
        {\popdisplay\fontsize{11}{12.5}\selectfont\color{white}\MakeUppercase{Passe à OnBuch+}}\par\vspace{4pt}
        {\raggedright\popsemi\fontsize{8.4}{11}\selectfont\color{white}Tous les cours signés, Tuteur IA illimité, fiches PDF — onglet OnBuch+ de l'app\par}
      \end{minipage}
    \end{tcolorbox}
  \end{minipage}
  \vspace{8pt}
  \popcontenu
  \vfill
  \signatureonbuch
  \restoregeometry
  \newpage
}

% Rangée « Ce cours contient » (page de garde)
\newcommand{\poptile}[3]{% #1 couleur #2 gros texte #3 légende
  \begin{tcolorbox}[enhanced,colback=#1,colframe=popink,boxrule=2pt,arc=9pt,shadow={2.5pt}{-2.5pt}{0pt}{popink},
      left=6pt,right=6pt,top=9pt,bottom=9pt,halign=center,valign=center,height=1.75cm,width=\linewidth,nobeforeafter]
    {\popdisplay\fontsize{12.5}{13}\selectfont\color{white}\MakeUppercase{#2}}\par\vspace{3pt}
    {\centering\popsemi\fontsize{7.8}{9.5}\selectfont\color{white}#3\par}
  \end{tcolorbox}}
\newcommand{\popcontenu}{%
  \par\noindent{\popxboldls\fontsize{7.5}{7.5}\selectfont\color{popmuted}CE COURS SIGNÉ CONTIENT}\par\vspace{7pt}
  \noindent\begin{minipage}[t]{0.235\linewidth}\poptile{popblue}{Cours}{complet, illustré, pas à pas}\end{minipage}\hfill
  \begin{minipage}[t]{0.235\linewidth}\poptile{poporange}{Méthodes}{et exercices résolus}\end{minipage}\hfill
  \begin{minipage}[t]{0.235\linewidth}\poptile{poppink}{Exercices}{type examen + corrigés}\end{minipage}\hfill
  \begin{minipage}[t]{0.235\linewidth}\poptile{popgreen}{Fiche bilan}{l'essentiel à retenir}\end{minipage}\par}

% Sommaire
\newtcolorbox{sommaire}{enhanced,colback=white,colframe=popink,boxrule=2.2pt,arc=10pt,
  drop shadow=popink,left=18pt,right=18pt,top=13pt,bottom=13pt,before skip=6pt,after skip=20pt,
  before upper={\poppill{Sommaire}{poppurple}{white}\par\vspace{9pt}\popsemi\fontsize{10.6}{14.5}\selectfont\color{popink}}}

% Signature de fin de cours — poussée tout en bas de la dernière page
\newcommand{\findecours}{%
  \par\vfill\nopagebreak
  \begin{center}\popsquiggle{poporange}\end{center}\vspace{4pt}
  \signatureonbuch
}
\AtBeginDocument{\def\llbracket{\mathopen{[\mkern-3.5mu[}}\def\rrbracket{\mathclose{]\mkern-3.5mu]}}} % intervalles d'entiers (glyphe ⟦ absent de la police)
% Glyphes absents de Fira Math : redéfinis à partir de glyphes disponibles
\AtBeginDocument{%
  \def\setminus{\mathbin{\backslash}}\let\smallsetminus\setminus
  \def\vdots{\mathord{\vbox{\baselineskip3.2pt\lineskiplimit0pt\kern4pt\hbox{.}\hbox{.}\hbox{.}}}}%
  \def\ddots{\mathinner{\mkern1mu\raise6.5pt\hbox{.}\mkern2mu\raise3.6pt\hbox{.}\mkern2mu\raise0.7pt\hbox{.}\mkern1mu}}%
  \def\triangle{\mathord{\scalebox{0.9}{$\symup{\Delta}$}}}%
  \def\nabla{\mathord{\raisebox{\depth}{\scalebox{1}[-1]{$\symup{\Delta}$}}}}%
  \def\checkmark{\popcheck{popdark}}%
  \def\star{\mathbin{\ast}}\def\bigstar{\mathord{\ast}}%
  \def\diamond{\mathbin{\scalebox{0.6}{$\Diamond$}}}%
  \def\bigcup{\mathop{\mathchoice{\raisebox{-0.3em}{\scalebox{1.7}{$\cup$}}}{\scalebox{1.25}{$\cup$}}{\cup}{\cup}}\displaylimits}%
  \def\bigcap{\mathop{\mathchoice{\raisebox{-0.3em}{\scalebox{1.7}{$\cap$}}}{\scalebox{1.25}{$\cap$}}{\cap}{\cap}}\displaylimits}%
  \def\lesseqqgtr{\mathrel{\lessgtr}}\def\bigoplus{\mathop{\oplus}\displaylimits}%
}
\providecommand{\ssi}{si et seulement si} % abréviation fréquente
\AtBeginDocument{\providecommand{\wideparen}{\overparen}} % arc de cercle

%% FIGFIX-BEGIN v3 — correctifs globaux des figures (docs/onbuch-generation/tools/figfix)
\usepackage{adjustbox}\usepackage{etoolbox}
% toute figure TikZ/pgfplots plus large que la ligne est réduite à la largeur disponible
\BeforeBeginEnvironment{tikzpicture}{\begin{adjustbox}{max width=\linewidth}}
\AfterEndEnvironment{tikzpicture}{\end{adjustbox}}
% légendes pgfplots lisibles par-dessus les courbes
\pgfplotsset{every axis/.append style={legend style={fill=white,fill opacity=0.92,text opacity=1,draw=black!15,inner sep=3pt}}}
% étiquettes d'arêtes (syntaxe edge["..."]) avec fond blanc
\tikzset{every edge quotes/.append style={fill=white,inner sep=1.2pt}}
% bulles de cartes mentales : texte à la ligne dans la bulle, bulle un peu plus grande
\tikzset{every concept/.append style={text width=2.0cm,align=center,minimum size=2.3cm,inner sep=2pt}}
%% FIGFIX-END
\begin{document}

\pagegarde{Leçon 18 — Grammaire : 被字句}{Chinois}{Tle A}{Module 3 : Citoyenneté et ouverture au monde}{ZH18}{2 h}{Cette leçon de grammaire présente la phrase en 被 (被字句), structure passive fondamentale du chinois moderne. Les élèves apprennent à la construire, à l'employer correctement, à la comparer avec le passif français et à éviter les erreurs typiques des francophones.
\vspace{4pt}
\begin{multicols}{2}\small
\begin{itemize}
\item À la fin de cette leçon, je sais reconnaître une phrase en 被 dans un texte
\item À la fin de cette leçon, je sais construire le schéma : patient + 被 + agent + verbe + complément
\item À la fin de cette leçon, je sais employer 被 pour exprimer une action subie, souvent défavorable
\item À la fin de cette leçon, je sais placer correctement la négation 没(有) et les adverbes avant 被
\item À la fin de cette leçon, je sais utiliser les variantes 叫 et 让 dans la langue parlée
\item À la fin de cette leçon, je sais transformer une phrase active en phrase passive et inversement
\end{itemize}\end{multicols}}

\begin{sommaire}
\begin{multicols}{2}
\begin{enumerate}
\item Observation d'exemples
\item Schéma de la structure
\item Emplois et exceptions
\item Comparaison avec le français et pièges
\item Exercices d'application et de transformation
\item Activité d'intégration
\item Méthodes et exercices résolus
\item Exercices
\item Corrigés des exercices
\item Fiche bilan
\end{enumerate}
\end{multicols}
\end{sommaire}

\begin{objectifs}
\begin{itemize}
\item À la fin de cette leçon, je sais reconnaître une phrase en 被 dans un texte
\item À la fin de cette leçon, je sais construire le schéma : patient + 被 + agent + verbe + complément
\item À la fin de cette leçon, je sais employer 被 pour exprimer une action subie, souvent défavorable
\item À la fin de cette leçon, je sais placer correctement la négation 没(有) et les adverbes avant 被
\item À la fin de cette leçon, je sais utiliser les variantes 叫 et 让 dans la langue parlée
\item À la fin de cette leçon, je sais transformer une phrase active en phrase passive et inversement
\item À la fin de cette leçon, je sais éviter les erreurs fréquentes des francophones (verbe seul sans complément, mauvaise place de la négation)
\item À la fin de cette leçon, je sais traduire correctement des phrases passives du français vers le chinois
\end{itemize}
\end{objectifs}

\begin{prerequis}
\begin{itemize}
\item Ordre de base de la phrase chinoise : sujet + verbe + objet (SVO)
\item Compléments de résultat et de degré (吃完, 打破, 洗干净)
\item La particule 了 pour l'accomplissement
\item La négation 没(有) pour les actions passées
\item Les compléments de direction (进来, 出去)
\end{itemize}
\end{prerequis}

\newpage

\coursec{Leçon 18 — Grammaire : 被字句}

\courssub{Situation d'accroche : au marché et à l'école}

Au marché Mokolo à Yaoundé, un vendeur raconte à ses voisins : « Mon sac a été volé ! ». À l'école, un élève se plaint : « Mon cahier a été pris par mon voisin. ». Dans ces deux situations, la personne qui parle n'est pas l'auteur de l'action : c'est elle qui la \emph{subit}. En français, on utilise naturellement le passif (« a été volé », « a été pris ») ou la tournure « se faire + infinitif » (« je me suis fait voler mon sac »).

En chinois, il existe une structure spéciale et très fréquente pour exprimer cela : la \cle{phrase en 被} (bèi). Comment dire en chinois « mon téléphone a été cassé par mon petit frère » ? Où placer celui qui subit, celui qui agit, et le petit mot 被 ? C'est ce que nous allons découvrir dans cette leçon, en observant d'abord cinq exemples authentiques.

\textbf{Problématique :} dans une phrase en 被, qui est le sujet ? Où se place 被 ? Et comment passe-t-on d'une phrase active à une phrase passive ?

\courssub{Lecture guidée des cinq exemples}

Lisons attentivement les cinq phrases suivantes. Pour chacune, demande-toi : \emph{qui subit l'action ? qui la fait ?}

\begin{textebox}[Exemple 1 : le téléphone cassé]
我的手机被弟弟打破了。\\
Wǒ de shǒujī bèi dìdi dǎpò le.\\
« Mon téléphone a été cassé par mon petit frère. »
\end{textebox}

\begin{itemize}
\item \textbf{Exemple 2 :} \zh{钱包被偷了。} \emph{Qiánbāo bèi tōu le.} « Le portefeuille a été volé. »
\item \textbf{Exemple 3 :} \zh{他被老师批评了。} \emph{Tā bèi lǎoshī pīpíng le.} « Il a été critiqué par le professeur. »
\item \textbf{Exemple 4 :} \zh{晚饭被妈妈做好了。} \emph{Wǎnfàn bèi māma zuò hǎo le.} « Le dîner a été préparé par maman. »
\item \textbf{Exemple 5 :} \zh{我的自行车没被偷走。} \emph{Wǒ de zìxíngchē méi bèi tōu zǒu.} « Mon vélo n'a pas été volé. »
\end{itemize}

\courssub{Questions d'observation}

Avant d'énoncer la règle, observons ensemble. Reprends chaque exemple et réponds aux trois questions :

\begin{enumerate}
\item \textbf{Qui subit l'action ?} Dans l'exemple 1, c'est 我的手机 (\emph{wǒ de shǒujī}, « mon téléphone ») qui subit la casse. Dans l'exemple 3, c'est 他 (\emph{tā}, « il ») qui subit la critique.
\item \textbf{Qui fait l'action ?} Dans l'exemple 1, c'est 弟弟 (\emph{dìdi}, « le petit frère ») qui casse. Dans l'exemple 3, c'est 老师 (\emph{lǎoshī}, « le professeur ») qui critique. Dans l'exemple 2, l'auteur du vol n'est même pas mentionné !
\item \textbf{Où se place 被 ?} Toujours \emph{après} le sujet et \emph{avant} l'agent et le verbe : 手机 \textbf{被} 弟弟 打破了。
\end{enumerate}

\begin{methode}[Comment observer une phrase en 被]
Pour analyser n'importe quelle phrase en 被, procède en trois gestes :
\begin{enumerate}
\item \textbf{Souligne} le patient, c'est-à-dire celui qui subit l'action : c'est le \emph{sujet} de la phrase (placé avant 被).
\item \textbf{Encadre} 被 + l'agent (celui qui fait l'action), placé juste après 被.
\item \textbf{Entoure} le verbe et son complément (résultat, direction, 了…), placés à la fin.
\end{enumerate}
Application à l'exemple 1 : \underline{我的手机} [被 弟弟] (打破了)。
\end{methode}

\courssub{Le constat : le sujet est le patient}

\begin{definition}[La phrase en 被 (被字句)]
La \cle{被字句} (\emph{bèi zì jù}, littéralement « phrase à caractère 被 ») est une phrase dans laquelle le \textbf{sujet subit l'action} exprimée par le verbe. Elle est introduite par la préposition \zh{被} (\emph{bèi}), qui marque le passif et introduit l'agent (l'auteur de l'action).
\end{definition}

Le constat essentiel de cette observation est le suivant : dans une phrase en 被, \textbf{le sujet de la phrase est le patient} (celui qui subit), et non l'agent (celui qui agit). C'est exactement l'inverse de la phrase active ordinaire.

Comparons les deux versions du même événement :

\begin{center}
\begin{tabularx}{\linewidth}{|X|X|X|}
\hline
\rowcolor{popblueL} Type de phrase & Phrase chinoise + pinyin & Traduction \\
\hline
Active & 弟弟打破了我的手机。Dìdi dǎpò le wǒ de shǒujī. & Mon petit frère a cassé mon téléphone. \\
\hline
Passive (被) & 我的手机被弟弟打破了。Wǒ de shǒujī bèi dìdi dǎpò le. & Mon téléphone a été cassé par mon petit frère. \\
\hline
\end{tabularx}
\end{center}

Dans la phrase active, l'agent 弟弟 est sujet et le patient 手机 est complément d'objet. Dans la phrase passive, les rôles s'inversent : le patient 手机 devient sujet, et l'agent 弟弟 est introduit par 被.

\begin{popfigure}
\begin{tikzpicture}[
  bloc/.style={rectangle, rounded corners=3pt, draw=popline, fill=popcream, align=center, font=\small, inner sep=4pt},
  agent/.style={bloc, fill=poporangeL, draw=poporange},
  patient/.style={bloc, fill=popblueL, draw=popblue},
  verbe/.style={bloc, fill=popgreenL, draw=popgreen},
  bei/.style={bloc, fill=poppinkL, draw=poppink},
  fleche/.style={-{Stealth[length=2.5mm]}, thick, popdark}
]
% Phrase active (haut)
\node[agent, align=center] (a1) at (0,2.6){弟弟\\ agent};
\node[verbe, align=center] (v1) at (3.2,2.6){打破了\\ verbe + compl.};
\node[patient, align=center] (p1) at (6.6,2.6){我的手机\\ patient};
\node[font=\small\bfseries, popdark] at (-1.6,2.6) {Active};
% Phrase passive (bas)
\node[patient, align=center] (p2) at (0,0){我的手机\\ patient};
\node[bei, align=center] (b2) at (3.2,0){被\\ marque};
\node[agent, align=center] (a2) at (5.0,0){弟弟\\ agent};
\node[verbe, align=center] (v2) at (7.4,0){打破了\\ verbe + compl.};
\node[font=\small\bfseries, popdark] at (-1.6,0) {Passive};
% Flèches de déplacement
\draw[fleche, popblue] (p1.south) to[bend left=25] node[midway, right, font=\scriptsize, align=center]{le patient\\ devient sujet} (p2.north);
\draw[fleche, poporange] (a1.south) to[bend right=20] node[midway, left, font=\scriptsize, align=center]{l'agent passe\\ après 被} (a2.north west);
\end{tikzpicture}
\legende{Transformation de la phrase active en phrase passive en 被 : le patient devient sujet, l'agent est introduit par 被, le verbe (avec son complément) reste à la fin.}
\end{popfigure}

\courssub{Tableau récapitulatif des cinq exemples observés}

\begin{center}
\begin{tabularx}{\linewidth}{|X|X|X|X|}
\hline
\rowcolor{popblueL} Patient (sujet) & 被 + agent & Verbe + complément & Traduction \\
\hline
我的手机 & 被弟弟 & 打破了 & Mon téléphone a été cassé par mon petit frère. \\
\hline
钱包 & 被 (agent omis) & 偷了 & Le portefeuille a été volé. \\
\hline
他 & 被老师 & 批评了 & Il a été critiqué par le professeur. \\
\hline
晚饭 & 被妈妈 & 做好了 & Le dîner a été préparé par maman. \\
\hline
我的自行车 & 没被 (agent omis) & 偷走 & Mon vélo n'a pas été volé. \\
\hline
\end{tabularx}
\end{center}

Ce tableau révèle trois faits importants que nous étudierons en détail dans la suite de la leçon :

\begin{itemize}
\item L'\textbf{agent peut être omis} (exemples 2 et 5) : 被 est alors directement suivi du verbe. On dit 钱包被偷了 quand on ne sait pas (ou peu importe) qui a volé.
\item Le verbe est \textbf{presque toujours accompagné} d'un élément après lui : 了, un complément de résultat (破, 好) ou de direction (走). On ne dit pas \zh{手机被弟弟打} seul.
\item La \textbf{négation} se fait avec 没 placé \emph{avant} 被 : 没被偷走 (exemple 5).
\end{itemize}

\courssub{Schéma de la structure et emplois}

\begin{propriete}[Structure canonique de la phrase en 被]
\emph{Structure de base :}
\begin{center}
\textbf{Patient (Sujet) + 被 + Agent + Verbe + Complément (Résultat / Direction / 了)}
\end{center}

\emph{Variante sans agent (agent inconnu, indéfini ou évident) :}
\begin{center}
\textbf{Patient (Sujet) + 被 + Verbe + Complément}
\end{center}

\emph{Négation (passé / accompli) :}
\begin{center}
\textbf{Patient + 没 + 被 + (Agent) + Verbe + Complément}
\end{center}

\emph{Remarques importantes :}
\begin{enumerate}
\item \textbf{Verbes concernés :} Verbes transitifs d'action (动词) qui peuvent prendre un complément de résultat (结果补语), de direction (趋向补语) ou le marqueur 了. Exemples : 打破, 偷走, 做好, 批评, 吃完, 看见.
\item \textbf{Verbes exclus :} Verbes d'état, adjectifs predicatifs (sauf cas figés comme 被困), verbes intransitifs purs. On ne dit pas *\zh{他被高} (Il a été grand).
\item \textbf{Registre :} 被 est le marqueur passif \textbf{standard, neutre à formel} (écrit, discours officiel, plainte). À l'oral familier, on préfère souvent 叫, 让 ou 给 (ex. : \zh{手机让弟弟打破了}).
\end{enumerate}
\end{propriete}

\courssub{Comparaison avec le français et pièges}

\begin{center}
\begin{tabularx}{\linewidth}{|X|X|X|}
\hline
\rowcolor{popblueL} Notion & Français & Chinois \\
\hline
Marque du passif & auxiliaire « être » + participe passé (« a été cassé ») & préposition 被 devant l'agent, verbe inchangé \\
\hline
Introduction de l'agent & « par » + agent & 被 + agent \\
\hline
Place du patient & sujet du verbe passif & sujet, avant 被 \\
\hline
Accord & participe accordé (« cassé-e-s ») & aucun accord, le verbe ne change jamais \\
\hline
Négation passé & « n'a pas été » + participe & 没 + 被 + V + compl. (pas de 有) \\
\hline
\end{tabularx}
\end{center}

Bonne nouvelle pour les francophones : la logique générale est la même qu'en français (le patient devient sujet, l'agent est introduit par une préposition). La différence majeure est que le verbe chinois \textbf{ne se conjugue pas et ne s'accorde pas} : 打破 reste identique dans les deux phrases.

\begin{attention}[Pièges fréquents des francophones]
\begin{enumerate}
\item \textbf{Inversion sujet/agent :} On ne dit \emph{jamais} \zh{弟弟被我的手机打破了} (cela signifierait « le petit frère a été cassé par mon téléphone » !). L'ordre est rigide : \textbf{Patient + 被 + Agent}.
\item \textbf{Verbe nu sans complément :} Une phrase en 被 se termine rarement par un verbe seul. \textbf{Obligatoire} : 了, complément de résultat (破, 好, 完, 错), de direction (走, 来, 去) ou de quantité. \quad \emph{Correct :} \zh{信被他寄走了。} \quad \emph{Incorrect :} *\zh{信被他寄。}
\item \textbf{Négation avec 不 au lieu de 没 :} Le passif 被 rapporte un événement accompli (passé). On utilise \textbf{没(有)}. \zh{没被偷走} (pas *\zh{不被偷走}). \emph{Note :} 不被 s'emploie pour une habitude ou une vérité générale (ex. : \zh{这本书不被重视} « Ce livre n'est pas pris au sérieux »).
\item \textbf{Oubli de l'agent quand il est nécessaire :} Si l'agent est un pronom personnel (我, 你, 他) ou un nom précis essentiel au sens, on le garde. On l'omet seulement s'il est générique (« on », « quelqu'un ») ou inconnu.
\end{enumerate}
\end{attention}

\begin{exemplebox}[Cas particulier : 被 sans agent avec verbe à deux objets]
Quand le verbe a deux objets (donner, envoyer, apprendre), le passif peut porter sur l'un ou l'autre objet.
\begin{itemize}
\item \zh{书被他给了我。} (Le livre a été donné à moi par lui.) — Patient = 书.
\item \zh{我被他给了书。} (J'ai été donné un livre par lui.) — Patient = 我. \emph{Traduction française naturelle :} « Il m'a donné un livre. »
\end{itemize}
\end{exemplebox}

\courssub{Vocabulaire de la leçon}

\begin{center}
\begin{tabularx}{\linewidth}{|X|X|X|X|}
\hline
\rowcolor{popblueL} Caractères & Pinyin & Nature & Traduction \\
\hline
被 & bèi & prép. & marque du passif (par) \\
\hline
打破 & dǎpò & verbe (V+R) & casser (en frappant) \\
\hline
偷 & tōu & verbe & voler \\
\hline
偷走 & tōu zǒu & verbe (V+D) & voler (et emporter) \\
\hline
批评 & pīpíng & verbe / nom & critiquer / critique \\
\hline
做好 & zuò hǎo & verbe (V+R) & bien faire, préparer (repas), terminer \\
\hline
做完 & zuò wán & verbe (V+R) & finir de faire \\
\hline
晚饭 & wǎnfàn & nom & dîner, repas du soir \\
\hline
自行车 & zìxíngchē & nom & vélo \\
\hline
钱包 & qiánbāo & nom & portefeuille \\
\hline
患者 & huànzhě & nom & patient (celui qui subit, terme grammatical) \\
\hline
施事者 & shīshìzhě & nom & agent (celui qui agit, terme grammatical) \\
\hline
\end{tabularx}
\end{center}

\courssub{Étude de caractère : 被}

\begin{definition}[Analyse du caractère 被 (bèi)]
\begin{itemize}
\item \textbf{Radical (clé) :} \zh{衣} (yī, « vêtement », forme compressée \zh{衤} à gauche). Indique le domaine sémantique : tissu, couverture, vêtement.
\item \textbf{Partie phonétique :} \zh{皮} (pí, « peau, cuir »). Donne l'indication de prononciation (pí → bèi, alternance courante).
\item \textbf{Sens originel :} Une couverture, un drap (\zh{被子} bèizi).
\item \textbf{Glissement sémantique :} « Se couvrir de » $\rightarrow$ « Subir, recevoir (une action) » $\rightarrow$ Marque grammaticale du passif.
\item \textbf{Nombre de traits :} 10 traits.
\item \textbf{Ordre des traits :} 1. 丶 (point haut gauche de 衤) 2. ㇇ (horizontal court) 3. 丿 (pié gauche) 4. 丶 (point bas gauche) 5. 丨 (verticale centrale de 皮) 6. (trait brisé) (horizontal + crochet vertical) 7. 一 (horizontal) 8. 丿 (pié) 9. ㇏ (na) 10. 丶 (point final). \emph{Règle :} composant gauche (衤) avant droit (皮) ; haut avant bas ; horizontal avant vertical quand ils se croisent.
\end{itemize}
\end{definition}

\begin{savaistu}[D'où vient 被 ?]
À l'origine, le caractère 被 désignait une \textbf{couverture} (un drap dont on se couvre). Son sens a ensuite glissé vers « couvrir », puis « recevoir, subir » — on « se couvre » d'une action qu'on endure. C'est de ce sens « subir » que vient son emploi grammatical comme marque du passif. Cette histoire explique pourquoi la phrase en 被 exprime souvent (mais pas toujours) un événement \emph{subi}, parfois désagréable : être volé, cassé, critiqué… On retrouve 被 dans des mots comme 被子 (\emph{bèizi}, « la couverture ») et 被动 (\emph{bèidòng}, « passif », « passivité »).
\end{savaistu}

\courssub{Exercices d'application et de transformation}

\begin{exercice}[Transformation active $\rightarrow$ passive (被)]
Transforme les phrases actives suivantes en phrases passives avec 被. Attention au complément du verbe (了, résultat, direction) et à l'omission de l'agent si nécessaire.
\begin{enumerate}
\item 妈妈做好了晚饭。 \emph{(Maman a préparé le dîner.)}
\item 小偷偷走了我的自行车。 \emph{(Le voleur a volé mon vélo.)}
\item 老师批评了他。 \emph{(Le prof l'a critiqué.)}
\item 弟弟打破了我的手机。 \emph{(Mon petit frère a cassé mon tel.)}
\item 谁打扫了教室？ \emph{(Qui a nettoyé la salle de classe ? $\rightarrow$ La salle de classe a été nettoyée par qui ?)}
\end{enumerate}
\end{exercice}

\begin{exercice}[Négation et traduction]
\begin{enumerate}
\item Mets les phrases suivantes à la forme négative (passé/accompli) :
      \begin{itemize}
      \item 钱包被偷了。
      \item 作业被我做完了。
      \item 信被他寄走了。
      \end{itemize}
\item Traduis en chinois (caractères + pinyin) :
      \begin{itemize}
      \item Mon portable a été pris par ma sœur.
      \item Ce plat n'a pas été mangé.
      \item Les devoirs ont été finis par les élèves.
      \end{itemize}
\end{enumerate}
\end{exercice}

\begin{experience}[Jeu de rôle : Au commissariat / À la vie scolaire]
\textbf{Contexte :} Tu es au commissariat du quartier (ou chez le surveillant général) pour déclarer un vol ou un dégât.
\textbf{Rôles :} A = Déclarant (victime) ; B = Agent (note les faits).
\textbf{Consigne :} A doit décrire 3 objets/subis en utilisant la structure 被 (agent connu ou omis). B reformule pour confirmer : « Ah, donc [objet] a été [verbe] par [agent] ? ».
\textbf{Exemple :}
A: \zh{我的钱包被偷了。} (Mon portefeuille a été volé.)
B: \zh{你的钱包被偷了，对吗？} (Ton portefeuille a été volé, c'est ça ?)
\textbf{Vocabulaire utile :} 手机, 自行车, 书包, 作业, 电脑, 丢了 (perdu), 坏了 (cassé), 弄脏了 (salit).
\end{experience}

\begin{exoresolu}[Corrigé des exercices d'application]
\tcblower
\textbf{Exercice 1 (Transformation) :}
\begin{enumerate}
\item 晚饭被妈妈做好了。 (Wǎnfàn bèi māma zuò hǎo le.)
\item 我的自行车被小偷偷走了。 / 我的自行车被偷走了。 (Wǒ de zìxíngchē bèi xiǎotōu tōu zǒu le. / bèi tōu zǒu le.)
\item 他被老师批评了。 (Tā bèi lǎoshī pīpíng le.)
\item 我的手机被弟弟打破了。 (Wǒ de shǒujī bèi dìdi dǎpò le.)
\item 教室被谁打扫了？ (Jiàoshì bèi shéi dǎsǎo le?) \emph{Note :} L'agent « qui » (谁) se place après 被.
\end{enumerate}

\textbf{Exercice 2 (Négation) :}
\begin{itemize}
\item 钱包\emph{没}被偷了。 (Qiánbāo \emph{méi} bèi tōu le.)
\item 作业\emph{没}被我做完。 (Zuòyè \emph{méi} bèi wǒ zuò wán.) \emph{Attention :} pas de 了 à la fin à la forme négative accomplie.
\item 信\emph{没}被他寄走。 (Xìn \emph{méi} bèi tā jì zǒu.)
\end{itemize}

\textbf{Exercice 2 (Traduction) :}
\begin{itemize}
\item 我的手机被姐姐拿走了。 (Wǒ de shǒujī bèi jiějie ná zǒu le.) \emph{ou} 被姐姐拿了。
\item 这个菜没被吃。 (Zhège cài méi bèi chī.) \emph{ou} 这道菜没被吃掉。 (Zhè dào cài méi bèi chī diào.)
\item 作业被学生们做完了。 (Zuòyè bèi xuéshēngmen zuò wán le.)
\end{itemize}
\end{exoresolu}

\begin{aretenir}[L'essentiel de la leçon 18 : 被字句]
\begin{itemize}
\item \textbf{Fonction :} Exprimer le passif : le sujet est le \emph{patient} (celui qui subit).
\item \textbf{Structure de base :} \textbf{Patient + 被 + Agent + Verbe + Complément (résultat / direction / 了)}.
\item \textbf{Agent omis :} \textbf{Patient + 被 + Verbe + Complément} (agent inconnu, générique ou inutile).
\item \textbf{Négation (passé) :} \textbf{Patient + 没 + 被 + (Agent) + Verbe + Complément (sans 了)}.
\item \textbf{Contraintes verbales :} Verbe transitif + complément obligatoire (pas de verbe nu).
\item \textbf{Registre :} Neutre / Formel. Oral familier $\rightarrow$ 叫 / 让 / 给.
\item \textbf{Erreurs à bannir :} Inverser Patient/Agent ; oublier le complément verbal ; utiliser 不 au lieu de 没 pour un fait accompli.
\end{itemize}
\end{aretenir}

\coursec{Leçon 18 — Grammaire : 被字句}

\courssub{Situation d'observation : Une mésaventure au marché Central}

\begin{textebox}[Dialogue : Au marché Central de Yaoundé]
\textbf{Contexte :} Marie et Paul discutent après les cours. Marie a eu un problème.
\zh{玛丽：哎，保罗，我的手机被偷了！}\\
\emph{Mǎlì : Āi, Bǎoluó, wǒ de shǒujī bèi tōu le!}\\
« Eh, Paul, mon portable a été volé ! »\\[4pt]
\zh{保罗：真的？怎么被偷的？}\\
\emph{Bǎoluó : Zhēn de? Zěnme bèi tōu de?}\\
« Vraiment ? Comment a-t-il été volé ? »\\[4pt]
\zh{玛丽：我放在桌子上，转身买 ndolé，回来就不见了。}\\
\emph{Mǎlì : Wǒ fàng zài zhuōzi shàng, zhuǎnshēn mǎi ndolé, huílái jiù bùjiàn le.}\\
« Je l'avais posé sur la table, je me suis retournée pour acheter du ndolé, en revenant il avait disparu. »\\[4pt]
\zh{保罗：报警了吗？}\\
\emph{Bǎoluó : Bàojǐng le ma?}\\
« As-tu porté plainte ? »\\[4pt]
\zh{玛丽：已经报了，警察说会帮我找。手机被找回来的希望不大。}\\
\emph{Mǎlì : Yǐjīng bào le, jǐngchá shuō huì bāng wǒ zhǎo. Shǒujī bèi zhǎo huílái de xīwàng bù dà.}\\
« Déjà fait, la police dit qu'elle m'aidera à chercher. L'espoir que le portable soit retrouvé est mince. »
\end{textebox}

\begin{aretenir}[Observations guidées]
\begin{enumerate}
\item Repérez les trois phrases contenant \zh{被} : \zh{手机被偷了}, \zh{怎么被偷的}, \zh{手机被找回来}.
\item Dans chaque cas, qui est le \textbf{sujet grammatical} (en tête de phrase) ? Est-ce celui qui \emph{fait} l'action ou celui qui la \emph{subit} ?
\item Où se place l'agent (le voleur, la police) ? Est-il toujours exprimé ?
\item Que remarquez-vous sur la forme du verbe qui suit \zh{被} ? Est-il nu (\zh{偷}) ou accompagné (\zh{偷了}, \zh{找回来}) ?
\end{enumerate}
\end{aretenir}

\courssub{Schéma de la structure}

\begin{definition}[La phrase en 被 (bèizìjù)]
La phrase en \zh{被} (chinois : \zh{被字句}, \emph{bèizìjù}, litt. « phrase à caractère 被 ») est la construction passive standard du chinois moderne. Elle place en position de sujet grammatical le \textbf{patient} (celui qui subit l'action) et introduit l'\textbf{agent} (celui qui fait l'action) à l'aide de la particule \zh{被} (\emph{bèi}).
\end{definition}

\begin{aretenir}[Schéma de base de la 被字句]
\begin{center}
\textbf{Patient + 被 + (Agent) + Verbe + Complément obligatoire + (了)}
\end{center}
\begin{itemize}
\item \textbf{Patient} : Sujet grammatical, défini/connu, placé en tête.
\item \zh{被} : Particule passive, invariable.
\item \textbf{Agent} : Entre \zh{被} et le verbe. \textbf{Omissible} s'il est inconnu, évident ou sans importance.
\item \textbf{Verbe + Complément} : \textbf{Règle d'or} : le verbe ne reste \textbf{jamais nu}. Il doit être suivi d'un complément de résultat, de direction, de \zh{了}, d'un objet, etc.
\end{itemize}
\end{aretenir}

\begin{popfigure}
\begin{tikzpicture}[node distance=0.5cm, every node/.style={font=\small}]
\node[draw=popblue, fill=popblueL, rounded corners, minimum width=2.4cm, minimum height=1cm, align=center] (pat) {\textbf{Patient}\\ 手机 \emph{shǒujī}};
\node[draw=poporange, fill=poporangeL, rounded corners, minimum width=2.8cm, minimum height=1cm, align=center, right=of pat] (bei) {\textbf{被 + Agent}\\ 被小偷 \emph{bèi xiǎotōu}};
\node[draw=popgreen, fill=popgreenL, rounded corners, minimum width=3.2cm, minimum height=1cm, align=center, right=of bei] (verb) {\textbf{Verbe + Compl.}\\ 偷走了 \emph{tōu zǒu le}};
\draw[-{Stealth}, thick, popdark] (pat) -- (bei);
\draw[-{Stealth}, thick, popdark] (bei) -- (verb);
\draw[-{Stealth}, very thick, poppink, bend right=40] (bei.south) to node[below, align=center, text=poppink]{action subie} (pat.south);
\end{tikzpicture}
\legende{Analyse en blocs : 手机被小偷偷走了。\emph{Shǒujī bèi xiǎotōu tōu zǒu le.} — Le portable a été emporté par le voleur.}
\end{popfigure}

\courssub{Le patient : défini et topicalisé}

Le patient en tête de phrase en \zh{被} est presque toujours un \textbf{thème connu} (défini, possessif, démonstratif). Le chinois, comme le français, répugne à mettre un indéfini (\zh{一本书}) en sujet de passif.

\begin{exemplebox}[Patient défini (correct)]
\zh{那本书被小明借走了。}\\
\emph{Nà běn shū bèi Xiǎomíng jiè zǒu le.}\\
« Ce livre-là a été emprunté par Xiaoming. »\\[4pt]
\zh{我的作业被老师批改了。}\\
\emph{Wǒ de zuòyè bèi lǎoshī pīgǎi le.}\\
« Mes devoirs ont été corrigés par le prof. »
\end{exemplebox}

\begin{attention}[Patient indéfini : éviter le passif]
\begin{center}
\begin{tabularx}{\linewidth}{|X|X|}
\hline
\rowcolor{popblueL} \textbf{Maladroit / Faux} & \textbf{Préféré (Actif / Existence)} \\
\hline
*\zh{一本书被借走了。} & \zh{有人借走了一本书。} (\emph{Yǒurén jiè zǒu le yī běn shū.}) \\
*\zh{一个包被偷了。} & \zh{偷走一个包。} (\emph{Tōu zǒu le yī gè bāo.}) ou \zh{包被偷了} (si \zh{包} est contextuellement défini). \\
\hline
\end{tabularx}
\end{center}
\cle{Règle} : Tête de phrase en \zh{被} = information donnée (connue des deux interlocuteurs).
\end{attention}

\courssub{L'agent : exprimé ou omis}

\begin{propriete}[Agent après 被]
\begin{itemize}
\item \textbf{Exprimé} : si l'agent est une personne/entité précise et pertinente. \zh{被 + Agent + V...}
\item \textbf{Omis} : si l'agent est inconnu (\zh{小偷}), générique (\zh{大家}, \zh{人们}) ou sans importance. \zh{被} se colle alors au verbe : \zh{被 + V...}
\end{itemize}
\end{propriete}

\begin{exemplebox}[Agent omis (fréquent)]
\zh{钱包被偷了。} (Agent inconnu) \\
\emph{Qiánbāo bèi tōu le.}\\[4pt]
\zh{会议被取消了。} (Agent évident : les organisateurs) \\
\emph{Huìyì bèi qǔxiāo le.}\\[4pt]
\zh{中文在这里被广泛使用。} (Agent générique) \\
\emph{Zhōngwén zài zhèlǐ bèi guǎngfàn shǐyòng.}
\end{exemplebox}

\begin{popfigure}
\begin{tikzpicture}[node distance=0.7cm and 1.2cm, every node/.style={font=\small, align=center}]
\node[draw=popdark, fill=popcream, rounded corners, minimum width=5cm, minimum height=1cm, align=center] (q){\textbf{L'agent est-il connu,}\\\textbf{pertinent ou nécessaire ?}};
\node[draw=popblue, fill=popblueL, rounded corners, minimum width=4.5cm, minimum height=1.3cm, below left=of q, align=center] (oui){\textbf{OUI}\\ Patient + 被 + \textbf{Agent} + V\\ \zh{门被风吹开了。}};
\node[draw=poporange, fill=poporangeL, rounded corners, minimum width=4.5cm, minimum height=1.3cm, below right=of q, align=center] (non){\textbf{NON}\\ Patient + \textbf{被} + V (+ 了)\\ \zh{钱包被偷了。}};
\draw[-{Stealth}, thick, popdark] (q) -| node[pos=0.25, above left, font=\tiny]{OUI} (oui);
\draw[-{Stealth}, thick, popdark] (q) -| node[pos=0.25, above right, font=\tiny]{NON} (non);
\end{tikzpicture}
\legende{Choix de l'expression de l'agent.}
\end{popfigure}

\courssub{Le verbe : l'obligation du complément (Règle d'or)}

C'est l'erreur n°1 des francophones : traduire « a été volé » par *\zh{被偷} (verbe nu).

\begin{propriete}[Le verbe de la 被字句 ne reste jamais nu]
Le prédicat doit indiquer un \textbf{résultat}, un \textbf{accomplissement} ou une \textbf{portée}. On ajoute obligatoirement :
\begin{enumerate}
\item \textbf{Particule \zh{了}} (accomplissement) : \zh{信被寄了。} (\emph{Xìn bèi jì le.} — La lettre a été envoyée.)
\item \textbf{Complément de résultat} : \zh{杯子被打破了。} (\emph{Bēizi bèi dǎ pò le.} — Verre cassé \rightarrow résultat \zh{破}.)
\item \textbf{Complément de direction} : \zh{书被拿走了。} (\emph{Shū bèi ná zǒu le.} — Livre emporté \rightarrow direction \zh{走}.)
\item \textbf{Objet / Complément de portée} : \zh{他被选为班长。} (\emph{Tā bèi xuǎn wéi bānzhǎng.} — Il a été élu délégué.)
\item \textbf{Verbe composé / Verbe + 得/地 + compl.} : \zh{作业被做完了。} (\emph{Zuòyè bèi zuò wán le.})
\end{enumerate}
\end{propriete}

\begin{attention}[Verbe nu = Phrase fautive]
\begin{center}
\begin{tabularx}{\linewidth}{|X|X|X|}
\hline
\rowcolor{popblueL} \textbf{Phrase} & \textbf{Analyse} & \textbf{Correction} \\
\hline
*\zh{杯子被打了。} & \zh{打} seul = action non aboutie & \zh{杯子被打破了。} / \zh{杯子被打碎了。} \\
\hline
*\zh{信被寄。} & Verbe nu, pas de \zh{了} & \zh{信被寄了。} / \zh{信被寄出去了。} \\
\hline
*\zh{他被批评。} & Verbe nu (sauf style télégraphique) & \zh{他被批评了。} / \zh{他被狠狠批评了一顿。} \\
\hline
\end{tabularx}
\end{center}
\cle{Pourquoi ?} Le passif chinois (\zh{被}) décrit un \textbf{état résultant} ou un \textbf{événement clos}. Un verbe nu laisse l'action « en suspens ». Le français « a été cassé » inclut le résultat ; le chinois doit l'expliciter (\zh{破}, \zh{碎}, \zh{了}, \zh{完}...).
\end{attention}

\courssub{Négation et adverbes : TOUJOURS avant 被}

\begin{propriete}[Position de la négation et des adverbes]
\begin{center}
\textbf{(Négation / Adverbe) + 被 + Agent + Verbe + Complément}
\end{center}
\begin{itemize}
\item \textbf{Négation passée} : \zh{没(有)} + \zh{被}... (le \zh{了} final disparaît).
\item \textbf{Négation habituelle/future} : \zh{不} + \zh{被}... (rare).
\item \textbf{Adverbes} : \zh{已经}, \zh{刚}, \zh{又}, \zh{常常}, \zh{都}, \zh{一定}... tous se placent \textbf{avant} \zh{被}.
\end{itemize}
\end{propriete}

\begin{exemplebox}[Négation et adverbes]
\zh{门没有被风吹开。} (Négation passée : \zh{没} + \zh{被}, pas de \zh{了})\\
\emph{Mén méiyǒu bèi fēng chuī kāi.}\\
« La porte n'a pas été ouverte par le vent. »\\[4pt]
\zh{作业已经被做完了。}\\
\emph{Zuòyè yǐjīng bèi zuò wán le.}\\
« Les devoirs ont déjà été finis. »\\[4pt]
\zh{他又被老师批评了。}\\
\emph{Tā yòu bèi lǎoshī pīpíng le.}\\
« Il a encore été critiqué par le prof. »
\end{exemplebox}

\begin{attention}[Piège n°2 : Calque du français « n'a pas été »]
\begin{center}
\begin{tabularx}{\linewidth}{|X|X|}
\hline
\rowcolor{popblueL} \textbf{Erreur typique (Francophone)} & \textbf{Correct} \\
\hline
*\zh{他被没批评。} & \zh{他没被批评。} \\
*\zh{门被没有吹开。} & \zh{门没被吹开。} \\
*\zh{被已经做完了。} & \zh{已经被做完了。} \\
\hline
\end{tabularx}
\end{center}
\cle{Astuce mnémotechnique} : \zh{被} est un « aimant » qui attire l'agent à sa droite ; la négation et les adverbes, repoussés, restent à sa gauche.
\end{attention}

\courssub{Emplois, nuances et exceptions (Le « quand » utiliser 被)}

\begin{propriete}[Emploi principal : Passif adversatif / Mésaventure]
Contrairement au français où le passif est neutre (\emph{« La tour Eiffel a été construite par Eiffel »}), le \zh{被} chinois porte très souvent une \textbf{nuance négative} : dommage, perte, punition, mésaventure subie involontairement.
\end{propriete}

\begin{exemplebox}[Nuance de mésaventure (Typique 被)]
\zh{我的自行车被偷了。} (Perte) \\
\zh{他被开除了。} (Punition : Il a été viré.) \\
\zh{房子被烧毁了。} (Dégât : Maison brûlée.) \\
\zh{衣服让雨淋湿了。} (Mésaventure : Vêtements mouillés par la pluie.)
\end{exemplebox}

\begin{attention}[Passif « neutre » ou honorifique : 被 est souvent maladroit]
En chinois, pour exprimer un passif neutre (construction, découverte, écriture) ou honorifique, on préfère :
\begin{enumerate}
\item \textbf{Phrase active à sujet indéfini} : \zh{有人建造了这座桥。} (\emph{Yǒurén jiànzào le zhè zuò qiáo.} — Quelqu'un a construit ce pont / Ce pont a été construit.)
\item \textbf{Structure 把 (bǎ) inversée} (si focus sur l'objet) : \zh{这首歌是周杰伦写的。} (\emph{Zhè shǒu gē shì Zhōu Jiélún xiě de.} — Cette chanson, c'est Jay Chou qui l'a écrite.) $\leftarrow$ \textbf{Très fréquent} (Structure \zh{是...的}).
\item \textbf{Verbes psychologiques / statifs} : \zh{爱}, \zh{恨}, \zh{知道}, \zh{懂}, \zh{像}, \zh{有}... \textbf{ne s'emploient pas en 被}.
\end{enumerate}
\end{attention}

\begin{aretenir}[Verbes incompatibles avec 被 (Liste non exhaustive)]
\begin{center}
\begin{tabularx}{\linewidth}{|X|X|X|}
\hline
\rowcolor{popblueL} \textbf{Catégorie} & \textbf{Exemples} & \textbf{Alternative} \\
\hline
Verbes psychologiques & \zh{喜欢, 爱, 讨厌, 害怕, 后悔} & Actif : \zh{大家都喜欢他。} \\
Verbes cognitifs & \zh{知道, 认识, 明白, 懂, 了解} & Actif : \zh{大家都知道这件事。} \\
Verbes d'état/relation & \zh{有, 是, 像, 属于, 等于} & Actif / \zh{是...的} \\
Verbes « subir » intrinsèques & \zh{受伤, 生病, 去世, 睡着} & Déjà passifs par nature : \zh{他受伤了。} (pas *\zh{被受伤}) \\
\hline
\end{tabularx}
\end{center}
\end{aretenir}

\courssub{Comparaison Français / Chinois : Tableau des pièges}

\begin{center}
\begin{tabularx}{\linewidth}{|X|X|X|}
\hline
\rowcolor{popblueL} \textbf{Point de grammaire} & \textbf{Français} & \textbf{Chinois (被字句)} \\
\hline
\textbf{Sujet passif} & Peut être indéfini (\emph{Un livre a été lu}) & \textbf{Doit être défini/connu} (\zh{那本书被读了}) \\
\hline
\textbf{Verbe} & Participe passé seul (\emph{a été lu}) & \textbf{Verbe + Complément obligatoire} (\zh{读完了}, \zh{读懂了}) \\
\hline
\textbf{Négation} & Autour de l'auxiliaire (\emph{n'a pas été lu}) & \textbf{Devant 被} (\zh{没被读完}) \\
\hline
\textbf{Adverbes} & Autour du verbe (\emph{a déjà été lu}) & \textbf{Devant 被} (\zh{已经被读完了}) \\
\hline
\textbf{Agent} & \emph{Par} + agent (souvent omis) & \textbf{Après 被} (omis si inconnu) ; \textbf{obligatoire avec 叫/让} \\
\hline
\textbf{Registre} & Neutre, administratif, scientifique & \textbf{Marqué : mésaventure, dommage} ; neutre $\rightarrow$ Actif / \zh{是...的} \\
\hline
\end{tabularx}
\end{center}

\courssub{Variantes orales : 叫, 让, 给}

À l'oral (Nord de la Chine surtout), \zh{被} est remplacé par \zh{叫} (\emph{jiào}), \zh{让} (\emph{ràng}) ou \zh{给} (\emph{gěi}).

\begin{propriete}[Contraintes des variantes orales]
\begin{itemize}
\item \textbf{Registre} : Familier / Oral.
\item \textbf{Agent OBLIGATOIRE} : On ne dit pas *\zh{钱包叫偷了}. Il faut \zh{钱包叫小偷偷了}.
\item \textbf{给} : Souvent suivi de l'agent, structure \zh{Patient + 给 + Agent + V...}.
\end{itemize}
\end{propriete}

\begin{center}
\begin{tabularx}{\linewidth}{|X|X|X|X|}
\hline
\rowcolor{popblueL} \textbf{Particule} & \textbf{Registre} & \textbf{Agent omis ?} & \textbf{Exemple} \\
\hline
\zh{被} & Neutre, Écrit/Oral & \textbf{Oui} & \zh{手机被偷了。} \\
\hline
\zh{叫} & Oral, Familier & \textbf{Non} & \zh{手机叫小偷偷了。} \\
\hline
\zh{让} & Oral, Familier & \textbf{Non} & \zh{手机让小偷偷了。} \\
\hline
\zh{给} & Oral, Très familier & \textbf{Non} & \zh{手机给小偷偷了。} \\
\hline
\end{tabularx}
\end{center}

\begin{exemplebox}[Variantes à l'oral]
\zh{我的自行车让人骑走了。}\\
\emph{Wǒ de zìxíngchē ràng rén qí zǒu le.}\\
« Mon vélo a été emmené par quelqu'un (on me l'a pris). »\\[4pt]
\zh{饭给狗吃了。}\\
\emph{Fàn gěi gǒu chī le.}\\
« Le riz a été mangé par le chien. »
\end{exemplebox}

\courssub{Vocabulaire clé de la leçon}

\begin{center}
\begin{tabularx}{\linewidth}{|X|X|X|X|}
\hline
\rowcolor{popblueL} \textbf{Caractères} & \textbf{Pinyin} & \textbf{Nature} & \textbf{Traduction} \\
\hline
被 & bèi & Part. (Passif) & (marque le passif) \\
叫 & jiào & Part. (Passif oral) & (marque le passif oral) \\
让 & ràng & Part. (Passif oral) / V & (marque le passif oral) / laisser, faire \\
给 & gěi & Part. (Passif oral) / V / Prép. & (marque le passif oral) / donner / pour \\
偷 & tōu & V & voler \\
打破 & dǎ pò & V (V+R) & casser (en frappant) \\
借走 & jiè zǒu & V (V+D) & emporter en empruntant \\
找回 & zhǎo huí & V (V+R) & retrouver, récupérer \\
批评 & pīpíng & V / N & critiquer / critique \\
开除 & kāichú & V & renvoyer, virer \\
取消 & qǔxiāo & V & annuler \\
广泛 & guǎngfàn & Adv / Adj & largement, étendu \\
已经 & yǐjīng & Adv & déjà \\
又 & yòu & Adv & encore, de nouveau \\
刚 & gāng & Adv & à l'instant, venir de \\
\hline
\end{tabularx}
\end{center}

\courssub{Écriture des caractères : 被, 叫, 让}

\begin{definition}[Analyse graphique et ordre des traits]
\begin{itemize}
\item \zh{被} (\emph{bèi}) : 10 traits. Clé : \zh{衣} (vêtement, radical 145). Haut : \zh{皮} (peau) ; Bas : \zh{衣}. Ordre : haut \rightarrow bas (写 \zh{皮} 然后 \zh{衣}). Sens originel : couverture de paille $\rightarrow$ subir, couvrir.
\item \zh{叫} (\emph{jiào}) : 10 traits. Clé : \zh{口} (bouche, radical 30). Phonétique : \zh{丩} (jiū). Ordre : \zh{口} (verticale, horizontales, fermeture) + \zh{丩}. Sens : crier, appeler $\rightarrow$ passif oral.
\item \zh{让} (\emph{ràng}) : 17 traits. Clé : \zh{讠} (parole, radical 149). Phonétique : \zh{上} + \zh{一} + \zh{八} (simplifié de \zh{譲}). Ordre : \zh{讠} (point, trait horizontal courbe) + partie droite (haut \rightarrow bas). Sens : céder, laisser $\rightarrow$ passif oral.
\end{itemize}
\end{definition}

\courssub{Méthode : Construire et transformer une phrase en 被}

\begin{methode}[De l'actif au passif (被) en 5 étapes]
\begin{enumerate}
\item \textbf{Identifier} : Sujet (Agent) / Verbe + Complément / Objet (Patient) de la phrase active. \textit{Ex: 弟弟 打破了 杯子。}
\item \textbf{Extraire le Patient} : Le placer en tête. \textit{Ex: 杯子...}
\item \textbf{Ajouter 被} : \textit{Ex: 杯子被...}
\item \textbf{Gérer l'Agent} : Le placer après \zh{被} (si connu/pertinent) ou omettre \zh{被} seul. \textit{Ex: 杯子被弟弟...}
\item \textbf{Vérifier le Verbe} : \textbf{JAMAIS NU}. Ajouter \zh{了} / Résultat / Direction / Objet. \textit{Ex: 杯子被弟弟打破了。}
\item \textbf{Négation ?} Placer \zh{没(有)} / \zh{不} \textbf{DEVANT 被}. \textit{Ex: 杯子没被弟弟打破。}
\end{enumerate}
\end{methode}

\begin{exoresolu}[Transformation complète : Actif $\rightarrow$ Passif $\rightarrow$ Négatif]
Énoncé : Transformer \zh{风把门吹开了。} (\emph{Fēng bǎ mén chuī kāi le.} — Le vent a ouvert la porte) en passif \zh{被}, puis mettre à la forme négative passée.
\tcblower
\textbf{Solution détaillée :}
\begin{enumerate}
\item \textbf{Analyse active} : Agent = \zh{风} (vent) ; Structure \zh{把} : \zh{风把门吹开} ; Complément de direction \zh{开} + \zh{了}.
\item \textbf{Passif} : Patient \zh{门} + \zh{被} + Agent \zh{风} + Verbe \zh{吹开} + \zh{了} $\rightarrow$ \zh{门被风吹开了。} (\emph{Mén bèi fēng chuī kāi le.})
\item \textbf{Négatif passé} : \zh{没(有)} devant \zh{被} ; \zh{了} final \textbf{disparaît} ; \zh{开} reste (complément de direction). $\rightarrow$ \zh{门没有被风吹开。} (\emph{Mén méiyǒu bèi fēng chuī kāi.})
\end{enumerate}
\textbf{Vérification} : Verbe non nu (\zh{吹开}), Négation avant \zh{被}, \zh{了} absent au négatif. \checkmark
\end{exoresolu}

\courssub{Exercices d'application et de transformation}

\begin{exercice}[Exercice 1 : Choisir la bonne structure]
Complétez par \zh{被}, \zh{叫}, \zh{让} ou \zh{给} (parfois plusieurs réponses possibles, précisez le registre).
\begin{enumerate}
\item 我的钱包 \_\_\_\_\_\_ 小偷偷了。 (Écrit / Oral)
\item 这件事 \_\_\_\_\_\_ 大家都知道了。 (Passif possible ?)
\item 作业 \_\_\_\_\_\_ 老师批改完了。 (Registre standard)
\item 门 \_\_\_\_\_\_ 风吹开了。 (Agent = force naturelle)
\item 手机 \_\_\_\_\_\_ 偷了。 (Agent inconnu)
\end{enumerate}
\end{exercice}

\begin{exercice}[Exercice 2 : Corriger les erreurs]
Chaque phrase contient une erreur (ordre, mot-outil, complément manquant, négation). Corrigez et expliquez.
\begin{enumerate}
\item *\zh{杯子被打了。}
\item *\zh{他被没批评。}
\item *\zh{一本书被借走了。}
\item *\zh{钱包叫偷了。}
\item *\zh{作业被已经做完了。}
\item *\zh{他被喜欢大家。}
\end{enumerate}
\end{exercice}

\begin{exercice}[Exercice 3 : Transformation Actif $\rightarrow$ Passif (被)]
Transformez les phrases actives suivantes en phrases passives en \zh{被}. Attention au complément du verbe !
\begin{enumerate}
\item \zh{小明吃完了那个苹果。} (\emph{Xiǎomíng chī wán le nà gè píngguǒ.})
\item \zh{老师批改了我的作业。}
\item \zh{大雨淋湿了我的衣服。} (\emph{Dàyǔ lín shī le wǒ de yīfu.})
\item \zh{有人打扫了教室。} (\emph{Yǒurén dǎsǎo le jiàoshì.})
\item \zh{警察抓住了小偷。} (\emph{Jǐngchá zhuā zhù le xiǎotōu.})
\end{enumerate}
\end{exercice}

\begin{exercice}[Exercice 4 : Transformation Négative]
Mettez les phrases passives suivantes à la forme négative (passé accompli).
\begin{enumerate}
\item \zh{信被寄出去了。}
\item \zh{窗户被打破了。}
\item \zh{他被选为班长了。}
\end{enumerate}
\end{exercice}

\begin{exercice}[Exercice 5 : Traduction thème (Français $\rightarrow$ Chinois)]
Traduisez en chinois (caractères + pinyin) en utilisant la structure \zh{被}.
\begin{enumerate}
\item Mon vélo a été volé (par un voleur).
\item La porte n'a pas été fermée par le vent.
\item Les devoirs ont déjà été finis par moi.
\item Ce plat a été mangé par le chien. (Oral : \zh{叫/让/给})
\item Son portable a encore été cassé.
\end{enumerate}
\end{exercice}

\courssub{Corrigés détaillés}

\begin{corrige}[Exercice 1]
\begin{enumerate}
\item \zh{被} (écrit/neutre), \zh{叫/让/给} (oral, agent \zh{小偷} exprimé).
\item \textbf{Piège} : \zh{知道} (verbe cognitif) $\nexists$ passif \zh{被}. Phrase active : \zh{大家都知道这件事。}
\item \zh{被} (standard). \zh{叫/让/给} possibles à l'oral.
\item \zh{被} (standard, force naturelle = agent). \zh{让} possible à l'oral (\zh{门让风吹开了}).
\item \zh{被} \textbf{seul} (agent omis = inconnu). \zh{叫/让/给} \textbf{impossibles} (agent obligatoire avec eux).
\end{enumerate}
\end{corrige}

\begin{corrige}[Exercice 2]
\begin{enumerate}
\item *\zh{杯子被打了。} $\rightarrow$ \zh{杯子被打破了。} / \zh{杯子被打碎了。} (Verbe nu $\rightarrow$ ajout complément de résultat \zh{破}/\zh{碎}).
\item *\zh{他被没批评。} $\rightarrow$ \zh{他没被批评。} (Négation \zh{没} devant \zh{被}).
\item *\zh{一本书被借走了。} $\rightarrow$ \zh{那本书被借走了。} (Patient défini) OU \zh{有人借走了一本书。} (Actif).
\item *\zh{钱包叫偷了。} $\rightarrow$ \zh{钱包叫小偷偷了。} (Agent obligatoire avec \zh{叫}) OU \zh{钱包被偷了。} (Agent omis avec \zh{被}).
\item *\zh{作业被已经做完了。} $\rightarrow$ \zh{作业已经被做完了。} (Adverbe \zh{已经} devant \zh{被}).
\item *\zh{他被喜欢大家。} $\rightarrow$ \zh{大家都喜欢他。} (\zh{喜欢} verbe psychologique, pas de passif \zh{被}).
\end{enumerate}
\end{corrige}

\begin{corrige}[Exercice 3]
\begin{enumerate}
\item \zh{那个苹果被小明吃完了。} (\emph{Nà gè píngguǒ bèi Xiǎomíng chī wán le.}) — Complément de résultat \zh{完}.
\item \zh{我的作业被老师批改了。} (\emph{Wǒ de zuòyè bèi lǎoshī pīgǎi le.}) — \zh{了} accomplissement.
\item \zh{我的衣服被大雨淋湿了。} (\emph{Wǒ de yīfu bèi dàyǔ lín shī le.}) — Complément de résultat \zh{湿}.
\item \zh{教室被打扫了。} (\emph{Jiàoshì bèi dǎsǎo le.}) — Agent \zh{有人} omis (générique) ; \zh{了} obligatoire.
\item \zh{小偷被警察抓住了。} (\emph{Xiǎotōu bèi jǐngchá zhuā zhù le.}) — Complément de résultat \zh{住}.
\end{enumerate}
\end{corrige}

\begin{corrige}[Exercice 4]
\begin{enumerate}
\item \zh{信没被寄出去。} (\emph{Xìn méi bèi jì chū qù.}) — \zh{没} + \zh{被}, pas de \zh{了}.
\item \zh{窗户没被打破。} (\emph{Chuānghu méi bèi dǎ pò.})
\item \zh{他没被选为班长。} (\emph{Tā méi bèi xuǎn wéi bānzhǎng.})
\end{enumerate}
\end{corrige}

\begin{corrige}[Exercice 5]
\begin{enumerate}
\item \zh{我的自行车被小偷偷了。} (\emph{Wǒ de zìxíngchē bèi xiǎotōu tōu le.}) / Oral : \zh{我的自行车叫小偷偷了。}
\item \zh{门没被风关上。} (\emph{Mén méi bèi fēng guān shàng.}) — Verbe \zh{关上} (fermer + direction haut/résultat).
\item \zh{作业已经被我做完了。} (\emph{Zuòyè yǐjīng bèi wǒ zuò wán le.}) — Agent \zh{我} (1ère pers.) accepté en chinois.
\item \zh{这道菜叫狗吃了。} (\emph{Zhè dào cài jiào gǒu chī le.}) — Oral, agent \zh{狗} obligatoire.
\item \zh{他的手机又被摔坏了。} (\emph{Tā de shǒujī yòu bèi shuāi huài le.}) — \zh{摔坏} (casser en faisant tomber), \zh{又} devant \zh{被}.
\end{enumerate}
\end{corrige}

\courssub{Élément socioculturel : Responsabilité et formulation du passif}

\begin{savaistu}[Pourquoi le passif « 被 » sonne-t-il souvent « négatif » en chinois ?]
En chinois, la responsabilité individuelle est culturellement mise en avant. La structure active \zh{谁做的？} (Qui a fait ça ?) est la norme. Le passif \zh{被} déplace le focus sur la \textbf{victime} et le \textbf{résultat subi}, ce qui implique souvent : « Je n'ai pas pu l'empêcher », « C'est arrivé \emph{à moi} ». C'est pourquoi on l'utilise pour les vols, accidents, punitions, catastrophes.
\par\medskip
\textbf{Au Cameroun :} En français camerounais, on utilise aussi le passif pour atténuer la responsabilité (« \emph{La porte a été fermée} » au lieu de « \emph{Quelqu'un a fermé la porte} »). En chinois, si vous dites \zh{门被关了}, cela sous-entend « La porte s'est retrouvée fermée (et ce n'est pas normal / c'est embêtant) ». Pour un fait neutre (ex: « Le pont a été construit en 2020 »), préférez : \zh{这座桥是2020年建的。} (\emph{Zhè zuò qiáo shì 2020 nián jiàn de.}) — Structure \zh{是...的} (focus sur le temps/agent/moyen), neutre et naturelle.
\end{savaistu}

\begin{experience}[Activité orale : « Enquête au commissariat »]
\textbf{Objectif :} Produire des phrases en \zh{被} (plainte, déclaration) et \zh{叫/让} (récit oral entre amis).
\textbf{Consigne :} Par binômes. A = Victime (dépose plainte, registre formel \zh{被}). B = Policier (note, pose questions \zh{怎么被...的？}). Puis inversion : A raconte l'histoire à un ami (registre oral \zh{叫/让/给}).
\textbf{Scénarios :} 1. Portable volé dans un taxi à Douala. 2. Devoirs mangés par le chien. 3. Vêtements mouillés par la pluie soudaine. 4. Vélo emprunté sans permission par le petit frère.
\textbf{Contraintes :} Minimum 3 phrases \zh{被} par scénario. Vérifier : Patient défini, Verbe accompagné, Négation devant \zh{被} si besoin.
\end{experience}

\coursec{Leçon 18 — Grammaire : 被字句}

\courssub{Observation d'exemples : le passif en situation}

Avant d'analyser la structure, observons comment le chinois exprime le fait de \emph{subir} une action dans des situations authentiques. Le dialogue suivant se déroule entre deux camarades de classe à Yaoundé après un orage violent.

\begin{textebox}[Dialogue : Après l'orage]
\zh{— 嘿，你看那棵大树！} \emph{Hēi, nǐ kàn nà kē dà shù!} — « Hé, regarde cet arbre ! »\\
\zh{— 唉，被风吹倒了。} \emph{Ài, bèi fēng chuīdǎo le.} — « Ah, il a été renversé par le vent. »\\
\zh{— 还有那辆自行车，也是新买的。} \emph{Hái yǒu nà liàng zìxíngchē, yě shì xīn mǎi de.} — « Et ce vélo-là, il est neuf en plus. »\\
\zh{— 被压坏了，真倒霉。} \emph{Bèi yā huài le, zhēn dǎoméi.} — « Il a été écrasé (et abîmé), quelle malchance. »\\
\zh{— 听说小林被选为班长了？} \emph{Tīngshuō Xiǎo Lín bèi xuǎn wéi bānzhǎng le?} — « J'ai entendu dire que Xiao Lin a été élu délégué ? »\\
\zh{— 对，他受到了老师的表扬。} \emph{Duì, tā shòudào le lǎoshī de biǎoyáng.} — « Oui, il a reçu les éloges du prof. »
\end{textebox}

\begin{aretenir}[Premiers constats]
\begin{itemize}
\item La marque \zh{被} \emph{bèi} place le \textbf{patient} (l'arbre, le vélo, Xiao Lin) en position de sujet.
\item L'agent (le vent, les électeurs, le prof) peut être exprimé après \zh{被} ou omis.
\item Le verbe est \textbf{toujours suivi d'un complément} (résultat, direction, aspect, \zh{了}…), jamais « nu ».
\item On rencontre aussi \zh{受到} \emph{shòudào} pour un événement favorable formel (éloges).
\end{itemize}
\end{aretenir}

\courssub{Schéma de la structure 被字句}

\begin{propriete}[Structure canonique]
La phrase en \zh{被} (passif marqué) suit le schéma immuable :
\[
\textbf{Patient + 被 + (Agent) + Verbe + Complément obligatoire}
\]
\begin{itemize}
\item \textbf{Patient} : l'entité qui subit l'action (sujet grammatical). Il est \textbf{défini} (démonstratif, possessif, contexte).
\item \textbf{被} \emph{bèi} : marqueur du passif (préposition d'agent).
\item \textbf{Agent} : l'auteur de l'action. \textbf{Peut être omis} s'il est inconnu, évident ou générique (\zh{人} \emph{rén} « on, les gens »).
\item \textbf{Verbe} : \textbf{verbe transitif} à sens dynamique (action concrète).
\item \textbf{Complément} : \textbf{Obligatoire}. Résultat (\zh{坏} \emph{huài}, \zh{好} \emph{hǎo}, \zh{完} \emph{wán}), direction (\zh{走} \emph{zǒu}, \zh{来} \emph{lái}), \zh{了} \emph{le} (aspect accompli), \zh{成} \emph{chéng} + nom, etc. \cle{Un verbe nu après 被 est incorrect.}
\end{itemize}
\end{propriete}

\begin{exemplebox}[Analyse de la structure]
\zh{我的自行车} \emph{Wǒ de zìxíngchē} \textbf{(Patient)} + \zh{被} \emph{bèi} + \zh{小偷} \emph{xiǎotōu} \textbf{(Agent)} + \zh{偷} \emph{tōu} \textbf{(Verbe)} + \zh{走了} \emph{zǒu le} \textbf{(Complément directionnel + aspect)}.\\
\zh{Mon vélo a été volé par un voleur.}
\end{exemplebox}

\begin{attention}[Erreur n°1 : verbe « nu » après 被]
\begin{itemize}
\item \textbf{Faux :} \zh{* 信被写。} \emph{* Xìn bèi xiě.} (Le verbe \zh{写} est nu.)
\item \textbf{Vrai :} \zh{信被写好了。} \emph{Xìn bèi xiě hǎo le.} (Complément de résultat \zh{好} + \zh{le}).
\item \textbf{Vrai :} \zh{信被写完了。} \emph{Xìn bèi xiě wán le.} (Complément \zh{完} + \zh{le}).
\end{itemize}
\emph{Règle d'or :} Après le verbe, il faut \textbf{toujours} un élément qui « referme » l'action (résultat, achèvement, déplacement, \zh{了}).
\end{attention}

\begin{exemplebox}[Agent omis : passif « agentless »]
\zh{作业被改错了。} \emph{Zuòyè bèi gǎi cuò le.} — Le devoir a été corrigé à tort (par quelqu'un).\\
\zh{门被关上了。} \emph{Mén bèi guān shàng le.} — La porte a été fermée (par quelqu'un / par le vent).\\
\zh{手机被偷了。} \emph{Shǒujī bèi tōu le.} — Le téléphone a été volé (agent inconnu).
\end{exemplebox}

\courssub{Emplois, nuances et exceptions}

\courssub{Emploi principal : l'action subie, souvent défavorable}

\begin{propriete}[被 et la « couleur » défavorable]
Traditionnellement, la phrase en \zh{被} exprime une action \textbf{subie} par le patient, le plus souvent \textbf{défavorable, inattendue ou regrettable} : vol, perte, casse, critique, punition, tromperie, maladie, catastrophe naturelle. Le locuteur signale ainsi que le patient n'a rien demandé et subit l'événement. En chinois moderne, sous l'influence des langues occidentales, l'emploi de \zh{被} s'est \textbf{élargi aux faits neutres ou favorables} : cet usage est aujourd'hui courant et accepté, surtout à l'écrit et dans les médias.
\end{propriete}

\begin{exemplebox}[Actions défavorables typiques (contexte camerounais)]
\zh{他的钱包在杜阿拉的市场被偷了。} \emph{Tā de qiánbāo zài Dù'ālā de shìchǎng bèi tōu le.} — Son portefeuille a été volé au marché de Douala.\\
\zh{小明被爸爸罚站了一个小时。} \emph{Xiǎo Míng bèi bàba fá zhàn le yí ge xiǎoshí.} — Xiao Ming a été puni par son père : une heure debout.\\
\zh{我们的球队被打败了。} \emph{Wǒmen de qiúduì bèi dǎbài le.} — Notre équipe a été battue.\\
\zh{那棵芒果树被风吹倒了。} \emph{Nà kē mángguǒ shù bèi fēng chuīdǎo le.} — Ce manguier a été renversé par le vent.\\
\zh{她被雨淋湿了。} \emph{Tā bèi yǔ lín shī le.} — Elle a été trempée par la pluie.
\end{exemplebox}

\begin{exemplebox}[Faits neutres ou favorables (usage moderne standard)]
\zh{他被选为班长。} \emph{Tā bèi xuǎn wéi bānzhǎng.} — Il a été élu délégué de classe.\\
\zh{这本书被翻译成好几种语言。} \emph{Zhè běn shū bèi fānyì chéng hǎo jǐ zhǒng yǔyán.} — Ce livre a été traduit en plusieurs langues.\\
\zh{她被雅温得的一所大学录取了。} \emph{Tā bèi Yǎwēndé de yì suǒ dàxué lùqǔ le.} — Elle a été admise dans une université de Yaoundé.\\
\zh{这个传统被保留下来了。} \emph{Zhè ge chuántǒng bèi bǎoliú xiàlai le.} — Cette tradition a été préservée.
\end{exemplebox}

\begin{attention}[Ne pas « passer tout au passif » comme en français]
Le français passe au passif presque toutes les phrases (« ce plat se mange froid », « la porte est ouverte à 8 h », « le pont a été construit en 2020 »). Le chinois, lui, préfère souvent la \textbf{phrase active} (sujet agent) ou le \textbf{passif sans marque} (voir §4). Traduire mécaniquement chaque passif français par \zh{被} est l'erreur n° 1 des francophones.\\
« Ce plat se mange froid. » → \zh{这个菜要冷着吃。} \emph{Zhè ge cài yào lěngzhe chī.} (Active : « On mange ce plat froid »)\\
« Les billets sont vendus ici. » → \zh{这儿卖票。} \emph{Zhèr mài piào.} (Active sans sujet : « Ici on vend des billets »)\\
« Le pont a été construit en 2020. » → \zh{这座桥是2020年建的。} \emph{Zhè zuò qiáo shì 2020 nián jiàn de.} (Construction \zh{是…的} pour insister sur le temps/agent).
\end{attention}

\courssub{Le patient doit être défini ou connu}

\begin{propriete}[Définition du patient]
En tête d'une phrase en \zh{被}, le \textbf{patient doit être défini ou connu} des interlocuteurs : il est précédé d'un démonstratif (\zh{这} \emph{zhè}, \zh{那} \emph{nà}), d'un possessif (\zh{我的} \emph{wǒ de}), ou désigne une chose unique dans le contexte (\zh{教室的窗户} \emph{jiàoshì de chuānghu}). On dit \zh{那本书被借走了。} « Ce livre-là a été emprunté », mais on évite de commencer par un patient indéfini du type « un livre » (\zh{* 一本书被借走了} est maladroit). C'est logique : la phrase en \zh{被} commente le sort d'une chose \emph{identifiée}.
\end{propriete}

\begin{exemplebox}[Patients définis]
\zh{我的手机被摔坏了。} \emph{Wǒ de shǒujī bèi shuāi huài le.} — Mon téléphone a été cassé (en tombant).\\
\zh{教室的窗户被打开了。} \emph{Jiàoshì de chuānghu bèi dǎkāi le.} — La fenêtre de la salle a été ouverte.\\
\zh{饭被妹妹做好了。} \emph{Fàn bèi mèimei zuò hǎo le.} — Le repas a été préparé par ma petite sœur.
\end{exemplebox}

\courssub{Le passif sans marque : quand le chinois n'utilise PAS 被}

\begin{propriete}[Phrases à sens passif sans 被 (不被字句)]
Quand le patient est en tête de phrase, que l'agent est absent ou évident, et que le verbe porte un complément de résultat/d'achèvement, le chinois emploie très souvent une \textbf{phrase à thème patient sans aucune marque passive}. C'est le passif le plus fréquent de la langue parlée courante. Structure : \textbf{Patient + Verbe + Complément (+ 了)}.
\end{propriete}

\begin{exemplebox}[Passif sans marque (constat de résultat)]
\zh{信写好了。} \emph{Xìn xiě hǎo le.} — La lettre est (a été) écrite.\\
\zh{作业做完了。} \emph{Zuòyè zuò wán le.} — Les devoirs sont finis.\\
\zh{门关上了。} \emph{Mén guān shàng le.} — La porte a été fermée.\\
\zh{票卖完了。} \emph{Piào mài wán le.} — Les billets sont épuisés (tous vendus).\\
\zh{菜切好了。} \emph{Cài qiē hǎo le.} — Les légumes sont coupés.
\end{exemplebox}

\begin{propriete}[Choix : 被 vs Passif sans marque]
\begin{itemize}
\item \zh{被} : insiste sur le fait que l'événement est \emph{subi} (souvent mal vécu, involontaire) et \textbf{permet de nommer l'agent}.
\item Passif sans marque : constate simplement un \emph{résultat} ou un \emph{état achevé}, agent irrecevable ou inutile.
\end{itemize}
\zh{杯子打破了。} \emph{Bēizi dǎpò le.} « Le verre est cassé » (constat neutre, résultat).\\
\zh{杯子被弟弟打破了。} \emph{Bēizi bèi dìdi dǎpò le.} « Le verre a été cassé par mon petit frère » (on souligne l'accident et le coupable).
\end{propriete}

\courssub{受到 / 得到 + nom verbal : le passif des actions favorables formelles}

\begin{definition}[受到 et 得到]
Pour les actions \textbf{favorables et formelles} (éloges, aide, accueil, soutien, attention, éducation, influence, encouragement), le chinois emploie plutôt \zh{受到} \emph{shòudào} « recevoir, être l'objet de » ou \zh{得到} \emph{dédào} « obtenir » suivis d'un \textbf{nom verbal} (souvent verbe + \zh{的} ou nom d'action), au lieu de \zh{被}. Structure : \textbf{Sujet + 受到/得到 + (agent 的) + Nom verbal}.
\end{definition}

\begin{exemplebox}[受到 / 得到 — registre formel / écrit]
\zh{他受到了老师的表扬。} \emph{Tā shòudào le lǎoshī de biǎoyáng.} — Il a reçu les éloges du professeur.\\
\zh{这位医生得到了病人的信任。} \emph{Zhè wèi yīshēng dédào le bìngrén de xìnrèn.} — Ce médecin a gagné la confiance de ses patients.\\
\zh{代表团受到了热烈的欢迎。} \emph{Dàibiǎotuán shòudào le rèliè de huānyíng.} — La délégation a reçu un accueil chaleureux.\\
\zh{这个提议得到了大家的支持。} \emph{Zhè ge tíyì dédào le dàjiā de zhīchí.} — Cette proposition a obtenu le soutien de tous.
\end{exemplebox}

\courssub{Verbes incompatibles avec 被}

\begin{attention}[Tous les verbes ne passent pas au passif en 被]
Les verbes qui n'ont \textbf{aucun effet concret} sur le patient — verbes d'état, de sentiment, de connaissance, d'existence — ne peuvent pas construire une phrase en \zh{被} :
\begin{itemize}
\item \zh{爱} \emph{ài} (aimer), \zh{喜欢} \emph{xǐhuan} (aimer), \zh{讨厌} \emph{tǎoyàn} (détester)
\item \zh{知道} \emph{zhīdào} (savoir), \zh{认识} \emph{rènshi} (connaître), \zh{了解} \emph{liǎojiě} (comprendre)
\item \zh{有} \emph{yǒu} (avoir), \zh{是} \emph{shì} (être), \zh{属于} \emph{shǔyú} (appartenir à)
\end{itemize}
On ne peut pas dire « il est aimé par tous » avec \zh{被}. Le chinois utilisera la \textbf{phrase active} : \zh{大家都喜欢他。} \emph{Dàjiā dōu xǐhuan tā.} « Tout le monde l'aime. » De même, les verbes intransitifs (\zh{来} \emph{lái}, \zh{去} \emph{qù}, \zh{睡觉} \emph{shuìjiào}, \zh{死} \emph{sǐ}) sont exclus, puisqu'ils n'ont pas de patient.
\end{attention}

\begin{center}
\begin{tabularx}{\linewidth}{|X|X|X|}
\hline
\rowcolor{popblueL} Situation & Exemple (caractères + pinyin) & Traduction \\
\hline
Action défavorable (agent exprimé) & \zh{我的自行车被偷了。} Wǒ de zìxíngchē bèi tōu le. & Mon vélo a été volé. \\
\hline
Fait favorable (usage moderne) & \zh{他被选为班长。} Tā bèi xuǎn wéi bānzhǎng. & Il a été élu délégué de classe. \\
\hline
Passif sans marque (constat) & \zh{信写好了。} Xìn xiě hǎo le. & La lettre est écrite. \\
\hline
Action favorable formelle & \zh{他受到表扬。} Tā shòudào biǎoyáng. & Il reçoit des éloges. \\
\hline
Verbe d'état/sentiment : pas de 被 & \zh{大家都喜欢他。} Dàjiā dōu xǐhuan tā. & Il est aimé de tous (Tout le monde l'aime). \\
\hline
Français passif $\rightarrow$ Actif chinois & \zh{这儿卖票。} Zhèr mài piào. & Les billets sont vendus ici (On vend des billets ici). \\
\hline
\end{tabularx}
\end{center}

\begin{popfigure}
\begin{tikzpicture}[mindmap, grow cyclic, every node/.style={concept, align=center, text=white},
root concept/.append style={concept color=popblue, font=\Large\bfseries},
level 1/.append style={level distance=3.8cm, sibling angle=120},
level 2/.append style={level distance=2.8cm, sibling angle=50, font=\small}]
\node[root concept, align=center]{被字句\\ bèi zì jù}
  child[concept color=poporange] { node[align=center]{Action\\ défavorable}
    child { node[align=center]{我的自行车\\ 被偷了。} }
    child { node {被风吹倒} }
  }
  child[concept color=popgreen] { node[align=center]{Agent omis\\ (被 + V + Compl.)}
    child { node[align=center]{杯子被\\ 打破了。} }
    child { node[align=center]{作业被\\ 改错了。} }
  }
  child[concept color=poppurple] { node[align=center]{Fait neutre /\\ favorable (moderne)}
    child { node[align=center]{他被选为\\ 班长。} }
    child { node[align=center]{书被翻译成\\ 外语。} }
  };
\end{tikzpicture}
\legende{Carte mentale : les trois grands emplois de la phrase en 被.}
\end{popfigure}

\courssub{Comparaison avec le français et pièges (Synthèse)}

\begin{methode}[Traduire un passif français vers le chinois : arbre de décision]
\begin{enumerate}
\item \textbf{Le verbe français exprime-t-il un sentiment/état/connaissance ?} (\emph{aimer, savoir, avoir, être}) $\rightarrow$ \textbf{Active chinoise} : Sujet (agent) + Verbe + Objet (patient). Ex: \emph{Il est connu de tous} $\rightarrow$ \zh{大家都认识他。}
\item \textbf{L'agent est-il important / l'événement est-il subi (souvent malheur) ?} $\rightarrow$ \textbf{被字句} : Patient + 被 + Agent + V + Compl. Ex: \emph{Le vélo a été volé par un voleur} $\rightarrow$ \zh{自行车被小偷偷走了。}
\item \textbf{Est-ce une action favorable formelle ?} (\emph{éloges, accueil, confiance, soutien}) $\rightarrow$ \textbf{受到 / 得到 + Nom verbal}. Ex: \emph{Il a été accueilli chaleureusement} $\rightarrow$ \zh{他受到了热烈欢迎。}
\item \textbf{Simple constat de résultat, pas d'agent, pas de malheur ?} $\rightarrow$ \textbf{Passif sans marque} : Patient + V + Complément (+ 了). Ex: \emph{Les devoirs sont faits} $\rightarrow$ \zh{作业做完了。}
\item \textbf{Passif français impersonnel / générique ?} (\emph{On vend des billets, Cela se fait ainsi}) $\rightarrow$ \textbf{Active chinoise sans sujet} : \zh{这儿卖票。} / \zh{这儿这么做。}
\end{enumerate}
\end{methode}

\begin{savaistu}[被, une structure bien vivante : l'ironie du net]
Dans la presse et sur les réseaux sociaux chinois, \zh{被} sert à créer des \textbf{néologismes ironiques} qui détournent la structure pour dénoncer une contrainte administrative : \zh{被就业} \emph{bèi jiùyè} « être déclaré employé » (quand une université gonfle ses statistiques d'emploi des diplômés), \zh{被代表} \emph{bèi dàibiǎo} « être représenté sans l'avoir demandé », \zh{被自愿} \emph{bèi zìyuàn} « être volontaire de force ». Ces jeux de langage, très populaires, montrent la vitalité de la structure : \zh{被} garde sa nuance fondamentale — \emph{subir quelque chose qu'on n'a pas choisi} — même quand on l'utilise pour plaisanter ou critiquer.
\end{savaistu}

\begin{aretenir}[L'essentiel des emplois — Fiche de révision]
\begin{itemize}
\item \zh{被} marque d'abord une action \textbf{subie}, souvent défavorable ; l'usage neutre ou favorable (\zh{被选为}, \zh{被录取}) est courant en chinois moderne.
\item Structure rigide : \textbf{Patient (défini) + 被 + (Agent) + Verbe transitif + Complément obligatoire}.
\item Pour un simple constat de résultat sans agent : \textbf{Passif sans marque} (Patient + V + Complément + 了).
\item Pour les actions favorables formelles : \textbf{受到/得到 + Nom verbal}.
\item \textbf{Interdits avec 被 :} verbes d'état/sentiment/connaissance (\zh{爱}, \zh{知道}, \zh{有}…), verbes intransitifs.
\item \textbf{Piège francophone :} ne pas calquer le passif français systématique $\rightarrow$ privilégier l'actif ou le passif sans marque.
\end{itemize}
\end{aretenir}

\courssub{Exercices d'application et de transformation}

\begin{exercice}[Transformation : Actif $\rightarrow$ 被字句]
Transformez les phrases actives en phrases en \zh{被} (gardez l'agent si pertinent, ajoutez le complément nécessaire). Mettez le pinyin et la traduction.
\begin{enumerate}
\item 爸爸修好了自行车。 \emph{Bàba xiū hǎo le zìxíngchē.}
\item 小偷偷走了我的手机。 \emph{Xiǎotōu tōu zǒu le wǒ de shǒujī.}
\item 老师批评了那个学生。 \emph{Lǎoshī pīpíng le nà ge xuésheng.}
\item 风吹倒了树。 \emph{Fēng chuīdǎo le shù.}
\item 他们选他当班长。 \emph{Tāmen xuǎn tā dāng bānzhǎng.}
\end{enumerate}
\end{exercice}

\begin{exercice}[Choix de la bonne tournure passive]
Choisissez la structure adaptée (\zh{被}, passif sans marque, \zh{受到}, phrase active) et écrivez la phrase chinoise complète (caractères + pinyin).
\begin{enumerate}
\item « La fenêtre a été ouverte (par quelqu'un / par le vent). » (Contexte : on entre dans la classe, la fenêtre est grande ouverte, on ne sait pas par qui.)
\item « Il a reçu une éducation rigoureuse. » (Registre formel.)
\item « Ce gâteau a été mangé par mon petit frère. » (Action subie, agent connu.)
\item « Le livre est écrit en chinois. » (Description d'état/résultat, pas d'agent.)
\item « Elle a été admise à l'Université de Pékin. » (Fait favorable, usage moderne.)
\end{enumerate}
\end{exercice}

\begin{exercice}[Traduction thème (Français $\rightarrow$ Chinois)]
Traduisez en chinois (caractères + pinyin).
\begin{enumerate}
\item Mon dictionnaire a été perdu. (Agent omis.)
\item La lettre a été écrite par ma mère.
\item Notre équipe a été battue par l'équipe adverse.
\item Il a été élu président de l'association des étudiants.
\item Ce problème a été résolu. (Constat de résultat.)
\item Les délégués ont reçu un accueil chaleureux. (Registre formel.)
\item Ce plat ne se mange pas froid. (Éviter le passif !)
\end{enumerate}
\end{exercice}

\begin{corrige}[Exercice 1 : Transformation Actif $\rightarrow$ 被]
\begin{enumerate}
\item \zh{自行车被爸爸修好了。} \emph{Zìxíngchē bèi bàba xiū hǎo le.} — Le vélo a été réparé par papa.
\item \zh{我的手机被小偷偷走了。} \emph{Wǒ de shǒujī bèi xiǎotōu tōu zǒu le.} — Mon téléphone a été volé par un voleur.
\item \zh{那个学生被老师批评了。} \emph{Nà ge xuésheng bèi lǎoshī pīpíng le.} — Cet élève a été critiqué par le prof.
\item \zh{树被风吹倒了。} \emph{Shù bèi fēng chuīdǎo le.} — L'arbre a été renversé par le vent.
\item \zh{他被他们选为班长。} \emph{Tā bèi tāmen xuǎn wéi bānzhǎng.} — Il a été élu délégué par eux. (Ou \zh{他被选为班长。} agent omis).
\end{enumerate}
\end{corrige}

\begin{corrige}[Exercice 2 : Choix de la tournure]
\begin{enumerate}
\item Contexte : constat de résultat, agent inconnu/irrelevant $\rightarrow$ \textbf{Passif sans marque} : \zh{窗户打开了。} \emph{Chuānghu dǎkāi le.} (Ou \zh{窗户被打开了。} si on insiste sur l'action subie, mais \zh{打开了} est plus naturel pour un état résultatif).
\item Action favorable formelle $\rightarrow$ \textbf{受到} : \zh{他受到了严格的教育。} \emph{Tā shòudào le yángé de jiàoyù.}
\item Action subie défavorable, agent connu $\rightarrow$ \textbf{被} : \zh{这个蛋糕被弟弟吃了。} \emph{Zhè ge dàngāo bèi dìdi chī le.}
\item Description d'état/résultat, pas d'agent $\rightarrow$ \textbf{Passif sans marque} (ou \zh{是…写的}) : \zh{这本书是用中文写的。} \emph{Zhè běn shū shì yòng Zhōngwén xiě de.} (Mieux que \zh{书写好了} qui implique « le livre a fini d'être écrit »).
\item Fait favorable, usage moderne $\rightarrow$ \textbf{被} : \zh{她被北京大学录取了。} \emph{Tā bèi Běijīng Dàxué lùqǔ le.}
\end{enumerate}
\end{corrige}

\begin{corrige}[Exercice 3 : Traduction thème]
\begin{enumerate}
\item \zh{我的字典丢了。} \emph{Wǒ de zìdiǎn diū le.} (Verbe intransitif \zh{丢} « se perdre, disparaître » $\rightarrow$ pas de \zh{被}. Active : « Mon dico a disparu ».) \textbf{Attention :} \zh{被丢了} est incorrect.
\item \zh{信被妈妈写好了。} \emph{Xìn bèi māma xiě hǎo le.} (Ou \zh{信是妈妈写的。} \emph{Xìn shì māma xiě de.} construction \zh{是…的} focus agent).
\item \zh{我们队被对手打败了。} \emph{Wǒmen duì bèi duìshǒu dǎbài le.}
\item \zh{他被选为学生会会长。} \emph{Tā bèi xuǎn wéi xuéshēnghuì huìzhǎng.}
\item \zh{这个问题解决了。} \emph{Zhè ge wèntí jiějué le.} (Passif sans marque, verbe composé \zh{解决} + \zh{了}).
\item \zh{代表们受到了热烈欢迎。} \emph{Dàibiǎomen shòudào le rèliè huānyíng.}
\item \zh{这个菜不冷吃。} \emph{Zhè ge cài bù lěng chī.} Ou \zh{这个菜要热着吃。} \emph{Zhè ge cài yào rèzhe chī.} (Active impersonnelle / modalité).
\end{enumerate}
\end{corrige}

\coursec{Leçon 18 — Grammaire : 被字句 (La construction passive avec \zh{被})}

\begin{textebox}[Situation de communication : Un incident à Yaoundé]
\textbf{Contexte :} Paul (élève à Yaoundé) et Li Ming (étudiant chinois en échange) discutent près du lycée après la classe.

\medskip
\noindent
\textbf{Paul :} \zh{哎，李明，我的自行车被偷走了！} \\ 
\emph{Āi, Lǐmíng, wǒ de zìxíngchē bèi tōuzǒu le!} \\
« Eh, Li Ming, mon vélo a été volé ! »

\medskip
\noindent
\textbf{Li Ming :} \zh{真的吗？有没有被找到？} \\ 
\emph{Zhēn de ma? Yǒu méi yǒu bèi zhǎodào?} \\
« Vraiment ? A-t-il été retrouvé ? »

\medskip
\noindent
\textbf{Paul :} \zh{没被找到。警察说：\textquotesingle 这事儿被记录了，会被调查的。\textquotesingle} \\ 
\emph{Méi bèi zhǎodào. Jǐngchá shuō: \textquotesingle Zhè shìr bèi jìlù le, huì bèi diàochá de.\textquotesingle} \\
« Pas retrouvé. La police a dit : "Cette affaire a été enregistrée, elle sera enquêtée." »

\medskip
\noindent
\textbf{Li Ming :} \zh{别难过。下次买个好锁，别让车被偷走。} \\ 
\emph{Bié nánguò. Xià cì mǎi gè hǎo suǒ, bié ràng chē bèi tōuzǒu.} \\
« Ne sois pas triste. La prochaine fois achète un bon antivol, ne laisse pas le vélo se faire voler. »
\end{textebox}

\courssub{Observation d'exemples authentiques}

\begin{exemplebox}[Analyse des phrases du dialogue]
Relevons les phrases en \zh{被} du dialogue et observons leur structure :
\begin{enumerate}
\item \zh{我的自行车被偷走了。} \emph{Wǒ de zìxíngchē bèi tōuzǒu le.} (Mon vélo a été volé.)
\item \zh{有没有被找到？} \emph{Yǒu méi yǒu bèi zhǎodào?} (A-t-il été retrouvé ?)
\item \zh{没被找到。} \emph{Méi bèi zhǎodào.} (Pas été retrouvé.)
\item \zh{这事儿被记录了。} \emph{Zhè shìr bèi jìlù le.} (Cette affaire a été enregistrée.)
\item \zh{会被调查的。} \emph{Huì bèi diàochá de.} (Sera enquêtée.)
\item \zh{别让车被偷走。} \emph{Bié ràng chē bèi tōuzǒu.} (Ne laisse pas le vélo se faire voler — \zh{让} + passif.)
\end{enumerate}
\textbf{Constat :} Le sujet de la phrase (\zh{自行车}, \zh{这事儿}, \zh{车}) \emph{subit} l'action. Le marqueur \zh{被} précède le verbe. L'agent (le voleur, la police) est souvent absent.
\end{exemplebox}

\courssub{Schéma de la structure et variations}

\begin{propriete}[Structure canonique du \zh{被}字句]
La phrase passive marquée par \zh{被} suit le schéma rigide suivant :

\medskip
\centerline{\textbf{Patient (Sujet) + 被 + Agent (optionnel) + Verbe + Complément obligatoire}}
\medskip

\begin{center}
\begin{tabularx}{\linewidth}{|X|X|X|X|}
\hline
\rowcolor{popblueL} \textbf{Patient} & \textbf{被} & \textbf{Agent} & \textbf{Verbe + Complément} \\
\hline
\zh{我的自行车} & \zh{被} & \zh{(小偷)} & \zh{偷走了} \\
\hline
\zh{这事儿} & \zh{被} & \zh{警察} & \zh{记录了} \\
\hline
\zh{窗户} & \zh{被} & \zh{风} & \zh{吹开了} \\
\hline
\end{tabularx}
\end{center}

\textbf{Règles de base :}
\begin{itemize}
\item \zh{被} (\emph{bèi}) est une préposition marquant l'agent (celui qui fait l'action). Elle se place \emph{après le sujet patient} et \emph{avant l'agent}.
\item L'agent peut être omis si le contexte le rend clair ou s'il est inconnu/indéfini (\zh{被偷走了}).
\item \textbf{Point crucial :} Le verbe \emph{ne peut pas rester nu} (surtout monosyllabique). Il doit être suivi d'un complément : \zh{了} (aspect accompli), complément de résultat (\zh{好, 完, 破, 走, 掉, 到}), complément de direction (\zh{出去, 进来}), ou complément de durée/quantité.
\end{itemize}
\end{propriete}

\begin{propriete}[Variations : Négation, Auxiliaires, Aspect]
\begin{center}
\begin{tabularx}{\linewidth}{|X|X|X|}
\hline
\rowcolor{popblueL} \textbf{Type} & \textbf{Structure} & \textbf{Exemple} \\
\hline
Négation passé & Sujet + \zh{没(有)} + \zh{被} + Agent + V + Compl. & \zh{他没被邀请。} \emph{Tā méi bèi yāoqǐng.} \\
\hline
Négation habitude/futur & Sujet + \zh{不} + \zh{被} + Agent + V + Compl. & \zh{这本书不被重视。} \emph{Zhè běn shū bù bèi zhòngshì.} \\
\hline
Auxiliaire (futur/probabilité) & Sujet + Aux. (\zh{会/要/应该/可能}) + \zh{被} + Agent + V + Compl. & \zh{作业可能会被检查。} \emph{Zuòyè kěnéng huì bèi jiǎnchá.} \\
\hline
Aspect accompli & Sujet + \zh{被} + Agent + V + \zh{了} / Résultat & \zh{信被寄出去了。} \emph{Xìn bèi jì chūqù le.} \\
\hline
Aspect expérientiel & Sujet + \zh{被} + Agent + V + \zh{过} + Compl. & \zh{我被骗过一次。} \emph{Wǒ bèi piàn guò yī cì.} \\
\hline
\end{tabularx}
\end{center}

\end{propriete}

\begin{attention}[Ordre impératif : Négation / Auxiliaire AVANT 被]
Jamais : \zh{✗ 他被没邀请。} \quad Toujours : \zh{✓ 他没被邀请。} \\
Jamais : \zh{✗ 他被会批评。} \quad Toujours : \zh{✓ 他会被批评。}
\end{attention}

\courssub{Emplois, restrictions et exceptions}

\begin{propriete}[Quand l'agent est-il obligatoire ?]
L'agent \emph{doit} être exprimé (ou sous-entendu très précis) si le verbe seul crée une ambiguïté ou si le verbe est \emph{intransitif} utilisé de façon causative.
\begin{itemize}
\item \zh{他被狗咬了。} \emph{Tā bèi gǒu yǎo le.} (Il a été mordu par le chien — \zh{咬} transitif, agent clair).
\item \zh{树被风吹倒了。} \emph{Shù bèi fēng chuīdǎo le.} (L'arbre a été renversé par le vent — causatif).
\item \emph{Mais :} \zh{自行车被偷走了。} (Agent inconnu/indéfini « un voleur » → omission naturelle).
\end{itemize}
\end{propriete}

\begin{propriete}[Verbes incompatibles avec 被 (Exceptions majeures)]
\zh{被} ne s'emploie \emph{pas} avec :
\begin{enumerate}
\item \textbf{Verbes d'état / psychologiques :} \zh{有, 是, 像, 爱, 喜欢, 讨厌, 知道, 认识, 理解}.\quad \zh{✗ 书被我有。} \zh{✗ 他被大家爱。}
\item \textbf{Verbes à double complément (donner) :} \zh{给, 送, 借, 还, 教, 告诉} (souvent). On préfère l'actif : \zh{他给了我一本书。} ou \zh{我被给了一本书} (rare, maladroit).
\item \textbf{Verbes intransitifs sans causative :} \zh{去, 来, 走, 死, 睡觉}.\quad \zh{✗ 他被死了。} (Sauf sens causative rare : \zh{被气死} « mis en colère à mort »).
\end{enumerate}
\end{propriete}

\begin{savaistu}[被 vs 叫 / 让 (Passif oral)]
À l'oral (Nord de la Chine surtout), \zh{叫} (\emph{jiào}) et \zh{让} (\emph{ràng}) remplacent \zh{被} :
\zh{我叫他骂了一顿。} \emph{Wǒ jiào tā mà le yī dùn.} « Je me suis fait engueuler par lui. »
\zh{别让自行车被偷了。} \emph{Bié ràng zìxíngchē bèi tōu le.} (Mélange \zh{让} + \zh{被}).
\textbf{Au écrit formel (examen, rédaction) : n'utilisez que \zh{被}.}
\end{savaistu}

\begin{propriete}[被 vs 受到 / 得到 (Passif favorable/neutre)]
Comme anticipé dans l'observation, \zh{被} porte une nuance \emph{adversative} (subir un malheur).
\begin{center}
\begin{tabularx}{\linewidth}{|X|X|X|}
\hline
\rowcolor{popblueL} \textbf{Nuance} & \textbf{Marque} & \textbf{Exemple} \\
\hline
Défavorable / Subi & \zh{被} & \zh{手机被偷了。} (Volé), \zh{被批评} (Critiqué), \zh{被雨淋湿} (Mouillé par la pluie) \\
\hline
Favorable (honneur, prix) & \zh{受到} + Nom abstrait & \zh{受到表扬} (Félicité), \zh{受到欢迎} (Accueilli chaleureusement), \zh{受到奖励} (Récompensé) \\
\hline
Obtention (droit, chance) & \zh{得到} + Nom & \zh{得到提拔} (Promu), \zh{得到锻炼} (A eu l'occasion de s'exercer) \\
\hline
Neutre / Résultat (pas d'agent) & Passif sans marque (SV) & \zh{作业写完了。} (Devoirs finis), \zh{饭做好了。} (Riz cuit), \zh{门修好了。} (Porte réparée) \\
\hline
\end{tabularx}
\end{center}
\end{propriete}

\coursec{Comparaison avec le français et pièges}

Dans les sections précédentes, nous avons établi le schéma du \zh{被}字句, ses variations (négation, auxiliaires, aspect) et ses restrictions sémantiques (verbes d'état, nuance adversative). Il reste maintenant à confronter cette structure à notre langue de départ, le français. C'est là que se joue l'essentiel des erreurs des élèves camerounais : le français possède un passif très fréquent, formé sur l'auxiliaire \emph{être}, et la tentation est grande de le « calquer » mot à mot en chinois. Or les deux systèmes ne se recouvrent que partiellement. Observons d'abord, puis dégageons les règles.

\courssub{Deux passifs, deux logiques}

\begin{exemplebox}[Observation : même sens, deux mécanismes]
Comparons les paires suivantes issues de situations camerounaises ou scolaires :
\begin{itemize}
\item \zh{手机被摔坏了。} \emph{Shǒujī bèi shuāihuài le.} « Le téléphone a été cassé (en tombant). »
\item \zh{他被老师批评了。} \emph{Tā bèi lǎoshī pīpíng le.} « Il a été critiqué par le professeur. »
\item \zh{我的自行车被修好了。} \emph{Wǒ de zìxíngchē bèi xiūhǎo le.} « Mon vélo a été réparé. »
\item \zh{他没被邀请。} \emph{Tā méi bèi yāoqǐng.} « Il n'a pas été invité. »
\end{itemize}
On observe : le français dit « \emph{être} + participe passé » (a été cassé, a été critiqué) avec accord ; le chinois dit simplement \cle{被} + verbe actif, sans verbe « être », sans accord, et \emph{avec un complément après le verbe} (\zh{了, 好, 掉}).
\end{exemplebox}

\begin{propriete}[Règle fondamentale : 被 n'est pas « être »]
Le passif français se construit avec l'auxiliaire \emph{être} + participe passé \emph{accordé} avec le sujet : « la porte \textbf{est} ferm\textbf{ée} », « les livres \textbf{ont été} rend\textbf{us} ».

Le passif chinois se construit avec la marque \cle{被} placée \emph{devant un verbe actif}, qui ne change jamais de forme :

\medskip
\centerline{\textbf{Patient + 被 (+ agent) + Verbe actif + Complément}}
\medskip

Conséquences directes pour le francophone :
\begin{itemize}
\item \textbf{Pas d'auxiliaire :} on ne met \emph{jamais} \zh{是} (\emph{shì}) devant \zh{被} pour former le passif.
\item \textbf{Pas d'accord :} le verbe chinois est invariable ; « cassé, cassée, cassés, cassées » se disent tous \zh{打破} (\emph{dǎpò}) ou \zh{摔坏} (\emph{shuāihuài}).
\item \textbf{Temps et aspect} se marquent par \zh{了}, \zh{过}, \zh{没} autour du verbe, pas par une conjugaison.
\end{itemize}
\end{propriete}

\begin{attention}[Piège 1 — Ne traduis jamais « être » par 是 dans un passif]
C'est l'erreur numéro un des francophones. « Le téléphone a été cassé » ne se traduit \emph{pas} par :

\medskip
\centerline{\zh{✗ 手机是被摔坏了。} \quad (calque de « le téléphone \textbf{est} cassé »)}
\medskip

La forme correcte est :

\medskip
\centerline{\zh{✓ 手机被摔坏了。} \emph{Shǒujī bèi shuāihuài le.} « Le téléphone a été cassé. »}
\medskip

\zh{被} seul suffit à marquer le passif. De même : \zh{他被打了。} \emph{Tā bèi dǎ le.} = « Il a été frappé » (et non \zh{✗ 他是被打了}).

\medskip
\emph{Nuance importante :} la structure \zh{是……的} existe en chinois, mais elle sert à \emph{mettre en relief} une circonstance (qui, quand, comment, par qui), pas à former le passif : \zh{这本书是在北京买的。} \emph{Zhè běn shū shì zài Běijīng mǎi de.} « C'est à Pékin que ce livre a été acheté (c'est là qu'on l'a acheté). »
\end{attention}

\courssub{Passif sans agent : un point commun, mais une exigence chinoise}

\begin{propriete}[L'agent peut être omis dans les deux langues]
Le français dit : « La vitre a été brisée » (sans dire par qui). Le chinois le permet aussi : \zh{被} peut être suivi \emph{directement du verbe}, sans agent exprimé :

\medskip
\centerline{\textbf{Patient + 被 + Verbe + Complément}}
\medskip

\begin{itemize}
\item \zh{钱包被偷走了。} \emph{Qiánbāo bèi tōuzǒu le.} « Le portefeuille a été volé. »
\item \zh{那封信被寄出去了。} \emph{Nà fēng xìn bèi jì chūqù le.} « La lettre a été envoyée. »
\end{itemize}
\textbf{Mais attention :} même sans agent, le chinois exige presque toujours un \emph{complément} après le verbe (voir Piège 2 ci-dessous), alors que le français accepte un participe seul (« la vitre a été brisée »).
\end{propriete}

\courssub{Faits favorables : le français passe au passif, le chinois hésite}

\begin{propriete}[被 garde une couleur « adversative »]
Le français met au passif n'importe quel événement, heureux ou malheureux : « Il a été élu », « Elle a été récompensée », « J'ai été invité ».

Le chinois, lui, réserve traditionnellement \zh{被} aux événements \emph{subis}, désagréables ou inattendus : être volé, critiqué, trompé, mouillé par la pluie, renversé par une voiture. Pour un fait favorable ou neutre, le chinois préfère d'autres tournures :

\begin{itemize}
\item \textbf{Passif sans marque (sujet patient + verbe + complément) :} \zh{作业写完了。} \emph{Zuòyè xiěwán le.} « Les devoirs sont finis (ont été finis). »
\item \textbf{Marques \zh{受到} (\emph{shòudào}) ou \zh{得到} (\emph{dédào}) :} \zh{他受到老师的表扬。} \emph{Tā shòudào lǎoshī de biǎoyáng.} « Il a été félicité par le professeur. » ; \zh{他得到提拔。} « Il a été promu. »
\item \textbf{Tournure active avec \zh{让} / \zh{叫} (oral) ou phrase active simple :} \zh{同学们选他当班长。} \emph{Tóngxuémen xuǎn tā dāng bānzhǎng.} « Les camarades l'ont élu chef de classe. »
\end{itemize}
Sous l'influence des langues européennes, le chinois moderne écrit (médias, administration) accepte de plus en plus \zh{被} pour des faits positifs (\zh{他被选为班长。} « Il a été élu chef de classe »), mais \textbf{au lycée, retiens la norme standard} : \emph{événement subi et défavorable = 被 ; événement neutre ou favorable = autre tournure}.
\end{propriete}

\begin{savaistu}[Les francophones « abusent » de 被]
Les statistiques des corpus d'apprenants (HSK, corpus universitaires) le montrent : les étudiants francophones (et anglophones) \emph{surutilisent} \zh{被}, par calque du passif de leur langue. Là où un Français dit spontanément « le repas a été préparé, la salle a été décorée, les invités ont été accueillis », un Chinois dira naturellement \zh{饭做好了}, \zh{房间布置好了}, \zh{客人来了} — des passifs sans marque ou des phrases actives. Un texte chinois authentique n'emploie \zh{被} qu'avec parcimonie, quand l'événement est vraiment subi. En rédaction, pose-toi toujours la question : « Un Chinois dirait-il vraiment \zh{被} ici ? »
\end{savaistu}

\courssub{Les quatre pièges majeurs des francophones}

\begin{attention}[Piège 2 — Le verbe seul ne suffit pas : il faut un complément]
En français, « le livre a été emprunté » est complet avec le seul participe. En chinois, un verbe monosyllabique nu après \zh{被} est \emph{grammaticalement incomplet} (sauf verbes dissyllabiques bien établis comme \zh{邀请}, \zh{批评}, \zh{表扬} qui peuvent parfois apparaître seuls dans l'écrit formel, mais \zh{了} reste préférable).

\medskip
\centerline{\zh{✗ 书被借。} \qquad \zh{✓ 书被借走了。} \emph{Shū bèi jièzǒu le.} « Le livre a été emprunté. »}
\medskip

Il faut ajouter après le verbe un élément qui précise le résultat ou l'achèvement :
\begin{itemize}
\item \zh{了} (aspect accompli) : \zh{门被关上了。} \emph{Mén bèi guānshàng le.}
\item Complément de résultat : \zh{好, 完, 破, 掉, 走, 到, 住, 开}…
\item Complément directionnel : \zh{出去, 进来, 上来, 下去}…
\item Complément de durée/quantité : \zh{被关了两个小时。}
\end{itemize}
Autres exemples corrigés :
\begin{itemize}
\item \zh{✗ 门被关。} → \zh{✓ 门被关上了。} « La porte a été fermée. »
\item \zh{✗ 饭被吃。} → \zh{✓ 饭被吃完了。} « Le riz a été mangé (tout fini). »
\item \zh{✗ 信被寄。} → \zh{✓ 信被寄出去了。} « La lettre a été envoyée. »
\end{itemize}
\end{attention}

\begin{attention}[Piège 3 — La négation se place AVANT 被, jamais après]
En français, la négation encadre l'auxiliaire : « il n'a \textbf{pas été} critiqué ». Le francophone place alors naturellement la négation après \zh{被}… et se trompe :

\medskip
\centerline{\zh{✗ 他被没批评。} \qquad \zh{✓ 他没被批评。} \emph{Tā méi bèi pīpíng.} « Il n'a pas été critiqué. »}
\medskip

Règle absolue : \cle{没(有)} (pour le passé/accompli) ou \cle{不} (pour l'habitude, le futur, l'état) précèdent toujours \zh{被} :

\medskip
\centerline{\textbf{Sujet + 没 / 不 + 被 (+ agent) + Verbe + Complément}}
\medskip

\begin{itemize}
\item \zh{他没被邀请。} \emph{Tā méi bèi yāoqǐng.} « Il n'a pas été invité (passé). »
\item \zh{这个地方不被人知道。} \emph{Zhège dìfang bù bèi rén zhīdào.} « Cet endroit n'est pas connu des gens (habitude/état). »
\end{itemize}
Même règle pour les auxiliaires modaux (\zh{会}, \zh{应该}, \zh{可能}, \zh{能}) : ils se placent \emph{avant} \zh{被} : \zh{他可能会被批评。} \emph{Tā kěnéng huì bèi pīpíng.} « Il risque d'être critiqué. »
\end{attention}

\begin{attention}[Piège 4 — Un état n'est pas une action : pas de 被 pour les états]
Le français utilise la forme « être + participe » (souvent adjectival) pour décrire de simples \emph{états} : « la porte est ouverte », « la fenêtre est fermée », « le magasin est ouvert ». Ici, aucune action n'est en cours ni accomplie récemment : on décrit un résultat stable. Le chinois n'emploie alors \emph{pas} \zh{被} :

\medskip
\centerline{\zh{✗ 门被开着。} \qquad \zh{✓ 门开着。} \emph{Mén kāizhe.} « La porte est ouverte. »}
\medskip

Pour les états résultatifs, le chinois utilise la structure \zh{verbe + 着} (\emph{zhe}) ou un adjectif prédicatif :

\begin{itemize}
\item \zh{窗户关着。} \emph{Chuānghu guānzhe.} « La fenêtre est fermée. » (État)
\item \zh{商店开着门。} \emph{Shāngdiàn kāizhe mén.} « Le magasin est ouvert (la porte est ouverte). »
\item \zh{灯亮着。} \emph{Dēng liàngzhe.} « La lumière est allumée. »
\end{itemize}
Réserve \zh{被} aux \emph{actions subies} (événement ponctuel) : \zh{门被风吹开了。} \emph{Mén bèi fēng chuīkāi le.} « La porte a été ouverte par le vent » (action !).
\end{attention}

\courssub{Tableau récapitulatif : français ↔ chinois (Bonnes et mauvaises traductions)}

\begin{center}
\begin{tabularx}{\linewidth}{|X|X|X|}
\hline
\rowcolor{popblueL} \textbf{Phrase française} & \textbf{Chinois correct} & \textbf{Traduction fautive (calque)} \\
\hline
Le verre a été cassé. & \zh{杯子被打破了。} & \zh{✗ 杯子是被打破了。} \\
\hline
Il n'a pas été invité. & \zh{他没被邀请。} & \zh{✗ 他被没邀请。} \\
\hline
Le livre a été emprunté. & \zh{书被借走了。} & \zh{✗ 书被借。} \\
\hline
La porte est ouverte. (état) & \zh{门开着。} & \zh{✗ 门被开着。} \\
\hline
Mon vélo a été réparé. & \zh{我的自行车被修好了。} & \zh{✗ 我的自行车是被修好。} \\
\hline
Il a été félicité. (favorable) & \zh{他受到老师的表扬。} & \zh{△ 他被老师表扬了。} \\
\hline
Les devoirs ont été finis. (neutre) & \zh{作业写完了。} & \zh{△ 作业被写完了。} \\
\hline
\end{tabularx}
\end{center}

\begin{popfigure}
\begin{tikzpicture}[
fr/.style={rectangle, rounded corners=3pt, draw=popblue, fill=popblueL, align=center, minimum width=5.5cm, minimum height=1.1cm, font=\small},
zhg/.style={rectangle, rounded corners=3pt, draw=popgreen, fill=popgreenL, align=center, minimum width=5.5cm, minimum height=1.1cm, font=\small},
zhb/.style={rectangle, rounded corners=3pt, draw=popink, fill=poppinkL, align=center, minimum width=5.5cm, minimum height=1.1cm, font=\small},
fleche/.style={-{Stealth[length=3mm]}, thick, popdark},
flechebad/.style={-{Stealth[length=3mm]}, thick, popink, dashed}]
\node[fr] (f1) at (0,0) {Le verre a été cassé.};
\node[zhg] (z1) at (8.5,0) {杯子被打破了。};
\draw[fleche] (f1) -- (z1);
\node[fr] (f2) at (0,-1.8) {Il n'a pas été invité.};
\node[zhg] (z2) at (8.5,-1.8) {他没被邀请。};
\draw[fleche] (f2) -- (z2);
\node[fr] (f3) at (0,-3.6) {La porte est ouverte.};
\node[zhg] (z3) at (8.5,-3.6) {门开着。};
\draw[fleche] (f3) -- (z3);
\node[zhb, align=center] (b1) at (8.5,-5.4){✗ 杯子是被打破了。\\ ✗ 门被开着。};
\node[font=\small\itshape, popmuted, align=center] at (0,-5.4) {Calques du français\\ à bannir};
\draw[flechebad] (0,-4.6) -- (b1.west);
\end{tikzpicture}
\legende{De la phrase française passive à l'équivalent chinois correct (flèches vertes pleines) ; en rose pointillé, les traductions mot à mot fautives.}
\end{popfigure}

\courssub{Méthode : traduire un passif français en chinois}

\begin{methode}[Traduire « être + participe passé » en quatre étapes]
\begin{enumerate}
\item \textbf{Repérer le patient et l'agent.} Dans « le vélo a été réparé par le mécanicien », le patient est \emph{le vélo} (\zh{我的自行车}), l'agent \emph{le mécanicien} (\zh{修理工}). Le patient devient le sujet de la phrase chinoise.
\item \textbf{Choisir la tournure.} L'événement est-il subi/défavorable (vol, casse, critique, pluie) ? $\rightarrow$ \zh{被}. Neutre ou résultatif sans agent (devoirs finis, riz cuit) ? $\rightarrow$ passif sans marque (\zh{作业写完了。}). Favorable (éloge, prix, élection) ? $\rightarrow$ \zh{受到}/\zh{得到} ou phrase active. Simple état (porte ouverte, lumière allumée) ? $\rightarrow$ \zh{verbe + 着}.
\item \textbf{Construire le squelette.} Patient + \zh{被} (+ agent) + verbe. Placer la négation (\zh{没}/\zh{不}) ou l'auxiliaire (\zh{会}/\zh{应该}) \emph{avant} \zh{被}. Ne jamais ajouter \zh{是}.
\item \textbf{Compléter le verbe.} Ajouter le complément obligatoire : \zh{了}, résultat (\zh{好, 完, 走, 破, 掉, 到, 住}), direction (\zh{出去, 回来})… Vérifier : un verbe monosyllabique nu après \zh{被} est presque toujours fautif.
\end{enumerate}
\end{methode}

\begin{exoresolu}[Traduction guidée]
Énoncé : traduis en chinois les deux phrases suivantes. (a) « Mon vélo a été réparé. » (b) « Il n'a pas été invité à la fête. »

\tcblower
Solution détaillée :

(a) \textbf{Étape 1 :} patient = \emph{mon vélo} \zh{我的自行车} ; agent non exprimé (contexte : au garage). \textbf{Étape 2 :} événement subi (le vélo était en panne) $\rightarrow$ \zh{被}, sans agent. \textbf{Étape 3 :} squelette \zh{我的自行车被修…}. \textbf{Étape 4 :} le verbe \zh{修} (monosyllabique) seul est insuffisant ; on ajoute le complément de résultat \zh{好} + \zh{了} (accompli).

\medskip
\centerline{\zh{我的自行车被修好了。} \emph{Wǒ de zìxíngchē bèi xiūhǎo le.}}
\medskip

(b) \textbf{Étape 1 :} patient = \emph{il} (\zh{他}) ; agent non exprimé. \textbf{Étape 2 :} événement subi (déception) $\rightarrow$ \zh{被}. \textbf{Étape 3 :} phrase négative au passé $\rightarrow$ \zh{没} placé \emph{avant} \zh{被} : \zh{他没被邀请}. \textbf{Étape 4 :} \zh{邀请} est un verbe dissyllabique, il peut se suffire dans la négation ; on ajoute le complément de lieu \zh{参加晚会} si l'on veut préciser.

\medskip
\centerline{\zh{他没被邀请（参加晚会）。} \emph{Tā méi bèi yāoqǐng (cānjiā wǎnhuì).}}
\medskip

\textbf{Erreurs à éviter :} \zh{✗ 他是被修好了} (\zh{是} en trop), \zh{✗ 他被没邀请} (négation mal placée), \zh{✗ 自行车被修} (verbe monosyllabique nu).
\end{exoresolu}

\begin{exercice}[À toi de jouer : transformation et traduction]
\begin{enumerate}
\item Traduis en chinois :
      \begin{enumerate}
      \item La fenêtre a été cassée par le ballon.
      \item Le voleur n'a pas été attrapé.
      \item La porte est fermée. (Attention : état ou action ?)
      \item Ce problème sera résolu.
      \item Il a été élu délégué de classe. (Choisis la meilleure tournure).
      \end{enumerate}
\item Corrige les phrases fautives suivantes (une erreur par phrase) :
      \begin{enumerate}
      \item \zh{✗ 他被没送到医院。}
      \item \zh{✗ 我的手机是被丢了。}
      \item \zh{✗ 饭被吃。}
      \item \zh{✗ 这本书被我有。}
      \item \zh{✗ 树被倒了。}
      \end{enumerate}
\item Transforme les phrases actives en phrases passives avec \zh{被} (si possible, sinon explique pourquoi) :
      \begin{enumerate}
      \item \zh{小偷偷走了我的钱包。}
      \item \zh{老师表扬了他。}
      \item \zh{大家喜欢这个老师。}
      \item \zh{风吹开了门。}
      \item \zh{他死了。}
      \end{enumerate}
\end{enumerate}
\end{exercice}

\begin{corrige}[Exercice « À toi de jouer »]
\begin{enumerate}
\item \begin{enumerate}
      \item \zh{窗户被球打破了。} \emph{Chuānghu bèi qiú dǎpò le.}
      \item \zh{小偷没被抓住。} \emph{Xiǎotōu méi bèi zhuāzhù.} (ou \zh{抓到})
      \item \zh{门关着。} \emph{Mén guānzhe.} (État $\rightarrow$ pas de \zh{被}, structure \zh{V+着}). Si action : \zh{门被风关上了。}
      \item \zh{这个问题会被解决的。} \emph{Zhège wèntí huì bèi jiějué de.}
      \item \zh{他被选为班长。} (Accepté en chinois moderne pour élection) \textbf{ou mieux :} \zh{同学们选他当班长。} (Actif naturel). \zh{受到} ne s'emploie pas pour « élire ».
      \end{enumerate}
\item \begin{enumerate}
      \item \zh{他没被送到医院。} (Négation \zh{没} avant \zh{被}).
      \item \zh{我的手机丢了。} (Verbe intransitif \zh{丢} « se perdre/être perdu » $\rightarrow$ pas de \zh{被}, pas de \zh{是}) \textbf{ou} \zh{我的手机被偷了。} (Si vol : verbe transitif \zh{偷}).
      \item \zh{饭被吃完了。} (Complément \zh{完了} obligatoire après \zh{吃} monosyllabique).
      \item \zh{这本书我有。} (\zh{有} verbe d'état $\rightarrow$ pas de passif \zh{被}). Structure active : Sujet + \zh{有} + Objet.
      \item \zh{树倒了。} (\zh{倒} intransitif « tomber » $\rightarrow$ pas de \zh{被} sans causative). Si causative : \zh{树被风吹倒了。}
      \end{enumerate}
\item \begin{enumerate}
      \item \zh{我的钱包被小偷偷走了。} (OK)
      \item \zh{他被老师表扬了。} (OK, \zh{表扬} dissyllabique, \zh{了} ajouté). Nuance : \zh{他受到老师的表扬。} (Plus formel/poli).
      \item \textbf{Impossible.} \zh{喜欢} verbe psychologique/état $\rightarrow$ pas de \zh{被}. Reste actif : \zh{大家喜欢这个老师。}
      \item \zh{门被风吹开了。} (OK, causative \zh{吹开}).
      \item \textbf{Impossible.} \zh{死} intransitif non causative $\rightarrow$ pas de \zh{被}. (Sauf \zh{被气死} « mis en colère à mort », sens figuré).
      \end{enumerate}
\end{enumerate}
\end{corrige}

\begin{aretenir}[L'essentiel de la leçon 18 — \zh{被}字句]
\begin{itemize}
\item \textbf{Structure :} Patient + \zh{被} (+ Agent) + Verbe + \textbf{Complément obligatoire} (\zh{了}, résultat, direction…).
\item \textbf{Passif $\neq$ Être :} \zh{被} marque le passif tout seul. \textbf{Jamais} \zh{是被}.
\item \textbf{Négation / Auxiliaires :} Se placent \textbf{AVANT} \zh{被} (\zh{没被}, \zh{不被}, \zh{会被}, \zh{应该被}).
\item \textbf{Sémantique :} \zh{被} = événement \emph{subi, défavorable, inattendu}. Favorable $\rightarrow$ \zh{受到}/\zh{得到} / Actif. Neutre/Résultatif $\rightarrow$ Passif sans marque (SV + Complément). État $\rightarrow$ \zh{V+着}.
\item \textbf{Verbes exclus :} États (\zh{有, 是, 爱, 知道}…), intransitifs purs (\zh{去, 死, 睡}), doubles compléments (\zh{给, 送}…).
\item \textbf{Réflexe anti-calque :} Avant d'écrire \zh{被}, demande-toi : « Est-ce un malheur subi ? Le verbe a-t-il un complément ? La négation est-elle avant \zh{被} ? »
\end{itemize}
\end{aretenir}

\coursec{Exercices d'application et de transformation}

Après avoir comparé la voix passive chinoise à la voix passive française et repéré les pièges les plus fréquents (verbe laissé seul, négation mal placée, oubli du complément), il est temps de passer à la pratique. Cette dernière partie de la leçon vous entraîne à \cle{transformer} les phrases (actif $\leftrightarrow$ passif), à \cle{choisir} la bonne particule passive (\zh{被}, \zh{叫}, \zh{让}), à \cle{corriger} des phrases fautives et à \cle{produire} vos propres phrases en \zh{被}字句. C'est exactement le type d'épreuve que vous rencontrerez à l'examen : transformations, thème et rédaction guidée.

\courssub{Rappel de méthode : transformer une phrase active en \zh{被}字句}

\begin{methode}[Transformer : actif $\rightarrow$ \zh{被}字句]
Pour transformer une phrase active du type « Sujet + Verbe + Objet (+ complément) » en phrase passive en \zh{被}, suivez les quatre étapes :
\begin{enumerate}
\item Repérer le \cle{patient} (l'objet qui subit l'action) : il devient le \textbf{sujet} de la nouvelle phrase.
\item Insérer \zh{被} juste après ce nouveau sujet.
\item Placer l'\cle{agent} (l'ancien sujet) juste après \zh{被} : c'est lui qui fait l'action.
\item Garder le verbe \textbf{avec son complément} (résultatif, directionnel, \zh{了}…) : le verbe ne reste jamais seul.
\end{enumerate}
Schéma : \textbf{Patient + \zh{被} + Agent + Verbe + Complément.}
\end{methode}

\begin{propriete}[\zh{把}字句 et \zh{被}字句 : des transformations miroir]
La phrase en \zh{把} et la phrase en \zh{被} sont les deux faces d'une même réalité, construites autour du \textbf{même verbe d'action} et du \textbf{même complément} :
\begin{itemize}
\item \zh{把}字句 : l'\textbf{agent} reste sujet, l'objet (patient) est \textbf{déplacé avant le verbe} grâce à \zh{把}. $\rightarrow$ phrase active.
\item \zh{被}字句 : le \textbf{patient} devient sujet, l'agent est introduit par \zh{被}. $\rightarrow$ phrase passive.
\end{itemize}
Dans les deux cas, le verbe porte obligatoirement un complément ou \zh{了}. On passe donc de l'une à l'autre en échangeant simplement les rôles : \zh{弟弟把杯子打破了。} $\leftrightarrow$ \zh{杯子被弟弟打破了。}
\end{propriete}

\begin{popfigure}
\begin{tikzpicture}[node distance=1.1cm, every node/.style={font=\small}]
\node[draw=popblue, fill=popblueL, rounded corners, align=center, text width=4.6cm] (base) {\textbf{Phrase de base}\\弟弟打破了杯子。\\\emph{Dìdi dǎpò le bēizi.}\\« Le petit frère a cassé le verre. »};
\node[draw=poporange, fill=poporangeL, rounded corners, align=center, text width=4.6cm, below left=1.6cm and -1.2cm of base] (ba) {\textbf{\zh{把}字句 (actif)}\\弟弟把杯子打破了。\\\emph{Dìdi bǎ bēizi dǎpò le.}\\« Le petit frère a cassé le verre. »};
\node[draw=poppurple, fill=poppurpleL, rounded corners, align=center, text width=4.6cm, below right=1.6cm and -1.2cm of base] (bei) {\textbf{\zh{被}字句 (passif)}\\杯子被弟弟打破了。\\\emph{Bēizi bèi dìdi dǎpò le.}\\« Le verre a été cassé par le petit frère. »};
\draw[-{Stealth[length=3mm]}, thick, poporange] (base.south west) -- (ba.north) node[midway, left, align=center, font=\footnotesize]{objet déplacé\\avant le verbe};
\draw[-{Stealth[length=3mm]}, thick, poppurple] (base.south east) -- (bei.north) node[midway, right, align=center, font=\footnotesize]{patient devient\\sujet};
\draw[{Stealth[length=3mm]}-{Stealth[length=3mm]}, thick, popdark] (ba.south) -- (bei.south) node[midway, below, font=\footnotesize]{transformation miroir : agent $\leftrightarrow$ patient};
\end{tikzpicture}
\legende{Le triangle des transformations : phrase de base, \zh{把}字句 et \zh{被}字句 autour du même verbe d'action.}
\end{popfigure}

\begin{exemplebox}[La paire miroir à retenir]
\zh{风把门吹开了。} \emph{Fēng bǎ mén chuīkāi le.} « Le vent a ouvert la porte (en soufflant). »\\
\zh{门被风吹开了。} \emph{Mén bèi fēng chuīkāi le.} « La porte a été ouverte par le vent. »\\
Même verbe \zh{吹} « souffler », même complément résultatif \zh{开} « ouvert », même particule \zh{了} : seuls les rôles de sujet et d'objet sont échangés.
\end{exemplebox}

\begin{attention}[Le verbe ne reste jamais seul]
Dans la \zh{被}字句 comme dans la \zh{把}字句, le verbe doit être \textbf{accompagné} : complément résultatif (\zh{打破}, \zh{吃完}), directionnel (\zh{拿走}), \zh{了}, ou un complément de manière. On ne dit jamais *\zh{杯子被摔} ni *\zh{弟弟把杯子打}. C'est la différence majeure avec le français, où « le verre a été cassé » se suffit à lui-même.
\end{attention}

\courssub{Vocabulaire clé de la leçon}

\begin{definition}[Mots-outils et verbes fréquents pour la \zh{被}字句]
\begin{center}
\begin{tabularx}{\linewidth}{|X|X|X|X|}
\hline
\rowcolor{popblueL} Caractères & Pinyin & Nature & Traduction \\
\hline
被 & bèi & part. & marqueur de la voix passive (écrit, formel) \\
叫 & jiào & part. & marqueur de la voix passive (oral, familier) \\
让 & ràng & part. & marqueur de la voix passive (oral, familier) \\
打破 & dǎpò & v. (res.) & casser (en frappant) \\
吃完 & chīwán & v. (res.) & manger entièrement, finir de manger \\
吹走 & chuīzǒu & v. (res.) & emporter en soufflant \\
吹开 & chuīkāi & v. (res.) & ouvrir en soufflant \\
打扫干净 & dǎsǎo gānjìng & v. (res.) & nettoyer jusqu'à propreté \\
拿走 & názǒu & v. (dir.) & emporter, prendre et partir \\
偷走 & tōuzǒu & v. (dir.) & voler (et emporter) \\
批评 & pīpíng & v. & critiquer, blâmer \\
录取 & lùqǔ & v. & admettre, recruter \\
冲垮 & chōngkuǎ & v. (res.) & emporter, détruire par l'eau (inondation) \\
淋湿 & línshī & v. (res.) & mouiller sous la pluie \\
砍倒 & kǎndǎo & v. (res.) & abattre (arbre), couper et faire tomber \\
修好 & xiūhǎo & v. (res.) & réparer (jusqu'à bon état) \\
看完 & kànwán & v. (res.) & finir de regarder / lire \\
\hline
\end{tabularx}
\end{center}
\end{definition}

\courssub{Écriture des caractères : \zh{被}, \zh{叫}, \zh{让}}

\begin{definition}[Ordre des traits et clés]
\begin{itemize}
\item \zh{被} (bèi) : 10 traits. Clé \zh{衣} (vêtement, radical 145). Ordre : haut \zh{皮} (peau) puis bas \zh{衣}. Haut $\rightarrow$ bas.
\item \zh{叫} (jiào) : 10 traits. Clé \zh{口} (bouche, radical 30). Ordre : \zh{口} à gauche, puis \zh{丩} (jiū) à droite. Gauche $\rightarrow$ droite.
\item \zh{让} (ràng) : 17 traits. Clé \zh{讠} (parole, radical 149 simplifié). Ordre : \zh{讠} à gauche, puis \zh{上} et \zh{寸} à droite. Gauche $\rightarrow$ droite ; horizontal avant vertical pour \zh{上}.
\end{itemize}
\end{definition}

\courssub{Exercice 1 — Transformation : actif $\rightarrow$ passif}

\begin{exoresolu}[Transformer en \zh{被}字句]
Énoncé : transformez la phrase active \zh{弟弟吃了我的蛋糕。} en phrase passive. \tcblower
Solution détaillée :
\begin{enumerate}
\item Patient (objet qui subit) : \zh{我的蛋糕} « mon gâteau » $\rightarrow$ devient sujet.
\item Insertion de \zh{被} : \zh{我的蛋糕被…}
\item Agent (ancien sujet) : \zh{弟弟} $\rightarrow$ \zh{我的蛋糕被弟弟…}
\item Verbe + \zh{了} : \zh{吃了}.
\end{enumerate}
Résultat : \zh{我的蛋糕被弟弟吃了。} \emph{Wǒ de dàngāo bèi dìdi chī le.} « Mon gâteau a été mangé par mon petit frère. »
\end{exoresolu}

\begin{exercice}[À vous : actif $\rightarrow$ passif]
Transformez les phrases suivantes en \zh{被}字句 :
\begin{enumerate}
\item \zh{弟弟写完了作业。}（弟弟 = agent, 作业 = patient）
\item \zh{风吹走了我的帽子。}
\item \zh{同学们打扫干净了教室。}
\item \zh{小偷拿走了他的手机。}
\end{enumerate}
\end{exercice}

\begin{corrige}[Exercice 1 : actif $\rightarrow$ passif]
\begin{enumerate}
\item \zh{作业被弟弟写完了。} \emph{Zuòyè bèi dìdi xiěwán le.} « Les devoirs ont été finis par le petit frère. »
\item \zh{我的帽子被风吹走了。} \emph{Wǒ de màozi bèi fēng chuīzǒu le.} « Mon chapeau a été emporté par le vent. »
\item \zh{教室被同学们打扫干净了。} \emph{Jiàoshì bèi tóngxuémen dǎsǎo gānjìng le.} « La salle de classe a été nettoyée par les élèves. »
\item \zh{他的手机被小偷拿走了。} \emph{Tā de shǒujī bèi xiǎotōu názǒu le.} « Son téléphone a été pris par le voleur. »
\end{enumerate}
\end{corrige}

\courssub{Exercice 2 — Transformation : passif $\rightarrow$ actif}

Pour repasser du passif à l'actif, on fait le chemin inverse : l'agent redevient sujet, et le patient est déplacé avant le verbe grâce à \zh{把}. C'est la transformation miroir vue plus haut.

\begin{exoresolu}[Transformer en \zh{把}字句]
Énoncé : transformez \zh{教室被学生打扫干净了。} en phrase active. \tcblower
Solution détaillée :
\begin{enumerate}
\item Agent (après \zh{被}) : \zh{学生} $\rightarrow$ redevient sujet.
\item Patient (sujet de la phrase passive) : \zh{教室} $\rightarrow$ déplacé après \zh{把}.
\item Verbe + complément inchangés : \zh{打扫干净了}.
\end{enumerate}
Résultat : \zh{学生把教室打扫干净了。} \emph{Xuésheng bǎ jiàoshì dǎsǎo gānjìng le.} « Les élèves ont nettoyé la salle de classe. »
\end{exoresolu}

\begin{exercice}[À vous : passif $\rightarrow$ actif]
Transformez en \zh{把}字句 :
\begin{enumerate}
\item \zh{我的蛋糕被弟弟吃了。}
\item \zh{门被风吹开了。}
\item \zh{杯子被妹妹打破了。}
\end{enumerate}
\end{exercice}

\begin{corrige}[Exercice 2 : passif $\rightarrow$ actif]
\begin{enumerate}
\item \zh{弟弟把我的蛋糕吃了。} « Mon petit frère a mangé mon gâteau. »
\item \zh{风把门吹开了。} « Le vent a ouvert la porte. »
\item \zh{妹妹把杯子打破了。} « Ma petite sœur a cassé le verre. »
\end{enumerate}
\end{corrige}

\courssub{Exercice 3 — Compléter avec \zh{被}, \zh{叫} ou \zh{让}}

Rappel : à l'oral et dans la langue familière, \zh{叫} et \zh{让} remplacent \zh{被}, mais ils exigent \textbf{toujours un agent exprimé} (on ne peut pas dire *\zh{书叫借走了}). \zh{被}, lui, peut être suivi directement du verbe : \zh{书被借走了。}

\begin{exoresolu}[Choisir la bonne particule]
Énoncé : complétez \zh{我的书\_\_\_同学借走了。} \tcblower
Solution : l'agent \zh{同学} est exprimé, donc les trois particules sont possibles. Au registre écrit et à l'examen, on choisit \zh{被} : \zh{我的书被同学借走了。} \emph{Wǒ de shū bèi tóngxué jièzǒu le.} « Mon livre a été emprunté par un camarade. » À l'oral : \zh{我的书叫同学借走了。} ou \zh{我的书让同学借走了。}
\end{exoresolu}

\begin{exercice}[Complétez avec \zh{被}, \zh{叫} ou \zh{让}]
\begin{enumerate}
\item \zh{他的自行车\_\_\_人偷走了。}
\item \zh{窗户\_\_\_风打破了。}
\item \zh{我的作业\_\_\_弟弟画了。}（langue parlée）
\item \zh{他\_\_\_老师批评了。}
\end{enumerate}
\end{exercice}

\begin{corrige}[Exercice 3 : \zh{被}, \zh{叫} ou \zh{让}]
\begin{enumerate}
\item \zh{被}（ou \zh{叫}/\zh{让}）: \zh{他的自行车被人偷走了。} \emph{Tā de zìxíngchē bèi rén tōuzǒu le.} « Son vélo a été volé (par quelqu'un). »
\item \zh{被} : \zh{窗户被风打破了。} \emph{Chuānghu bèi fēng dǎpò le.} « La fenêtre a été cassée par le vent. »
\item \zh{叫} ou \zh{让}（registre oral）: \zh{我的作业叫弟弟画了。} \emph{Wǒ de zuòyè jiào dìdi huà le.} « Mes devoirs ont été faits (dessins) par mon petit frère. »
\item \zh{被} : \zh{他被老师批评了。} \emph{Tā bèi lǎoshī pīpíng le.} « Il a été critiqué par le professeur. »
\end{enumerate}
\end{corrige}

\courssub{Exercice 4 — Correction de phrases fautives}

\begin{attention}[Les deux erreurs classiques des francophones]
\begin{itemize}
\item \textbf{Négation mal placée} : la négation \zh{没(有)} se met \textbf{avant} \zh{被}, jamais après. *\zh{他被没录取} $\rightarrow$ \zh{他没被录取。} « Il n'a pas été admis. »
\item \textbf{Verbe laissé seul} : *\zh{杯子被摔} $\rightarrow$ \zh{杯子被摔破了。} « Le verre a été cassé. » Il faut le complément résultatif \zh{破} et \zh{了}.
\end{itemize}
\end{attention}

\begin{exoresolu}[Corriger les phrases]
Énoncé : corrigez *\zh{他被没录取。} et *\zh{杯子被摔。} \tcblower
Solution :
\begin{enumerate}
\item *\zh{他被没录取。} $\rightarrow$ \zh{他没被录取。} \emph{Tā méi bèi lùqǔ.} « Il n'a pas été admis. » Règle : \zh{没} + \zh{被} + agent + verbe.
\item *\zh{杯子被摔。} $\rightarrow$ \zh{杯子被摔破了。} \emph{Bēizi bèi shuāipò le.} « Le verre a été cassé. » Règle : le verbe ne reste jamais seul ; on ajoute le complément \zh{破} et \zh{了}.
\end{enumerate}
\end{exoresolu}

\begin{exercice}[Corrigez les phrases fautives]
\begin{enumerate}
\item *\zh{我的手机被没偷走。}
\item *\zh{饭被弟弟吃。}
\item *\zh{他被老师批评了没。}
\end{enumerate}
\end{exercice}

\begin{corrige}[Exercice 4 : correction]
\begin{enumerate}
\item \zh{我的手机没被偷走。} \emph{Wǒ de shǒujī méi bèi tōuzǒu.} « Mon téléphone n'a pas été volé. »
\item \zh{饭被弟弟吃完了。} \emph{Fàn bèi dìdi chīwán le.} « Le riz a été fini par mon petit frère. »
\item \zh{他没被老师批评。} \emph{Tā méi bèi lǎoshī pīpíng.} « Il n'a pas été critiqué par le professeur. »
\end{enumerate}
\end{corrige}

\courssub{Exercice 5 — Thème : traduction du français vers le chinois}

\begin{exoresolu}[Thème guidé]
Énoncé : traduisez « La maison a été détruite par la pluie. » et « Mon sac n'a pas été volé. » \tcblower
Solution détaillée :
\begin{enumerate}
\item « La maison » = patient $\rightarrow$ sujet : \zh{房子}. « Par la pluie » = agent : \zh{被雨}. « A été détruite » = verbe + complément : \zh{冲垮了}. Résultat : \zh{房子被雨冲垮了。} \emph{Fángzi bèi yǔ chōngkuǎ le.}
\item Phrase négative : \zh{没} avant \zh{被}. « Mon sac » : \zh{我的包} ; « volé » : \zh{偷走}. Résultat : \zh{我的包没有被偷走。} \emph{Wǒ de bāo méiyǒu bèi tōuzǒu.}
\end{enumerate}
\end{exoresolu}

\begin{exercice}[Thème]
Traduisez en chinois :
\begin{enumerate}
\item « La voiture a été réparée par le mécanicien. »
\item « Les arbres n'ont pas été coupés. »
\item « Le match a été regardé par tous les élèves. »
\end{enumerate}
\end{exercice}

\begin{corrige}[Exercice 5 : thème]
\begin{enumerate}
\item \zh{汽车被修车师傅修好了。} \emph{Qìchē bèi xiūchē shīfu xiūhǎo le.}
\item \zh{树没有被砍倒。} \emph{Shù méiyǒu bèi kǎndǎo.}
\item \zh{比赛被所有学生看完了。} \emph{Bǐsài bèi suǒyǒu xuésheng kànwán le.}
\end{enumerate}
\end{corrige}

\courssub{Exercice 6 — Rédaction guidée : raconter un incident de la journée}

\begin{methode}[Rédiger un court récit avec la \zh{被}字句]
\begin{enumerate}
\item Racontez un petit incident de votre journée (à la maison, au lycée, au marché de Yaoundé ou de Douala) en \textbf{quatre phrases}.
\item Utilisez \textbf{au moins deux phrases en \zh{被}} ; les autres peuvent être actives ou en \zh{把}.
\item Vérifiez chaque phrase passive : patient en sujet, \zh{被} + agent, verbe + complément, négation éventuelle avant \zh{被}.
\end{enumerate}
\end{methode}

\begin{exemplebox}[Modèle de rédaction]
\zh{今天下午，我的雨伞被同学借走了。} \emph{Jīntiān xiàwǔ, wǒ de yǔsǎn bèi tóngxué jièzǒu le.} « Cet après-midi, mon parapluie a été emprunté par un camarade. »\\
\zh{回家的时候，我的衣服被雨淋湿了。} \emph{Huí jiā de shíhou, wǒ de yīfu bèi yǔ línshī le.} « En rentrant à la maison, mes vêtements ont été mouillés par la pluie. »\\
\zh{妈妈把热水准备好了。} \emph{Māma bǎ rèshuǐ zhǔnbèi hǎo le.} « Maman avait préparé de l'eau chaude. »\\
\zh{幸好我没有感冒。} \emph{Xìnghǎo wǒ méiyǒu gǎnmào.} « Heureusement, je n'ai pas attrapé froid. »
\end{exemplebox}

\begin{experience}[Production écrite puis orale]
Rédigez vos quatre phrases, puis échangez votre texte avec un voisin : chacun vérifie les deux phrases en \zh{被} de l'autre (place de \zh{被}, complément du verbe, négation). Lisez ensuite votre récit à la classe ; les camarades doivent repérer et compter les phrases passives entendues.
\end{experience}

\courssub{Tableau récapitulatif et bilan}

\begin{center}
\begin{tabularx}{\linewidth}{|X|X|X|}
\hline
\rowcolor{popblueL} \textbf{Structure} & \textbf{Exemple (caractères + pinyin)} & \textbf{Traduction} \\
\hline
Phrase de base : S + V + O & 弟弟打破了杯子。\emph{Dìdi dǎpò le bēizi.} & Le petit frère a cassé le verre. \\
\hline
\zh{把}字句 : Agent + 把 + Patient + V + compl. & 弟弟把杯子打破了。\emph{Dìdi bǎ bēizi dǎpò le.} & Le petit frère a cassé le verre. (actif) \\
\hline
\zh{被}字句 : Patient + 被 + Agent + V + compl. & 杯子被弟弟打破了。\emph{Bēizi bèi dìdi dǎpò le.} & Le verre a été cassé par le petit frère. (passif) \\
\hline
\zh{被}字句 sans agent & 杯子被打破了。\emph{Bēizi bèi dǎpò le.} & Le verre a été cassé. \\
\hline
\zh{被}字句 négative : 没 + 被 & 杯子没被打破。\emph{Bēizi méi bèi dǎpò.} & Le verre n'a pas été cassé. \\
\hline
Registre oral : \zh{叫} / \zh{让} + agent & 杯子叫弟弟打破了。\emph{Bēizi jiào dìdi dǎpò le.} & Le verre a été cassé par le petit frère. (oral) \\
\hline
\end{tabularx}
\end{center}

\begin{aretenir}[L'essentiel pour l'examen]
\begin{itemize}
\item Transformation actif $\rightarrow$ passif : le patient devient sujet, on insère \zh{被} + agent, le verbe garde son complément.
\item \zh{把}字句 et \zh{被}字句 sont des transformations \cle{miroir} : savoir passer de l'une à l'autre est une compétence d'examen.
\item Le verbe n'est \textbf{jamais seul} : toujours un complément ou \zh{了}.
\item La négation \zh{没(有)} se place \textbf{avant} \zh{被}.
\item \zh{叫} et \zh{让} (oral) exigent un agent exprimé ; \zh{被} peut s'en passer.
\end{itemize}
\end{aretenir}

\newpage

\coursec{Activité d'intégration}

\coursec{Leçon 18 — Grammaire : 被字句}

\courssub{Objectifs de la leçon}

À l'issue de cette leçon, tu seras capable de :
\begin{itemize}
\item identifier et analyser la structure \cle{被字句} (phrase passive avec \zh{被}) dans un texte authentique ;
\item construire correctement des phrases passives : affirmative, \cle{négative} (\zh{没有被}), à \cle{agent omis}, avec modal (\zh{要被}, \zh{会被}) ;
\item respecter les contraintes syntaxiques obligatoires : place de la négation, présence d'un complément après le verbe (\zh{了}, complément de résultat, de direction, de lieu) ;
\item distinguer l'emploi du \zh{被} (événement subi, souvent fâcheux) de la voix active ou de la structure \zh{把} ;
\item transformer une phrase passive en phrase active et inversement ;
\item produire un mini-récit cohérent (5--6 phrases) mobilisant au moins trois formes de \zh{被字句} et le présenter à l'oral avec des tons corrects.
\end{itemize}

\courssub{Observation et découverte : un incident au marché Mokolo}

\begin{textebox}[Témoignage : samedi matin au marché Mokolo]
星期六早上，莫科洛市场很热闹。突然刮起了大风，一个卖西红柿的摊子被风吹倒了，西红柿被撒了一地。一位女商贩的手机被一个孩子碰掉在地上，摔坏了。她的一袋洋葱也不见了：可能被偷走了。那个孩子后来被市场主任批评了。\\
\emph{Xīngqīliù zǎoshang, Mòkēluò shìchǎng hěn rènao. Tūrán guāqǐ le dàfēng, yí ge mài xīhóngshì de tānzi bèi fēng chuīdǎo le, xīhóngshì bèi sā le yídì. Yí wèi nǚ shāngfàn de shǒujī bèi yí ge háizi pēng diào zài dìshang, shuāihuài le. Tā de yí dài yángcōng yě bùjiàn le : kěnéng bèi tōuzǒu le. Nà ge háizi hòulái bèi shìchǎng zhǔrèn pīpíng le.}\\
« Samedi matin, le marché Mokolo était très animé. Soudain, un grand vent s'est levé : un étal de tomates a été renversé par le vent et les tomates ont été répandues par terre. Le téléphone portable d'une vendeuse a été heurté par un enfant, est tombé par terre et s'est cassé. Un de ses sacs d'oignons a aussi disparu : il a probablement été volé. Cet enfant a ensuite été réprimandé par le directeur du marché. »
\end{textebox}

\begin{experience}[Activité guidée : Repérage et analyse (10 min)]
Travail en binôme. Relis le texte ci-dessus.
\begin{enumerate}
\item Souligne toutes les phrases contenant \zh{被}.
\item Pour chaque occurrence, remplis le tableau d'analyse ci-dessous.
\item Vérifie : le verbe qui suit \zh{被} est-il « nu » ou suivi d'un complément ?
\end{enumerate}

\begin{center}
\begin{tabularx}{\linewidth}{|X|X|X|X|}
\hline
\rowcolor{popblueL} Phrase (extrait) & Patient (Sujet) & Agent (après 被) & Verbe + Complément \\
\hline
\zh{摊子被风吹倒了} & \zh{摊子} & \zh{风} & \zh{吹倒了} (V + rés. + \zh{了}) \\
\hline
\zh{西红柿被撒了一地} & \zh{西红柿} & (omis) & \zh{撒了一地} (V + compl. de quantité/lieu) \\
\hline
\zh{手机被一个孩子碰掉在地上} & \zh{手机} & \zh{一个孩子} & \zh{碰掉在地上} (V + compl. de direction/lieu) \\
\hline
\zh{可能被偷走了} & \zh{洋葱} (repris) & (omis) & \zh{偷走了} (V + dir. + \zh{了}) \\
\hline
\zh{被市场主任批评了} & \zh{那个孩子} & \zh{市场主任} & \zh{批评了} (V + \zh{了}) \\
\hline
\end{tabularx}
\end{center}

\textbf{Mise en commun :} On observe que le patient (ce qui subit l'action) est placé en \textbf{sujet} (début de phrase), \zh{被} introduit l'agent (le faiseur), et le verbe est \textbf{toujours complété} (\zh{了}, \zh{倒}, \zh{一地}, \zh{走}, \zh{掉在地上}). L'agent peut être omis (inconnu, indifférent).
\end{experience}

\courssub{Schéma formel et règle grammaticale}

\begin{propriete}[Structure de base du 被字句]
\textbf{Schéma :} \cle{Patient + 被 + (Agent) + Verbe + Complément obligatoire}

\begin{center}
\begin{tabularx}{\linewidth}{|X|X|X|}
\hline
\rowcolor{popblueL} Composant & Rôle & Exemple \\
\hline
Patient & Sujet grammatical (subit l'action) & \zh{手机} \\
\hline
\zh{被} & Marque du passif (préposition verbalisée) & \zh{被} \\
\hline
Agent & Faiseur de l'action (peut être omis) & \zh{孩子} / (omis) \\
\hline
Verbe + Complément & Noyau prédicatif \textbf{jamais nu} & \zh{摔坏了} / \zh{偷走了} / \zh{批评了} \\
\hline
\end{tabularx}
\end{center}

\textbf{Règle d'or :} En chinois, un verbe dans une phrase \zh{被} ne peut \textbf{jamais} rester seul. Il doit être suivi de :
\begin{itemize}
\item \zh{了} (aspect accompli) : \zh{被批评了} ;
\item un \textbf{complément de résultat} : \zh{打破} (casser), \zh{做完} (finir), \zh{看懂} (comprendre en lisant) ;
\item un \textbf{complément de direction} : \zh{拿走} (emporter), \zh{送来} (apporter), \zh{跑掉} (s'enfuir) ;
\item un \textbf{complément de lieu/quantité} : \zh{放在桌上} (mis sur la table), \zh{撒了一地} (répandu par terre).
\end{itemize}
\end{propriete}

\begin{aretenir}[Variantes principales]
\begin{enumerate}
\item \textbf{Affirmative simple :} \zh{作业被老师批改了。} \emph{Zuòyè bèi lǎoshī pīgǎi le.} (Les devoirs ont été corrigés par le prof.)
\item \textbf{Négative :} \zh{没有} (ou \zh{没}) \textbf{avant} \zh{被}. \zh{作业没有被批改。} \emph{Zuòyè méiyǒu bèi pīgǎi.} (Les devoirs n'ont pas été corrigés.) \textbf{Jamais} \zh{被没有} ou \zh{被不}.
\item \textbf{Agent omis :} \zh{自行车被偷走了。} \emph{Zìxíngchē bèi tōuzǒu le.} (Le vélo a été volé.) Agent inconnu ou sans importance.
\item \textbf{Avec modal (futur, obligation, possibilité) :} Le modal précède \zh{被}. \zh{信要被寄走。} \emph{Xìn yào bèi jì zǒu.} (La lettre doit être envoyée.) \zh{他会被选上。} \emph{Tā huì bèi xuǎn shàng.} (Il risque d'être sélectionné / sera probablement sélectionné.)
\item \textbf{被 vs 叫 / 让 :} À l'oral (Nord de la Chine), \zh{叫} (jiào) et \zh{让} (ràng) remplacent souvent \zh{被}. \zh{我叫他骂了一顿。} \emph{Wǒ jiào tā mà le yí dùn.} (Je me suis fait engueuler par lui.) \textbf{À l'écrit et à l'examen : utilise \zh{被}.}
\end{enumerate}
\end{aretenir}

\courssub{Emplois, restrictions et nuances sociolinguistiques}

\begin{propriete}[Emploi : la contrainte sémantique « événement subi »]
Contrairement au passif français (\emph{être + participe passé}) qui est neutre, le \zh{被字句} chinois marque presque toujours une \textbf{affectation négative, un dommage, un imprévu ou une contrainte} subie par le patient.

\begin{center}
\begin{tabularx}{\linewidth}{|X|X|}
\hline
\rowcolor{popblueL} \textbf{Contexte typique (BON usage de 被)} & \textbf{Contexte évité (MAUVAIS usage de 被)} \\
\hline
Perte, vol : \zh{钱包被偷了} & Cadeau, honneur : \zh{*我被送了一份礼物} $\rightarrow$ \zh{别人送了我一份礼物} \\
\hline
Dégât, casse : \zh{杯子被打破了} & Action volontaire positive : \zh{*他被邀请参加晚会} $\rightarrow$ \zh{他受邀参加晚会} / \zh{有人邀请他} \\
\hline
Réprimande, punition : \zh{被批评了}, \zh{被罚款了} & Promotion, élection : \zh{*他被选为班长} $\rightarrow$ \zh{他当选班长了} \\
\hline
Catastrophe naturelle : \zh{房子被风吹倒了} & Situation neutre descriptive : \zh{*这本书被很多人看} $\rightarrow$ \zh{很多人看这本书} \\
\hline
\end{tabularx}
\end{center}

\textbf{Exception :} Verbes purement formels/administratifs sans charge affective : \zh{被录取} (être admis), \zh{被批准} (être approuvé), \zh{被授予} (se voir décerner). Ces emplois sont figés (langue administrative).
\end{propriete}

\begin{attention}[Pièges fréquents des francophones — \textbf{À apprendre par cœur}]
\begin{itemize}
\item \textbf{Place de la négation :} \zh{没有} / \zh{没} \textbf{TOUJOURS} avant \zh{被}.
  \begin{itemize}
  \item \textcolor{red}{Faux :} \zh{手机被没有偷走。} \quad \textcolor{red}{Faux :} \zh{手机被没偷走。}
  \item \textcolor{popgreen}{Vrai :} \zh{手机没有被偷走。} \emph{Shǒujī méiyǒu bèi tōuzǒu.}
  \end{itemize}
\item \textbf{Verbe « nu » après 被 :} Interdit.
  \begin{itemize}
  \item \textcolor{red}{Faux :} \zh{西红柿被撒了。} (Verbe nu + \zh{了} insuffisant si pas de résultat clair)
  \item \textcolor{popgreen}{Vrai :} \zh{西红柿被撒了一地 / 撒在地上了。}
  \item \textcolor{red}{Faux :} \zh{作业被做了。} (Ambigu : fait par qui ? fini ?)
  \item \textcolor{popgreen}{Vrai :} \zh{作业被做完了。} / \zh{作业被做好了。}
  \end{itemize}
\item \textbf{Agent = Instrument / Moyen :} On n'utilise pas \zh{被} pour l'instrument.
  \begin{itemize}
  \item \textcolor{red}{Faux :} \zh{菜被刀切了。} (Le couteau n'est pas un agent volontaire)
  \item \textcolor{popgreen}{Vrai :} \zh{用刀切菜。} / \zh{菜被切了。} (Agent omis)
  \end{itemize}
\item \textbf{Confusion 被 / 把 :} \zh{把} = active (sujet AGIT sur objet) ; \zh{被} = passive (sujet SUBIT).
  \begin{itemize}
  \item \zh{我把手机弄坏了。} (J'ai cassé mon tel -- MA faute)
  \item \zh{手机被我弄坏了。} (Mon tel a été cassé par moi -- FOCUS sur le tel)
  \end{itemize}
\item \textbf{« Être + adjectif » n'est pas un passif :} \zh{他被很高兴。} \textcolor{red}{IMPOSSIBLE}. \zh{高兴} est adjectif/verbe d'état. \rightarrow \zh{他很高兴。}
\end{itemize}
\end{attention}

\courssub{Comparaison français / chinois : stratégies de traduction}

\begin{methode}[Traduire le passif français vers le chinois]
\begin{enumerate}
\item \textbf{Identifie le verbe français} : est-ce un événement subi / fâcheux ? $\rightarrow$ \zh{被} possible.
\item \textbf{Identifie l'agent} : Exprimé ? $\rightarrow$ le placer après \zh{被}. Inconnu / indifférent ? $\rightarrow$ omettre l'agent.
\item \textbf{Complète le verbe chinois} : Le participe passé français cache souvent un résultat. Ajoute \zh{了}, \zh{完}, \zh{坏}, \zh{走}, \zh{好}, \zh{在...}, \zh{给...} selon le sens.
\item \textbf{Vérifie la négation} : « n'a pas été... » $\rightarrow$ \zh{没有被...} (pas \zh{被没有}).
\item \textbf{Si c'est une bonne nouvelle / neutre} : Évite \zh{被}. Utilise la voix active (Sujet agent + Verbe + Objet patient) ou structure \zh{受} (formel : \zh{受邀}, \zh{受表彰}).
\end{enumerate}
\end{methode}

\begin{exemplebox}[Exemples de traduction]
\begin{itemize}
\item \fr{« Le livre a été lu par lui. »} (Neutre) $\rightarrow$ \zh{他读了这本书。} (Actif) \quad \emph{Tā dú le zhè běn shū.}
\item \fr{« Le livre a été déchiré (par lui). »} (Fâcheux) $\rightarrow$ \zh{书被(他)撕坏了。} \emph{Shū bèi (tā) sīhuài le.}
\item \fr{« Il a été invité. »} (Heureux) $\rightarrow$ \zh{他受邀了。} / \zh{有人邀请他。} \quad \emph{Tā shòuyāo le. / Yǒurén yāoqǐng tā.}
\item \fr{« Les devoirs n'ont pas été faits. »} $\rightarrow$ \zh{作业没有被做完。} \emph{Zuòyè méiyǒu bèi zuòwán.}
\item \fr{« La lettre sera envoyée demain. »} $\rightarrow$ \zh{信明天要被寄走。} \emph{Xìn míngtiān yào bèi jì zǒu.}
\end{itemize}
\end{exemplebox}

\courssub{Entraînement guidé et production}

\begin{exercice}[Exercices d'application — Travail individuel (15 min)]
\begin{enumerate}
\item \textbf{Transformation Passif $\rightarrow$ Actif.} Transforme en phrase active (sujet = agent). Donne le pinyin et la traduction.
  \begin{enumerate}
  \item \zh{窗户被风吹开了。}
  \item \zh{我的自行车被偷走了。}
  \item \zh{作业没有被学生做完。}
  \item \zh{那个坏人被警察抓住了。}
  \end{enumerate}

\item \textbf{Transformation Actif $\rightarrow$ Passif (avec 被).} Transforme en phrase passive. Attention au complément verbal !
  \begin{enumerate}
  \item \zh{小偷偷走了我的钱包。} (Ajoute \zh{了} ou \zh{走})
  \item \zh{妈妈做好了晚饭。} (Résultat \zh{好})
  \item \zh{没有人打扫这个房间。} (Négation + Agent omis $\rightarrow$ \zh{没有被})
  \item \zh{老师会批评那个学生。} (Modal \zh{会} + \zh{被})
  \end{enumerate}

\item \textbf{Correction d'erreurs.} Chaque phrase contient une erreur (ordre, négation, complément manquant, emploi abusif). Corrige et explique.
  \begin{enumerate}
  \item \zh{我被给了一本书。} (Bonne nouvelle)
  \item \zh{苹果被吃。} (Verbe nu)
  \item \zh{作业被没有做完。} (Place négation)
  \item \zh{树被斧头砍倒了。} (Instrument comme agent)
  \item \zh{他被高兴了。} (Adjectif traité comme passif)
  \end{enumerate}
\end{enumerate}
\end{exercice}

\begin{corrige}[Exercices d'application]
\begin{enumerate}
\item \textbf{Passif $\rightarrow$ Actif :}
  \begin{enumerate}
  \item \zh{风吹开了窗户。} \emph{Fēng chuīkāi le chuānghu.} (Le vent a ouvert la fenêtre.)
  \item \zh{(有人)偷走了我的自行车。} \emph{(Yǒurén) tōuzǒu le wǒ de zìxíngchē.} (Quelqu'un a volé mon vélo.)
  \item \zh{学生没有做完作业。} \emph{Xuéshēng méiyǒu zuòwán zuòyè.} (Les élèves n'ont pas fini les devoirs.)
  \item \zh{警察抓住了那个坏人。} \emph{Jǐngchá zhuāzhù le nà ge huàirén.} (La police a attrapé ce méchant.)
  \end{enumerate}

\item \textbf{Actif $\rightarrow$ Passif :}
  \begin{enumerate}
  \item \zh{我的钱包被小偷偷走了。} \emph{Wǒ de qiánbāo bèi xiǎotōu tōuzǒu le.}
  \item \zh{晚饭被妈妈做好了。} \emph{Wǎnfàn bèi māma zuòhǎo le.}
  \item \zh{这个房间没有被打扫。} \emph{Zhè ge fángjiān méiyǒu bèi dǎsǎo.} (Agent omis, négation correcte)
  \item \zh{那个学生会被老师批评。} \emph{Nà ge xuéshēng huì bèi lǎoshī pīpíng.}
  \end{enumerate}

\item \textbf{Corrections :}
  \begin{enumerate}
  \item \zh{我被给了一本书。} $\rightarrow$ \zh{别人给了我一本书。} / \zh{我收到了一本书。} (\zh{被} réservé aux événements fâcheux).
  \item \zh{苹果被吃。} $\rightarrow$ \zh{苹果被吃完了。} / \zh{苹果被吃掉了。} (Verbe nu $\rightarrow$ ajouter compl. résultat \zh{完} ou dir. \zh{掉}).
  \item \zh{作业被没有做完。} $\rightarrow$ \zh{作业没有被做完。} (\zh{没有} avant \zh{被}).
  \item \zh{树被斧头砍倒了。} $\rightarrow$ \zh{树被砍倒了。} (Agent omis) ou \zh{用斧头把树砍倒了。} (Structure \zh{把} avec instrument \zh{用}).
  \item \zh{他被高兴了。} $\rightarrow$ \zh{他很高兴。} / \zh{他高兴极了。} (Pas de passif sur état psychologique).
  \end{enumerate}
\end{enumerate}
\end{corrige}

\courssub{Production intégrée : Mini-récit « Incident au lycée »}

\begin{experience}[Tâche finale : Rédaction et présentation orale (25 min)]
\textbf{Consigne :} En groupes de 3. Choisis une situation ci-dessous. Rédige un récit de \textbf{5 à 6 phrases} en chinois.
\begin{itemize}
\item[(a)] Un étal de mangues renversé par une moto devant le lycée.
\item[(b)] Un portable cassé pendant la bousculade à la cantine.
\item[(c)] Le devoir de chinois oublié dans le taxi (modèle fourni).
\item[(d)] Un élève grondé par le proviseur devant la classe.
\item[(e)] Situation libre (incident réel ou imaginé au Cameroun).
\end{itemize}

\textbf{Contraintes obligatoires :}
\begin{enumerate}
\item Au moins \textbf{3 phrases en 被字句}.
\item Parmi elles : \textbf{1 affirmative}, \textbf{1 négative (\zh{没有被})}, \textbf{1 à agent omis}.
\item Chaque verbe passif \textbf{doit avoir un complément} (\zh{了}, \zh{走}, \zh{坏}, \zh{完}, \zh{在...}, \zh{给...}).
\item Le récit se termine par une chute (souvent positive) en \textbf{voix active} (éventuellement avec \zh{把}).
\end{enumerate}

\end{experience}

\begin{exoresolu}[Modèle commenté : Situation (b) — Portable cassé à la cantine]
\textbf{Production :}\\
中午在食堂，人太多了。我的手机被挤掉在地上了。屏幕被摔碎了。幸好，手机里的照片没有被删掉。手机可能被同学踩坏了。最后，同学帮我修好了屏幕，我很感激。\\
\emph{Zhōngwǔ zài shítáng, rén tài duō le. Wǒ de shǒujī bèi jǐ diào zài dìshang le. Píngmù bèi shuāisuì le. Xìnghǎo, shǒujī lǐ de zhàopiàn méiyǒu bèi shān diào. Shǒujī kěnéng bèi tóngxué cǎihuài le. Zuìhòu, tóngxué bāng wǒ xiūhǎo le píngmù, wǒ hěn gǎnjī.}

\textbf{Analyse des structures 被 :}
\begin{enumerate}
\item \zh{手机被挤掉在地上了} : Affirmative, agent omis (la foule), compl. direction \zh{掉在地上} + \zh{了}.
\item \zh{屏幕被摔碎了} : Affirmative, agent omis, compl. résultat \zh{碎} + \zh{了}.
\item \zh{照片没有被删掉} : \textbf{Négative}, \zh{没有} avant \zh{被}, compl. direction \zh{掉}.
\item \zh{手机可能被同学踩坏了} : Agent exprimé (\zh{同学}), modal \zh{可能}, compl. résultat \zh{坏} + \zh{了}.
\item \zh{同学帮我修好了屏幕} : Chute active (\zh{帮} + \zh{把} implicite), bonne nouvelle sans \zh{被}.
\end{enumerate}
\end{exoresolu}

\courssub{Réception orale et interaction}

\begin{experience}[Présentation et écoute active (10 min)]
\begin{enumerate}
\item Chaque groupe lit son récit (un élève lit, les autres jouent les bruitages si possible).
\item Les autres groupes notent sur une fiche les \textbf{phrases en 被} entendues (caractères ou pinyin).
\item \textbf{Défi transformation :} Le prof désigne une phrase notée $\rightarrow$ un élève la transforme à l'oral en phrase active.
\item Vote de la classe : « Récit le plus correct grammaticalement » + « Récit le plus vivant ».
\end{enumerate}
\end{experience}

\courssub{Lexique de la leçon (HSK 3--4)}

\begin{center}
\begin{tabularx}{\linewidth}{|X|X|X|X|}
\hline
\rowcolor{popblueL} Caractères & Pinyin & Nature & Traduction \\
\hline
被 & bèi & prep. & marque du passif (par) \\
\hline
摊子 & tānzi & n. & étal (de marché) \zh{[CL: 个]} \\
\hline
刮风 & guāfēng & v. & venter (il y a du vent) \\
\hline
吹倒 & chuīdǎo & v.+rés. & renverser par le vent \\
\hline
撒 & sā & v. & répandre, semer, renverser (liquide, objets menus) \\
\hline
碰掉 & pēngdiào & v.+dir. & heurter et faire tomber \\
\hline
摔碎 & shuāisuì & v.+rés. & casser en tombant (écran, verre) \\
\hline
偷走 & tōuzǒu & v.+dir. & voler (et emporter) \\
\hline
删掉 & shāndiào & v.+dir. & supprimer (fichier, message) \\
\hline
踩坏 & cǎihuài & v.+rés. & casser en marchant dessus / piétiner \\
\hline
批评 & pīpíng & v. & critiquer, réprimander \\
\hline
主任 & zhǔrèn & n. & directeur, responsable, chef de service \\
\hline
丢 & diū & v. & perdre, égarer \\
\hline
热闹 & rènao & adj. & animé, bruyant, vivant \\
\hline
着急 & zháojí & adj./v. & inquiet, anxieux / s'inquiéter \\
\hline
送回 & sònghuí & v.+dir. & ramener, rapporter (qch à qn) \\
\hline
遗忘 & yíwàng & v. & oublier (formel, écrit, utilisable en passif) \\
\hline
录取 & lùqǔ & v. & admettre (examen, concours) — \zh{被录取} \\
\hline
授予 & shòuyǔ & v. & décerner (titre, médaille) — \zh{被授予} \\
\hline
\end{tabularx}
\end{center}

\courssub{Point culture : Le passif dans la presse et l'administration}

\begin{savaistu}[Le passif chinois dans les médias]
En chinois moderne, le \zh{被字句} a connu une extension d'usage notable sur Internet et dans la presse depuis les années 2000 : le \textbf{« 被 + verbe monosyllabique »} sans complément apparent, souvent pour dénoncer une action subie de force par l'autorité.
\begin{itemize}
\item \zh{被自杀} (bèi zìshā) : « être suicidé » = présenté comme suicide par la police, mais meurtre suspecté.
\item \zh{被失踪} (bèi shīzōng) : « être disparu » = disparition forcée.
\item \zh{被就业} (bèi jiùyè) : « être employé » = compté comme employé dans les stats alors qu'on est au chômage.
\item \zh{被代表} (bèi dàibiǎo) : « être représenté » = un délégué nommé d'office sans élection.
\end{itemize}
Ces néologismes satiriques détournent la règle du « complément obligatoire » pour créer un effet d'ironie : le verbe nu devient un \textbf{nom d'action subie}. C'est un marqueur de \emph{langue de l'espace public chinois contemporain}, à connaître pour la compréhension de l'écrit (niveau HSK 5+ / presse), mais \textbf{interdit à l'examen officiel} (BAC, HSK 4) où la norme standard (verbe + complément) reste exigée.
\end{savaistu}

\courssub{Bilan et devoirs}

\begin{aretenir}[Check-list de révision — 被字句]
\noindent
\begin{tabularx}{\linewidth}{|X|X|}
\hline
\rowcolor{popblueL} \textbf{Point clé} & \textbf{Teste-toi : sais-tu... ?} \\
\hline
Schéma & Patient + \zh{被} + (Agent) + V + \textbf{Complément} ? \\
\hline
Négation & \zh{没有/没} + \zh{被} + Agent + V + Compl. ? (Jamais \zh{被没有}) \\
\hline
Agent omis & \zh{被} + V + Compl. directement ? \\
\hline
Modal & \zh{要/会/能/必须} + \zh{被} + ... ? \\
\hline
Sémantique & Réserver \zh{被} aux événements \textbf{subis / fâcheux} ? \\
\hline
Instrument & Ne pas mettre l'outil après \zh{被} (utiliser \zh{用}) ? \\
\hline
Traduction FR $\rightarrow$ CN & Convertir « bonne nouvelle passive » en actif / \zh{受} ? \\
\hline
\end{tabularx}
\end{aretenir}

\textbf{Devoirs (à rendre au prochain cours) :}
\begin{enumerate}
\item Exercices du cahier d'activités : p. 42--43 (Transformation, Complétion, Traduction 5 phrases FR $\rightarrow$ CN).
\item \textbf{Production écrite personnelle :} Rédige un court paragraphe (6--8 phrases) racontant \textbf{un vrai petit incident} qui t'est arrivé récemment (au marché, au lycée, dans le transport, à la maison). Utilise \textbf{au moins 4 phrases en 被} (affirmative, négative, agent omis, avec modal). Termine par la résolution (voix active). Écris en caractères, pinyin et français.
\item \textbf{Préparation orale :} Entraîne-toi à lire ton paragraphe à voix haute (enregistrement possible sur OnBuch+). Attention aux tons de \zh{被} (4e ton) et des compléments de résultat (\zh{倒} 3, \zh{坏} 4, \zh{走} 3, \zh{完} 2).
\end{enumerate}

\newpage

\coursec{Méthodes et exercices résolus}

\coursec{Leçon 18 — Grammaire : 被字句 (La construction passive avec 被)}

\begin{textebox}[Situation de communication : Un incident au marché de Mokolo]
\zh{— 怎么了？你的手机呢？} \\
\zh{— 我的手机被小偷偷走了。} \\
\zh{— 啊？那你报警了吗？} \\
\zh{— 报了，警察说会帮我找。可是钱包也被偷了，里面的钱被花完了。} \\
\zh{— 真倒霉！下次小心点儿。}
\par
\textbf{Pinyin :} \\
\zh{— Zěnme le? Nǐ de shǒujī ne?} \\
\zh{— Wǒ de shǒujī bèi xiǎotōu tōu zǒu le.} \\
\zh{— A? Nà nǐ bàojǐng le ma?} \\
\zh{— Bào le, jǐngchá shuō huì bāng wǒ zhǎo. Kěshì qiánbāo yě bèi tōu le, lǐmiàn de qián bèi huā wán le.} \\
\zh{— Zhēn dǎoméi! Xiàcì xiǎoxīn diǎnr.}
\par
\textbf{Traduction :} \\
« — Qu'est-ce qui se passe ? Ton téléphone ? — Mon téléphone a été volé par un pickpocket. — Ah ? Tu as appelé la police ? — Oui, la police a dit qu'elle m'aiderait à le retrouver. Mais le portefeuille a aussi été volé, et l'argent dedans a été tout dépensé. — Vraiment pas de chance ! Fais attention la prochaine fois. »
\end{textebox}

\courssub{Observation et analyse}

\begin{propriete}[Schéma fondamental de la 被字句]
La structure \cle{被字句} (bèi zì jù) permet d'exprimer la voix passive : le \textbf{sujet grammatical} (le \cle{patient}) \textbf{subit} l'action effectuée par le \cle{complément d'agent} (introduit par 被).

\begin{center}
\begin{tabularx}{\linewidth}{|X|X|X|X|X|}
\hline
\rowcolor{popblueL} \textbf{Patient (Sujet)} & \textbf{被} & \textbf{Agent} & \textbf{Verbe (+ Complément)} & \textbf{Particule / Complément final} \\
\hline
\zh{我的自行车} (Wǒ de zìxíngchē) & \zh{被} & \zh{小偷} (xiǎotōu) & \zh{偷走} (tōu zǒu) & \zh{了} (le) \\
\hline
\zh{窗户的玻璃} (Chuānghu de bōli) & \zh{被} & \zh{孩子们} (háizimen) & \zh{打破} (dǎpò) & \zh{了} (le) \\
\hline
\zh{作业} (Zuòyè) & \zh{被} & \zh{老师} (lǎoshī) & \zh{批改} (pīgǎi) & \zh{完} (wán) \\
\hline
\end{tabularx}
\end{center}

\textbf{Formule canonique :} \textbf{Patient + 被 + Agent + Verbe + Complément (résultat / direction / 了 / quantificateur...)}
\end{propriete}

\begin{definition}[Notions clés : Patient, Agent, Verbe transitif]
\begin{itemize}
\item \textbf{Patient (受事 zhòushì) :} L'entité qui \emph{subit} l'action. Il devient le sujet de la phrase passive.
\item \textbf{Agent (施事 shīshì) :} L'entité qui \emph{accomplit} l'action. Il est introduit par 被 (ou 叫/让).
\item \textbf{Verbe transitif (及物动词 jíwù dòngcí) :} Verbe qui appelle un objet direct (打, 吃, 偷, 批评, 破坏, 写, 建…). \textbf{Seuls les verbes transitifs d'action peuvent entrer dans une 被字句.}
\end{itemize}
\end{definition}

\begin{attention}[Condition absolue : Le verbe ne doit pas rester « nu »]
En chinois, le verbe d'une 被字句 \textbf{ne peut presque jamais apparaître seul}. Il doit être suivi d'au moins un des éléments suivants :
\begin{enumerate}
\item La particule d'accompli \zh{了} (le) : \zh{信被寄了。} (Xìn bèi jì le.)
\item Un \textbf{complément de résultat} (结果补语) : \zh{打破} (dǎpò), \zh{吃完} (chīwán), \zh{偷走} (tōuzǒu), \zh{修好} (xiūhǎo).
\item Un \textbf{complément de direction} (趋向补语) : \zh{拿走} (ná zǒu), \zh{送来} (sòng lái).
\item Un \textbf{complément d'objet} (si le verbe en prend un second) : \zh{他被选为班长。} (Tā bèi xuǎn wéi bānzhǎng.)
\item Un \textbf{quantificateur / durative} : \zh{被批评了一顿。} (bèi pīpíng le yí dùn).
\end{enumerate}
\emph{Exemple fautif :} \zh{*书被他买。} \rightarrow \textbf{Correct :} \zh{书被他买走了。} (Shū bèi tā mǎi zǒu le.)
\end{attention}

\courssub{Emplois, variantes et exceptions}

\begin{propriete}[Les trois marques du passif : 被, 叫, 让]
\begin{center}
\begin{tabularx}{\linewidth}{|X|X|X|}
\hline
\rowcolor{popblueL} \textbf{Marque} & \textbf{Registre / Nuance} & \textbf{Contrainte sur l'agent} \\
\hline
\zh{被} (bèi) & Standard, neutre, écrit \& oral. \textbf{Forme de référence pour l'examen.} & Agent \textbf{optionnel} (peut être omis si inconnu/général). \\
\hline
\zh{叫} (jiào) & Familier, oral (Nord de la Chine). & Agent \textbf{obligatoire}. \\
\hline
\zh{让} (ràng) & Familier, oral (Sud de la Chine / standard). & Agent \textbf{obligatoire}. \\
\hline
\end{tabularx}
\end{center}

\textbf{Exemples comparés :}
\begin{itemize}
\item Standard : \zh{我的自行车被偷了。} (Wǒ de zìxíngchē bèi tōu le.) — Agent omis.
\item Familier (叫) : \zh{我的自行车叫人偷走了。} (Wǒ de zìxíngchē jiào rén tōu zǒu le.) — Agent \zh{人} (rén) obligatoire.
\item Familier (让) : \zh{我的自行车让人偷走了。} (Wǒ de zìxíngchē ràng rén tōu zǒu le.)
\end{itemize}

\textbf{Variante orale renforcée :} \zh{被……给……} (bèi... gěi...). Ex : \zh{钱包被人给偷了。} (Qiánbāo bèi rén gěi tōu le.) — Insiste sur le subit soudain/inattendu.
\end{propriete}

\begin{propriete}[Cas particuliers : Agent omis / Agent général / Sujet inanimé]
\begin{enumerate}
\item \textbf{Agent inconnu ou indifférent :} On omet l'agent ou on utilise \zh{人} (rén, « quelqu'un ») / \zh{大家} (dàjiā).
\zh{门被(人)打开了。} (Mén bèi (rén) dǎkāi le.) « La porte a été ouverte. »
\item \textbf{Sujet inanimé + verbe d'action :} Très fréquent. \zh{这棵树被风吹倒了。} (Zhè kē shù bèi fēng chuīdǎo le.) « Cet arbre a été renversé par le vent. »
\item \textbf{Verbes à double objet (donner, envoyer, apprendre...) :} Seul l'objet \textbf{direct} (la chose) peut devenir sujet de 被. L'objet indirect (le destinataire) reste après le verbe.
\zh{这本书被老师送给我了。} (Zhè běn shū bèi lǎoshī sòng gěi wǒ le.) « Ce livre a été donné par le prof à moi. » (Pas : *我被老师送这本书了。)
\end{enumerate}
\end{propriete}

\begin{aretenir}[Verbes INTERDITS dans une 被字句]
Ne JAMAIS utiliser 被 avec :
\begin{itemize}
\item \textbf{Verbes d'état / adjectifs verbaux :} \zh{是, 有, 喜欢, 爱, 讨厌, 知道, 认识, 想, 会, 能...}
\emph{Faux :} \zh{*我被喜欢。} / \zh{*他被是学生。}
\item \textbf{Verbes intransitifs (sans objet) :} \zh{去, 来, 走, 睡觉, 死, 笑, 哭...}
\emph{Faux :} \zh{*他被死了。} (On dit : \zh{他死了。} ou \zh{他被害死了。} avec un verbe causatif).
\item \textbf{Verbes psychologiques / sensoriels purs :} \zh{见, 闻, 听见, 看见} (souvent remplacés par \zh{被看到/被听到}).
\end{itemize}
\end{aretenir}

\courssub{Négation, adverbes et ordre des mots}

\begin{propriete}[Position de la négation et des adverbes : TOUT avant 被]
Dans une 被字句, \textbf{tout élément qui n'est ni le Patient, ni l'Agent, ni le Verbe+Complément se place AVANT 被}.

\begin{center}
\begin{tabularx}{\linewidth}{|X|X|}
\hline
\rowcolor{popblueL} \textbf{Élément} & \textbf{Position correcte} \\
\hline
Négation \zh{没(有)} (méi(yǒu)) & \zh{作业 \textbf{没有} 被老师收走。} (Zuòyè \textbf{méiyǒu} bèi lǎoshī shōu zǒu.) \\
\hline
Adverbe de temps & \zh{他 \textbf{昨天} 被批评了。} (Tā \textbf{zuótiān} bèi pīpíng le.) \\
\hline
Adverbe de fréquence & \zh{他 \textbf{常常} 被表扬。} (Tā \textbf{chángcháng} bèi biǎoyáng.) \\
\hline
Adverbe de degré / modal & \zh{这事 \textbf{一定} 被他知道了。} (Zhè shì \textbf{yídìng} bèi tā zhīdào le.) \\
\hline
Adverbe de portée \zh{都/全} (dōu/quán) & \zh{剩饭 \textbf{都} 被吃完了。} (Shèngfàn \textbf{dōu} bèi chī wán le.) \\
\hline
\end{tabularx}
\end{center}

\emph{Règle mnémotechnique :} \textbf{Patient + [Nég/Adv/Temps] + 被 + Agent + V + Compl.}
\end{propriete}

\begin{exemplebox}[Exemples complets avec négation et adverbes]
\begin{itemize}
\item \zh{这封信没有被寄出去。} (Zhè fēng xìn méiyǒu bèi jì chūqù.) « Cette lettre n'a pas été envoyée. »
\item \zh{足球赛因为下雨被取消了。} (Zúqiú sài yīnwèi xiàyǔ bèi qǔxiāo le.) « Le match de foot a été annulé à cause de la pluie. »
\item \zh{所有的苹果都被小明吃光了。} (Suǒyǒu de píngguǒ dōu bèi Xiǎomíng chī guāng le.) « Toutes les pommes ont été mangées entièrement par Xiaoming. »
\end{itemize}
\end{exemplebox}

\courssub{Transformation Active (把) $\leftrightarrow$ Passive (被)}

\begin{methode}[Transformer une phrase en 把 en phrase en 被 (et inversement)]
Les structures \zh{把} (bǎ) et \zh{被} (bèi) sont \textbf{miroirs} l'une de l'autre. Même verbe, mêmes compléments, rôles inversés.

\begin{enumerate}
\item \textbf{Identifier la phrase en 把 :} Sujet-Agent + \zh{把} + Objet-Patient + Verbe + Complément.
\item \textbf{Échanger Sujet et Objet :} L'Objet-Patient devient Sujet ; le Sujet-Agent devient Complément d'agent après 被.
\item \textbf{Conserver intacts :} Le verbe, ses compléments (résultat, direction, 了, quantificateurs), les adverbes placés avant 把 (ils passent avant 被).
\item \textbf{Remplacer 把 par 被.}
\end{enumerate}

\begin{center}
\begin{tikzpicture}[node distance=1.2cm, every node/.style={font=\small}]
\node[draw, rounded corners, fill=popblueL] (active) {Actif (把) : \zh{大雨 把 路 冲坏 了}};
\node[draw, rounded corners, fill=poporangeL, below=of active] (passive) {Passif (被) : \zh{路 被 大雨 冲坏 了}};
\draw[<->, thick, popdark] (active.south) -- (passive.north) node[midway, right, fill=white, inner sep=2pt] {\textbf{Inversion Agent $\leftrightarrow$ Patient}};
\end{tikzpicture}
\end{center}
\end{methode}

\begin{exoresolu}[Transformation systématique 把 $\rightarrow$ 被]
\textbf{Énoncé.} Transforme en 被字句 et traduis.

a) \zh{小偷把我的钱包偷走了。} (Xiǎotōu bǎ wǒ de qiánbāo tōu zǒu le.)
b) \zh{妈妈把剩菜都热好了。} (Māma bǎ shèngcài dōu rè hǎo le.)
c) \zh{台风把房顶吹飞了。} (Táifēng bǎ fángdǐng chuī fēi le.)

\tcblower
\textbf{Solution détaillée.}

a) Patient : \zh{我的钱包} ; Agent : \zh{小偷} ; V+Compl : \zh{偷走了}.
\zh{我的钱包被小偷偷走了。} (Wǒ de qiánbāo bèi xiǎotōu tōu zǒu le.) « Mon portefeuille a été volé par un pickpocket. »

b) Patient : \zh{剩菜} ; Adverbe \zh{都} (portée sur le patient) $\rightarrow$ avant 被 ; Agent : \zh{妈妈} ; V+Compl : \zh{热好了}.
\zh{剩菜都被妈妈热好了。} (Shèngcài dōu bèi māma rè hǎo le.) « Les restes ont tous été réchauffés par maman. »

c) Patient : \zh{房顶} ; Agent : \zh{台风} ; V+Compl : \zh{吹飞了}.
\zh{房顶被台风吹飞了。} (Fángdǐng bèi táifēng chuī fēi le.) « Le toit a été soufflé (emporté) par le typhon. »
\end{exoresolu}

\courssub{Traduction : Le piège du calque français}

\begin{attention}[Le passif français $\neq$ 被字句 automatique]
Le français utilise le passif (\emph{être + participe passé}) pour : focaliser sur le patient, masquer l'agent, style administratif, objectivité scientifique. Le chinois \textbf{préfère souvent l'actif} ou des structures sans marque passive.

\begin{center}
\begin{tabularx}{\linewidth}{|X|X|X|}
\hline
\rowcolor{popblueL} \textbf{Type de passif français} & \textbf{Stratégie chinoise privilégiée} & \textbf{Exemple} \\
\hline
Passif « notional » / neutre (sujet inanimé, pas d'agent, action achevée) & \textbf{Actif à sujet inanimé} ou \textbf{Verbe + 好/完/了} sans 被 & \zh{信写好了。} (Xìn xiě hǎo le.) « La lettre est écrite. » \\
\hline
Passif défavorable / subi (vol, critique, accident, dommage) & \textbf{被字句} (standard) & \zh{手机被偷了。} (Shǒujī bèi tōu le.) « Le téléphone a été volé. » \\
\hline
Passif favorable / honorifique (nomination, récompense, élection) & \textbf{被字句} (usage moderne) OU \textbf{是……的} & \zh{他被评为优秀学生。} / \zh{他是被评为优秀学生的。} \\
\hline
Passif avec agent exprimé (style admin./technique) & \textbf{被……所} (formel écrit) OU \textbf{Active} & \zh{这座桥是中国工程师建的。} (Shì...de) \\
\hline
« Se faire + infinitif » (malheur) & \textbf{被} & \zh{他被骗了钱。} (Tā bèi piàn le qián.) « Il s'est fait voler de l'argent. » \\
\hline
\end{tabularx}
\end{center}
\end{attention}

\begin{exoresolu}[Choix de la bonne structure : Thème ciblé Bac]
\textbf{Énoncé.} Traduis en chinois en justifiant le choix de la structure.

1. « Le déjeuner a été préparé par ma mère. » (Contexte : description routine, neutre)
2. « Mon vélo a été volé hier soir. » (Contexte : plainte, malheur)
3. « Ce roman a été écrit par Mo Yan. » (Contexte : présentation littéraire)
4. « Les devoirs n'ont pas été faits. » (Contexte : constat neutre)

\tcblower
\textbf{Solution détaillée.}

1. Neutre, agent connu, routine $\rightarrow$ \textbf{Actif naturel} : \zh{午饭是妈妈做的。} (Wǔfàn shì māma zuò de.) — Structure \zh{是……的} pour focaliser sur l'agent. Ou simplement : \zh{妈妈做了午饭。}
2. Malheur, agent inconnu $\rightarrow$ \textbf{被字句 obligatoire} : \zh{我的自行车昨晚被偷了。} (Wǒ de zìxíngchē zuówǎn bèi tōu le.)
3. Présentation littéraire, agent célèbre $\rightarrow$ \textbf{是……的} (élégant) : \zh{这部小说是莫言写的。} (Zhè bù xiǎoshuō shì Mòyán xiě de.) — \zh{被} possible mais lourd : \zh{这部小说被莫言写了。}
4. Constat neutre, pas d'agent $\rightarrow$ \textbf{Sans 被} : \zh{作业没做完。} (Zuòyè méi zuò wán.) — \zh{作业没有被做完} = calque, peu naturel.
\end{exoresolu}

\courssub{Vocabulaire clé de la leçon}

\begin{center}
\begin{tabularx}{\linewidth}{|X|X|X|X|}
\hline
\rowcolor{popblueL} \textbf{Caractères} & \textbf{Pinyin} & \textbf{Nature} & \textbf{Traduction} \\
\hline
被 & bèi & prép. & marque du passif (par) \\
叫 & jiào & prép. & marque du passif (familier, par) \\
让 & ràng & prép. & marque du passif (familier, par) \\
给 & gěi & part. & renforce le passif (oral : 被…给…) \\
偷 & tōu & v. & voler \\
偷走 & tōu zǒu & v.+compl. & voler (et emporter) \\
打破 & dǎ pò & v.+compl. & casser (en frappant) \\
冲坏 & chōng huài & v.+compl. & endommager par l'eau (inonder) \\
批评 & pī píng & v./n. & critiquer / critique \\
表扬 & biǎo yáng & v./n. & louer, féliciter / éloge \\
评为 & píng wéi & v. & élire, désigner comme \\
建 & jiàn & v. & construire, bâtir \\
修好 & xiū hǎo & v.+compl. & réparer (avec succès) \\
倒霉 & dǎo méi & adj. & malchanceux, poisseux \\
小偷 & xiǎo tōu & n. & pickpocket, voleur \\
警察 & jǐng chá & n. & police, policier \\
台风 & tái fēng & n. & typhon \\
剩饭 / 剩菜 & shèng fàn / shèng cài & n. & restes (riz / plats) \\
零花钱 & líng huā qián & n. & argent de poche \\
\hline
\end{tabularx}
\end{center}

\courssub{Exercices d'application}

\begin{exercice}[Exercice 1 : Compléter par 被, 叫 ou 让 (selon le registre indiqué)]
Complète les phrases. Indique si l'agent est omis ou non.
\begin{enumerate}
\item (Standard) 这封信 \_\_\_\_ 寄出去了。 (Agent omis)
\item (Familier) 我的笔 \_\_\_\_ 弟弟 拿走 了。
\item (Standard) 窗户 \_\_\_\_ 风 吹开 了。
\item (Familier, Sud) 作业 \_\_\_\_ 狗 吃 了。
\item (Renforcé oral) 钱 \_\_\_\_ 人 \_\_\_\_ 拿 走了。
\end{enumerate}
\end{exercice}

\begin{exercice}[Exercice 2 : Transformer l'actif (把) en passif (被)]
Transforme et traduis.
\begin{enumerate}
\item \zh{哥哥把苹果吃完了。} (Gēge bǎ píngguǒ chī wán le.)
\item \zh{大雨把衣服淋湿了。} (Dàyǔ bǎ yīfu lín shī le.)
\item \zh{老师把作业都批改完了。} (Lǎoshī bǎ zuòyè dōu pīgǎi wán le.)
\item \zh{谁把门关上了？} (Shéi bǎ mén guān shàng le?) $\rightarrow$ Passif : \zh{门被谁关上了？}
\end{enumerate}
\end{exercice}

\begin{exercice}[Exercice 3 : Placer la négation / l'adverbe au bon endroit]
Réécris la phrase en ajoutant l'élément entre parenthèses.
\begin{enumerate}
\item \zh{自行车被偷走了。} (没) $\rightarrow$ \zh{自行车 \textbf{没} 被偷走。}
\item \zh{他被批评了。} (昨天) $\rightarrow$ \zh{\textbf{昨天} 他被批评了。}
\item \zh{饭被吃完了。} (都) $\rightarrow$ \zh{饭 \textbf{都} 被吃完了。}
\item \zh{信被寄了。} (已经) $\rightarrow$ \zh{信 \textbf{已经} 被寄了。}
\end{enumerate}
\end{exercice}

\begin{exercice}[Exercice 4 : Traduction thème (Français $\rightarrow$ Chinois) — Pièges]
Traduis en choisissant la structure adéquate (被 / Actif / 是…的 / Sans marque).
\begin{enumerate}
\item « La route a été coupée par la inondation. » (Contexte : info météo Douala)
\item « Ce gâteau a été fait par ma sœur. » (Contexte : fierté, agent important)
\item « Il s'est fait mordre par un chien. » (Malheur)
\item « La réunion a été annulée. » (Neutre, pas d'agent)
\item « Le livre a été lu par tout le monde. » (Portée \zh{都})
\end{enumerate}
\end{exercice}

\begin{exercice}[Exercice 5 : Correction d'erreurs (Type Bac)]
Corrige les phrases fautives suivantes (grammaire, ordre, choix du verbe).
\begin{enumerate}
\item \zh{* 我被高兴了。}
\item \zh{* 书被他买。}
\item \zh{* 作业被没有做完。}
\item \zh{* 我被狗咬。}
\item \zh{* 信被写好了由他。} (Agent mal placé)
\end{enumerate}
\end{exercice}

\courssub{Corrigés détaillés}

\begin{corrige}[Exercice 1]
\begin{enumerate}
\item \zh{被} (Agent omis). \zh{这封信被寄出去了。}
\item \zh{叫} (ou \zh{让}). Agent \zh{弟弟} obligatoire. \zh{我的笔叫弟弟拿走了。}
\item \zh{被}. \zh{窗户被风吹开了。}
\item \zh{让}. Agent \zh{狗} obligatoire. \zh{作业让狗吃了。}
\item \zh{被 / 给}. \zh{钱被人给拿走了。}
\end{enumerate}
\end{corrige}

\begin{corrige}[Exercice 2]
\begin{enumerate}
\item \zh{苹果被哥哥吃完了。} (Píngguǒ bèi gēge chī wán le.) « La pomme a été mangée entièrement par le grand frère. »
\item \zh{衣服被大雨淋湿了。} (Yīfu bèi dàyǔ lín shī le.) « Les vêtements ont été trempés par la grosse pluie. »
\item \zh{作业都被老师批改完了。} (Zuòyè dōu bèi lǎoshī pīgǎi wán le.) « Les devoirs ont tous été corrigés par le prof. » (\zh{都} avant 被 car porte sur le patient \zh{作业}).
\item \zh{门被谁关上了？} (Mén bèi shéi guān shàng le?) « La porte a été fermée par qui ? »
\end{enumerate}
\end{corrige}

\begin{corrige}[Exercice 3]
\begin{enumerate}
\item \zh{自行车没被偷走。} (Zìxíngchē méi bèi tōu zǒu.) « Le vélo n'a pas été volé. »
\item \zh{他昨天被批评了。} (Tā zuótiān bèi pīpíng le.) « Il a été critiqué hier. »
\item \zh{饭都被吃完了。} (Fàn dōu bèi chī wán le.) « Le riz a tout été mangé. »
\item \zh{信已经被寄了。} (Xìn yǐjīng bèi jì le.) « La lettre a déjà été envoyée. »
\end{enumerate}
\end{corrige}

\begin{corrige}[Exercice 4]
\begin{enumerate}
\item \zh{路被洪水冲断了。} (Lù bèi hóngshuǐ chōng duàn le.) — Passif défavorable, agent naturel. \zh{冲断} (emporter/couper par l'eau).
\item \zh{这个蛋糕是我姐姐做的。} (Zhè ge dàngāo shì wǒ jiějie zuò de.) — \textbf{是……的} : focalisation sur l'agent (fierté). \zh{被} possible mais moins naturel pour une action positive volontaire.
\item \zh{他被狗咬了。} (Tā bèi gǒu yǎo le.) — Malheur, \zh{被} obligatoire. Verbe \zh{咬} (mordre) + \zh{了}.
\item \zh{会议取消了。} (Huìyì qǔxiāo le.) — Neutre, pas d'agent $\rightarrow$ Actif sujet inanimé / Verbe d'état. \zh{*会议被取消了} = calque admin.
\item \zh{这本书被大家都读过了。} (Zhè běn shū bèi dàjiā dōu dú guò le.) — Agent \zh{大家}, portée \zh{都} avant 被. \zh{读过} (verbe + aspect expérientiel).
\end{enumerate}
\end{corrige}

\begin{corrige}[Exercice 5]
\begin{enumerate}
\item \zh{* 我被高兴了。} $\rightarrow$ \textbf{Impossible.} \zh{高兴} est adjectif/verbe d'état. \rightarrow \zh{我很高兴。} (Wǒ hěn gāoxìng.) « Je suis content. »
\item \zh{* 书被他买。} $\rightarrow$ \textbf{Verbe nu.} Il faut un complément. \rightarrow \zh{书被他买走了。} (Shū bèi tā mǎi zǒu le.) / \zh{书被他买好了。}
\item \zh{* 作业被没有做完。} $\rightarrow$ \textbf{Négation mal placée.} \zh{没(有)} avant 被. \rightarrow \zh{作业没有被做完。} (Zuòyè méiyǒu bèi zuò wán.) Mieux (neutre) : \zh{作业没做完。}
\item \zh{* 我被狗咬。} $\rightarrow$ \textbf{Verbe nu.} \rightarrow \zh{我被狗咬了。} (Wǒ bèi gǒu yǎo le.) / \zh{我被狗咬了一口。} (Wǒ bèi gǒu yǎo le yì kǒu.)
\item \zh{* 信被写好了由他。} $\rightarrow$ \textbf{Ordre Agent/Verbe inversé.} Agent après 被, avant verbe. \rightarrow \zh{信被他写好了。} (Xìn bèi tā xiě hǎo le.) Ou \zh{信是他写好的。}
\end{enumerate}
\end{corrige}

\courssub{Production écrite guidée}

\begin{experience}[Atelier rédaction : \zh{倒霉的一天} (Une journée de malchance)]
\textbf{Consigne :} Rédige un court récit (6-8 phrases) intitulé \zh{倒霉的一天}. Contraintes :
\begin{itemize}
\item Utiliser \textbf{au moins 3 phrases en 被} (dont une avec négation, une avec adverbe de temps, une avec 叫/让).
\item Intégrer les mots : \zh{自行车} (vélo), \zh{雨} (pluie), \zh{老师} (prof), \zh{作业} (devoirs), \zh{手机} (téléphone).
\item Alterner narration (actif) et passif (被) pour le rythme.
\item Utiliser les connecteurs : \zh{先……然后……最后} (d'abord… ensuite… enfin), \zh{因为……所以……}, \zh{结果}.
\end{itemize}

\textbf{Amorce :} \zh{昨天早上，我骑自行车去学校……}
\end{experience}

\begin{exoresolu}[Modèle de production attendue (Niveau Terminale)]
\zh{昨天真是倒霉的一天。早上，我骑自行车去学校，半路\textbf{被大雨淋湿了}，结果迟到了。\textbf{被老师批评了一顿}。更糟的是，放学时发现\textbf{手机被同桌叫拿走了}（原来是他开玩笑），\textbf{作业也没被带回来}，在教室里。最后回家，\textbf{自行车让人把锁偷了}，只好走路回家。真希望明天别再倒霉了！}

\textbf{Pinyin :} Zuótiān zhēnshi dǎoméi de yì tiān. Zǎoshang, wǒ qí zìxíngchē qù xuéxiào, bànlù \textbf{bèi dàyǔ lín shī le}, jiéguǒ chídào le. \textbf{Bèi lǎoshī pīpíng le yí dùn}. Gèng zāo de shì, fàngxué shí fāxiàn \textbf{shǒujī bèi tóngzuò jiào ná zǒu le} (yuánlái shì tā kāiwánxiào), \textbf{zuòyè yě méi bèi dài huílái}, zài jiàoshì lǐ. Zuìhòu huí jiā, \zh{自行车让人把锁偷了}, zhǐhǎo zǒulù huí jiā. Zhēn xīwàng míngtiān bié zài dǎoméi le!

\textbf{Analyse grammaticale :} 4 phrases en 被/叫/让. 1) \zh{被大雨淋湿了} (V+Résultat+了). 2) \zh{被老师批评了一顿} (V+Quantificateur). 3) \zh{叫拿走了} (Registre familier, Agent obligatoire). 4) \zh{没被带回来} (Négation avant 被). 5) \zh{让人把锁偷了} (Structure \zh{让} + Agent + \zh{把} + Objet + V — passif causatif complexe, niveau excellence).
\end{exoresolu}

\begin{savaistu}[Sociolinguistique : Le passif en chinois, en français et au Cameroun]
\begin{itemize}
\item \textbf{Fréquence :} Le passif \zh{被} est \textbf{beaucoup moins fréquent} en chinois qu'en français. Le chinois préfère l'actif (\zh{有人修好了路} vs \zh{La route a été réparée}) ou la topicalisation (\zh{路，修好了}).
\item \textbf{Valeur affective :} Historiquement, \zh{被} marque un \textbf{événement subi, souvent négatif} (malheur, dommage). L'extension aux événements positifs (\zh{被评为}, \zh{被录取}) est un phénomène moderne (XXe s.), influencé par la traduction de l'anglais (\emph{be elected}) et du français.
\item \textbf{Contexte camerounais :} Dans l'administration (français), le passif est roi (\emph{« Il est décidé que... », « Les dossiers sont traités »}). En chinois standard, on dirait : \zh{决定……} (Jué dìng...), \zh{档案处理完了。} (Dàng'àn chǔlǐ wán le.) Apprendre à \textbf{désactiver le réflexe passif français} est la clé pour un chinois naturel.
\item \textbf{注 :} Au Cameroun francophone, l'usage du \emph{« on »} impersonnel (\emph{« On a volé mon téléphone »}) correspond souvent au chinois \zh{被人偷了} (agent \zh{人} générique) ou à l'actif \zh{手机被偷了} (agent omis).
\end{itemize}
\end{savaistu}

\begin{aretenir}[Check-list finale pour l'examen (Bac) — 被字句]
\begin{enumerate}
\item \textbf{Structure :} Patient + [Adv/Nég] + 被 + Agent + Verbe + \textbf{Complément obligatoire} (了 / Résultat / Direction / Quantificateur).
\item \textbf{Verbe :} Transitif d'action uniquement. Pas d'état, pas d'intransitif.
\item \textbf{Négation :} \zh{没(有)} \textbf{avant} 被. Jamais \zh{被没}.
\item \textbf{Adverbes (temps, fréquence, degré, 都) :} \textbf{Avant} 被.
\item \textbf{Agent :} Omis si inconnu/général avec \zh{被}. \textbf{Obligatoire} avec \zh{叫/让}.
\item \textbf{Traduction :} Ne calque pas le passif français. Demande-toi : « Malheur/Subi ? $\rightarrow$ 被. Neutre/Positif/Agent focus ? $\rightarrow$ Actif / \zh{是……的} / Sans marque. »
\item \textbf{Caractères :} \zh{被} (radical 衤 衣字旁, 4ᵉ ton) $\neq$ \zh{披} (radical 扌 提手旁, 1ʳᵉ ton).
\end{enumerate}
\end{aretenir}

\newpage

\coursec{Exercices}

\courssub{Je vérifie mes connaissances}

\begin{exercice}[QCM : la bonne structure]
Pour chaque question, entoure la bonne réponse.
\begin{enumerate}
\item La structure correcte de la phrase passive en chinois est :
\begin{enumerate}
\item Sujet (objet) + 被 + agent + verbe
\item Agent + 被 + sujet + verbe
\item Sujet + verbe + 被 + agent
\end{enumerate}
\item \zh{我的自行车被偷了。} signifie :
\begin{enumerate}
\item J'ai volé mon vélo.
\item Mon vélo a été volé.
\item Mon vélo vole.
\end{enumerate}
\item Dans la phrase \zh{杯子被弟弟打破了}, l'agent est :
\begin{enumerate}
\item \zh{杯子} \item \zh{打破} \item \zh{弟弟}
\end{enumerate}
\item Quelle phrase est correcte ?
\begin{enumerate}
\item \zh{门被开了。} (sens : on a ouvert la porte)
\item \zh{他被老师批评了。}
\item \zh{他被没老师批评。}
\end{enumerate}
\end{enumerate}
\end{exercice}

\begin{exercice}[Vrai ou faux ? Justifie ta réponse]
Dis si chaque affirmation est vraie ou fausse, puis justifie en une phrase en français.
\begin{enumerate}
\item Dans une phrase en \zh{被}, le sujet subit l'action.
\item On peut toujours supprimer l'agent après \zh{被}.
\item La négation de la phrase passive se fait avec \zh{没(有)} placé avant \zh{被}.
\item \zh{被} peut être remplacé par \zh{叫} ou \zh{让} dans un style familier.
\item Le verbe d'une phrase en \zh{被} peut rester seul, sans complément ni \zh{了}.
\end{enumerate}
\end{exercice}

\begin{exercice}[Vocabulaire : les verbes de la passive]
Complète le tableau en donnant le pinyin et la traduction de chaque verbe souvent employé dans une phrase en \zh{被}.
\begin{center}
\begin{tabularx}{\linewidth}{|X|X|X|}\hline
\rowcolor{popblueL} Caractères & Pinyin & Traduction \\ \hline
\zh{偷} & & \\ \hline
\zh{打破} & & \\ \hline
\zh{批评} & & \\ \hline
\zh{借走} & & \\ \hline
\zh{吃完} & & \\ \hline
\zh{咬伤} & & \\ \hline
\end{tabularx}
\end{center}
\end{exercice}

\begin{exercice}[Pinyin et caractères : associe]
Associe chaque expression en caractères à son pinyin.
\begin{center}
\begin{tabularx}{\linewidth}{|X|X|}\hline
\rowcolor{popblueL} Caractères & Pinyin \\ \hline
1. \zh{被雨淋湿了} & a. bèi lǎoshī biǎoyáng le \\ \hline
2. \zh{被老师表扬了} & b. bèi gǒu yǎo le \\ \hline
3. \zh{被狗咬了} & c. bèi yǔ lín shī le \\ \hline
4. \zh{被风吹倒了} & d. bèi fēng chuī dǎo le \\ \hline
\end{tabularx}
\end{center}
\end{exercice}

\courssub{Je m'entraîne}

\begin{exercice}[Traduction : du français au chinois]
Traduis les phrases suivantes en chinois (caractères), en utilisant une phrase en \zh{被}.
\begin{enumerate}
\item Mon portefeuille a été volé au marché.
\item Il n'a pas été puni par le professeur.
\item La lettre a déjà été écrite.
\item Le ballon a été perdu par notre équipe.
\end{enumerate}
\end{exercice}

\begin{exercice}[Transformation : de l'actif au passif]
Réécris chaque phrase active en phrase passive avec \zh{被}. Écris la phrase complète en caractères.
\begin{enumerate}
\item \zh{小狗咬破了孩子的鞋。} (Xiǎo gǒu yǎo pò le háizi de xié.)
\item \zh{弟弟吃完了我的饭。} (Dìdi chī wán le wǒ de fàn.)
\item \zh{大风刮倒了大树。} (Dà fēng guā dǎo le dà shù.)
\item \zh{同学借走了我的词典。} (Tóngxué jiè zǒu le wǒ de cídiǎn.)
\item \zh{妈妈洗干净了衣服。} (Māma xǐ gānjìng le yīfu.)
\end{enumerate}
\end{exercice}

\begin{exercice}[Texte à trous : les mots-outils]
Complète le texte suivant avec les mots-outils qui conviennent : \zh{被}, \zh{了}, \zh{没}, \zh{得}.
\begin{textebox}[Un mauvais jour]
昨天我很倒霉。我的手机 \_\_\_\_ 人偷走 \_\_\_\_，我的作业 \_\_\_\_ 雨淋湿了，我还 \_\_\_\_ 老师批评 \_\_\_\_。晚上回家，我发现门 \_\_\_\_ 锁上，钥匙也找不到了。唉，真倒霉！\\[4pt]
\emph{Zuótiān wǒ hěn dǎoméi. Wǒ de shǒujī \_\_\_\_ rén tōu zǒu \_\_\_\_, wǒ de zuòyè \_\_\_\_ yǔ lín shī le, wǒ hái \_\_\_\_ lǎoshī pīpíng \_\_\_\_. Wǎnshang huí jiā, wǒ fāxiàn mén \_\_\_\_ suǒ shàng, yàoshi yě zhǎo bu dào le. Ai, zhēn dǎoméi!}\\[4pt]
« Hier, j'ai vraiment eu de la malchance. Mon téléphone a été volé, mes devoirs ont été trempés par la pluie, et j'ai aussi été critiqué par le professeur. Le soir, en rentrant, j'ai découvert que la porte était fermée à clé et je ne trouvais plus mes clés. Quelle malchance ! »
\end{textebox}
\end{exercice}

\begin{exercice}[Remise en ordre]
Remets les mots dans le bon ordre pour former une phrase passive correcte, puis traduis-la en français.
\begin{enumerate}
\item \zh{批评 / 被 / 他 / 了 / 老师}
\item \zh{打破了 / 杯子 / 妹妹 / 被}
\item \zh{被 / 我的书 / 借走了 / 同学}
\item \zh{没 / 被 / 他 / 发现 / 大家}
\end{enumerate}
\end{exercice}

\courssub{Je me prépare au Bac}

\begin{exercice}[Type Bac — Compréhension de texte et questions de langue]
Lis le texte suivant, puis réponds aux questions.
\begin{textebox}[Un incident au marché de Mokolo]
上星期六，阿里去莫科洛市场买东西。人很多，很热闹。他看水果的时候，钱包被人偷走了。他很着急，马上告诉了市场旁边的警察。警察帮他找，可是没有找到。更倒霉的是，他的自行车也被风吹倒了，车灯打破了。回到家，他把这件事告诉了妈妈。他没有被妈妈批评，妈妈只说：「下次小心一点！」\\[4pt]
\emph{Shàng ge xīngqīliù, Ālǐ qù Mòkēluò shìchǎng mǎi dōngxi. Rén hěn duō, hěn rènao. Tā kàn shuǐguǒ de shíhou, qiánbāo bèi rén tōu zǒu le. Tā hěn zháojí, mǎshàng gàosu le shìchǎng pángbiān de jǐngchá. Jǐngchá bāng tā zhǎo, kěshì méiyǒu zhǎodào. Gèng dǎoméi de shì, tā de zìxíngchē yě bèi fēng chuī dǎo le, chēdēng dǎ pò le. Huí dào jiā, tā bǎ zhè jiàn shì gàosu le māma. Tā méiyǒu bèi māma pīpíng, māma zhǐ shuō : « Xià cì xiǎoxīn yìdiǎn! »}\\[4pt]
« Samedi dernier, Ali est allé au marché de Mokolo faire des courses. Il y avait beaucoup de monde, c'était très animé. Pendant qu'il regardait les fruits, son portefeuille a été volé. Il était très inquiet et l'a tout de suite dit au policier près du marché. Le policier l'a aidé à chercher, mais sans succès. Comble de malchance, son vélo a aussi été renversé par le vent et le phare a été cassé. Rentré à la maison, il a raconté cela à sa mère. Il n'a pas été grondé par sa mère ; elle a seulement dit : "Fais plus attention la prochaine fois !" »
\end{textebox}
\textbf{A. Questions de compréhension} (réponds en français)
\begin{enumerate}
\item Où et quand l'histoire se passe-t-elle ?
\item Quels sont les deux incidents qui arrivent à Ali ? Qui ou quoi en est l'agent ?
\item Pourquoi la mère d'Ali ne le gronde-t-elle pas ?
\end{enumerate}
\textbf{B. Questions de langue}
\begin{enumerate}
\item Relève dans le texte trois phrases en \zh{被} et analyse-les (sujet patient, agent, verbe + complément).
\item Pourquoi dit-on \zh{没有被妈妈批评} et non \zh{被没有妈妈批评} ? Explique la règle.
\item Transforme en phrase passive : \zh{人偷走了他的钱包。}
\end{enumerate}
\end{exercice}

\begin{exercice}[Type Bac — Thème et correction d'erreurs]
\textbf{A. Thème.} Traduis en chinois (caractères) les phrases suivantes.
\begin{enumerate}
\item Mon portefeuille a été volé au marché.
\item Il n'a pas été puni par le professeur.
\item La lettre a déjà été écrite.
\item Les mangues ont toutes été mangées par les enfants.
\end{enumerate}
\textbf{B. Correction d'erreurs.} Chaque phrase contient une faute. Recopie la phrase corrigée et explique l'erreur en français.
\begin{enumerate}
\item \zh{他被没老师批评。}
\item \zh{杯子被打了。} (sens voulu : le verre a été cassé)
\item \zh{门是被开了。} (sens voulu : la porte a été ouverte)
\item \zh{我的书被同学借。}
\end{enumerate}
\end{exercice}

\begin{exercice}[Type Bac — Choix multiple et expression écrite]
\textbf{A. Choix multiple.} Complète chaque phrase avec le mot qui convient : \zh{被}, \zh{把}, \zh{叫} ou \zh{让}.
\begin{enumerate}
\item 我的词典 \_\_\_\_ 同桌借走了。
\item 他 \_\_\_\_ 门关上，就回家了。
\item 我的新鞋 \_\_\_\_ 小狗咬破了。(style familier)
\item 妈妈 \_\_\_\_ 我去市场买西红柿。
\item 教室 \_\_\_\_ 我们打扫得干干净净。
\end{enumerate}
\textbf{B. Expression écrite.} Rédige en chinois un court paragraphe de \textbf{60 à 80 caractères} intitulé \zh{倒霉的一天} (« Une journée de malchance »). Tu raconteras un incident (réel ou imaginaire) qui t'est arrivé à Yaoundé ou à Douala. Consignes obligatoires :
\begin{itemize}
\item utiliser \textbf{au moins deux phrases en \zh{被}} ;
\item utiliser au moins une fois la négation \zh{没有被} ;
\item soigner l'ordre des mots et la ponctuation chinoise 。
\end{itemize}
\end{exercice}

\newpage

\coursec{Corrigés des exercices}

\coursec{Corrigés détaillés — Leçon 18 : 被字句}

\begin{corrige}[Exercice 1 — QCM : la bonne structure]
\begin{enumerate}
\item Réponse \textbf{a} : Sujet (objet patient) + 被 + agent + verbe. C'est l'ordre fondamental de la \zh{被字句} : celui qui \emph{subit} l'action est placé en tête, l'agent suit \zh{被}.
\item Réponse \textbf{b} : « Mon vélo a été volé. » \zh{我的自行车} (sujet patient) + \zh{被} + agent sous-entendu + \zh{偷了}.
\item Réponse \textbf{c} : L'agent est \zh{弟弟} (le petit frère), placé juste après \zh{被}. \zh{杯子} est le sujet patient, \zh{打破} le verbe avec son complément de résultat.
\item Réponse \textbf{b} : \zh{他被老师批评了。} \emph{Tā bèi lǎoshī pīpíng le.} « Il a été critiqué par le professeur. » La phrase a est fautive car le verbe \zh{开} reste seul, sans complément ni résultat ; la phrase c est fautive car la négation \zh{没} doit se placer \emph{avant} \zh{被} : \zh{他没被老师批评。}
\end{enumerate}
\end{corrige}

\begin{corrige}[Exercice 2 — Vrai ou faux ? Justifie ta réponse]
\begin{enumerate}
\item \textbf{Vrai.} Dans une phrase en \zh{被}, le sujet est le \emph{patient} : il subit l'action exprimée par le verbe (ex. \zh{杯子被打破了} \emph{Bēizi bèi dǎ pò le} : c'est le verre qui subit la casse).
\item \textbf{Faux.} On ne peut supprimer l'agent que lorsqu'il est inconnu, évident ou sans importance (\zh{钱包被偷了。} \emph{Qiánbāo bèi tōu le.}). Si l'agent est précis et pertinent, il doit être mentionné : \zh{他被老师批评了。} \emph{Tā bèi lǎoshī pīpíng le.}
\item \textbf{Vrai.} La négation se fait avec \zh{没(有)} placé \emph{avant} \zh{被} : \zh{他没有被老师批评。} \emph{Tā méiyǒu bèi lǎoshī pīpíng.} « Il n'a pas été critiqué par le professeur. »
\item \textbf{Vrai.} À l'oral et dans un style familier, \zh{叫} \emph{jiào} et \zh{让} \emph{ràng} peuvent remplacer \zh{被} : \zh{我的鞋叫狗咬破了。} \emph{Wǒ de xié jiào gǒu yǎo pò le.} = \zh{我的鞋被狗咬破了。} (mais jamais dans un écrit formel).
\item \textbf{Faux.} Le verbe d'une \zh{被字句} ne peut presque jamais rester seul : il doit être suivi de \zh{了}, d'un complément de résultat (\zh{打破}, \zh{吃完}) ou de direction (\zh{借走}). On dit \zh{杯子被打破了}, jamais \zh{杯子被打} seul.
\end{enumerate}
\end{corrige}

\begin{corrige}[Exercice 3 — Vocabulaire : les verbes de la passive]
\begin{center}
\begin{tabularx}{\linewidth}{|X|X|X|X|}\hline
\rowcolor{popblueL} Caractères & Pinyin & Nature & Traduction \\ \hline
\zh{偷} & tōu & Verbe & voler, dérober \\ \hline
\zh{打破} & dǎ pò & Verbe composé (V + R) & casser (verbe + résultat « brisé ») \\ \hline
\zh{批评} & pīpíng & Verbe & critiquer, réprimander \\ \hline
\zh{借走} & jiè zǒu & Verbe composé (V + D) & emprunter (et emporter) \\ \hline
\zh{吃完} & chī wán & Verbe composé (V + R) & finir de manger, tout manger \\ \hline
\zh{咬伤} & yǎo shāng & Verbe composé (V + R) & blesser en mordant \\ \hline
\end{tabularx}
\end{center}
\medskip
\noindent \textbf{Point de vigilance :} \zh{打破}, \zh{借走}, \zh{吃完}, \zh{咬伤} sont des verbes à \cle{complément de résultat (R)} ou \cle{de direction (D)} ; c'est ce complément qui rend la phrase en \zh{被} grammaticalement complète. \zh{偷} et \zh{批评} acceptent \zh{了} seul pour marquer l'accompli.
\end{corrige}

\begin{corrige}[Exercice 4 — Pinyin et caractères : associe]
\begin{enumerate}
\item \textbf{c} : \zh{被雨淋湿了} \emph{bèi yǔ lín shī le} « avoir été trempé par la pluie ».
\item \textbf{a} : \zh{被老师表扬了} \emph{bèi lǎoshī biǎoyáng le} « avoir été félicité par le professeur ».
\item \textbf{b} : \zh{被狗咬了} \emph{bèi gǒu yǎo le} « avoir été mordu par un chien ».
\item \textbf{d} : \zh{被风吹倒了} \emph{bèi fēng chuī dǎo le} « avoir été renversé par le vent ».
\end{enumerate}
\end{corrige}

\begin{corrige}[Exercice 5 — Traduction : du français au chinois]
\begin{enumerate}
\item \zh{我的钱包在市场被偷了。} \emph{Wǒ de qiánbāo zài shìchǎng bèi tōu le.} — Le complément de lieu \zh{在市场} se place \emph{avant} \zh{被}.
\item \zh{他没有被老师惩罚。} \emph{Tā méiyǒu bèi lǎoshī chéngfá.} — Négation : \zh{没有} \emph{avant} \zh{被}.
\item \zh{信已经被写好了。} \emph{Xìn yǐjīng bèi xiě hǎo le.} — \zh{已经} « déjà » se place \emph{avant} \zh{被} ; \zh{写好} = verbe + résultat « terminé ».
\item \zh{球被我们队丢了。} \emph{Qiú bèi wǒmen duì diū le.} — L'agent \zh{我们队} « notre équipe » suit \zh{被}.
\end{enumerate}
\end{corrige}

\begin{corrige}[Exercice 6 — Transformation : de l'actif au passif]
\textbf{Méthode :} l'objet de la phrase active devient sujet ; l'ancien sujet devient agent après \zh{被} ; on conserve le verbe avec son complément et \zh{了}.
\begin{enumerate}
\item \zh{孩子的鞋被小狗咬破了。} \emph{Háizi de xié bèi xiǎo gǒu yǎo pò le.} « La chaussure de l'enfant a été trouée par le petit chien. »
\item \zh{我的饭被弟弟吃完了。} \emph{Wǒ de fàn bèi dìdi chī wán le.} « Mon repas a été entièrement mangé par mon petit frère. »
\item \zh{大树被大风刮倒了。} \emph{Dà shù bèi dà fēng guā dǎo le.} « Le grand arbre a été renversé par le grand vent. »
\item \zh{我的词典被同学借走了。} \emph{Wǒ de cídiǎn bèi tóngxué jiè zǒu le.} « Mon dictionnaire a été emprunté (emporté) par un camarade. »
\item \zh{衣服被妈妈洗干净了。} \emph{Yīfu bèi māma xǐ gānjìng le.} « Les vêtements ont été lavés bien propres par maman. »
\end{enumerate}
\medskip
\noindent \textbf{Point de vigilance :} ne jamais perdre le complément (\zh{破}, \zh{完}, \zh{倒}, \zh{走}, \zh{干净}) ni \zh{了} dans la transformation.
\end{corrige}

\begin{corrige}[Exercice 7 — Texte à trous : les mots-outils]
Texte complété :\\[4pt]
\zh{昨天我很倒霉。我的手机被人偷走了，我的作业被雨淋湿了，我还被老师批评了。晚上回家，我发现门没锁上，钥匙也找不到了。唉，真倒霉！}\\[4pt]
\emph{Zuótiān wǒ hěn dǎoméi. Wǒ de shǒujī bèi rén tōu zǒu le, wǒ de zuòyè bèi yǔ lín shī le, wǒ hái bèi lǎoshī pīpíng le. Wǎnshang huí jiā, wǒ fāxiàn mén méi suǒ shàng, yàoshi yě zhǎo bu dào le. Ai, zhēn dǎoméi!}
\begin{itemize}
\item \zh{被} (trois fois) : marque du passif devant les agents \zh{人}, \zh{雨}, \zh{老师}.
\item \zh{了} (trois fois dans les phrases en \zh{被}) : après \zh{偷走}, \zh{淋湿} et \zh{批评}, il marque l'accomplissement de l'action subie.
\item \zh{没} : négation de \zh{锁上} « la porte n'était pas fermée à clé » (action non réalisée au passé).
\item \zh{得} n'est pas utilisé ici : c'était un distracteur ; il sert dans les compléments de manière (\zh{打扫得干干净净}), absents de ce texte.
\end{itemize}
\end{corrige}

\begin{corrige}[Exercice 8 — Remise en ordre]
\begin{enumerate}
\item \zh{他被老师批评了。} \emph{Tā bèi lǎoshī pīpíng le.} « Il a été critiqué par le professeur. »
\item \zh{杯子被妹妹打破了。} \emph{Bēizi bèi mèimei dǎ pò le.} « Le verre a été cassé par (ma) petite sœur. »
\item \zh{我的书被同学借走了。} \emph{Wǒ de shū bèi tóngxué jiè zǒu le.} « Mon livre a été emprunté par un camarade. »
\item \zh{他没被大家发现。} \emph{Tā méi bèi dàjiā fāxiàn.} « Il n'a pas été découvert par tout le monde. » — \zh{没} se place avant \zh{被}.
\end{enumerate}
Rappel du schéma : Sujet patient + (\zh{没}) + \zh{被} + agent + verbe + complément / \zh{了}.
\end{corrige}

\begin{corrige}[Exercice 9 — Type Bac : compréhension de texte et questions de langue]
\textbf{Barème indicatif (sur 10 pts)} : A. compréhension 4 pts (1 + 2 + 1) ; B. langue 6 pts (3 + 2 + 1).

\textbf{A. Questions de compréhension}
\begin{enumerate}
\item L'histoire se passe au marché de Mokolo (à Yaoundé), samedi dernier (\zh{上星期六}).
\item Deux incidents : d'abord, son portefeuille a été volé — l'agent est \zh{人} « quelqu'un » (un voleur inconnu) ; ensuite, son vélo a été renversé et le phare cassé — l'agent est \zh{风} « le vent ».
\item Sa mère ne le gronde pas parce qu'elle est compréhensive : elle se contente de lui conseiller d'être plus prudent la prochaine fois (\zh{下次小心一点}).
\end{enumerate}
\textbf{B. Questions de langue}
\begin{enumerate}
\item Trois phrases en \zh{被} du texte :
\begin{itemize}
\item \zh{钱包被人偷走了。} \emph{Qiánbāo bèi rén tōu zǒu le.} — sujet patient : \zh{钱包} ; agent : \zh{人} ; verbe + complément : \zh{偷走} + \zh{了}.
\item \zh{他的自行车也被风吹倒了。} \emph{Tā de zìxíngchē yě bèi fēng chuī dǎo le.} — sujet patient : \zh{自行车} ; agent : \zh{风} ; verbe + complément : \zh{吹倒} + \zh{了}.
\item \zh{他没有被妈妈批评。} \emph{Tā méiyǒu bèi māma pīpíng.} — sujet patient : \zh{他} ; agent : \zh{妈妈} ; verbe : \zh{批评}, avec négation \zh{没有} avant \zh{被}.
\end{itemize}
\item On dit \zh{没有被妈妈批评} parce que la négation \zh{没(有)} se place obligatoirement \emph{avant} \zh{被} : c'est toute la construction passive qui est niée, pas le verbe seul. \zh{被没有妈妈批评} est impossible.
\item \zh{他的钱包被人偷走了。} \emph{Tā de qiánbāo bèi rén tōu zǒu le.} « Son portefeuille a été volé. »
\end{enumerate}
\textbf{Points de vigilance} : bien identifier l'agent même quand il est vague (\zh{人}) ; vérifier la place de la négation ; conserver \zh{了} en fin de verbe.
\end{corrige}

\begin{corrige}[Exercice 10 — Type Bac : thème et correction d'erreurs]
\textbf{Barème indicatif (sur 10 pts)} : A. thème 5 pts (1,25 pt par phrase) ; B. correction 5 pts (0,75 pt la phrase corrigée + 0,5 pt l'explication).

\textbf{A. Thème}
\begin{enumerate}
\item \zh{我的钱包在市场被偷了。} \emph{Wǒ de qiánbāo zài shìchǎng bèi tōu le.}
\item \zh{他没有被老师惩罚。} \emph{Tā méiyǒu bèi lǎoshī chéngfá.}
\item \zh{信已经被写好了。} \emph{Xìn yǐjīng bèi xiě hǎo le.}
\item \zh{芒果都被孩子们吃完了。} \emph{Mángguǒ dōu bèi háizimen chī wán le.} — \zh{都} « toutes » se place avant \zh{被}.
\end{enumerate}
\textbf{B. Correction d'erreurs}
\begin{enumerate}
\item Faute : \zh{他被没老师批评。} → Correction : \zh{他没有被老师批评。} La négation \zh{没有} doit précéder \zh{被}, jamais le suivre.
\item Faute : \zh{杯子被打了。} → Correction : \zh{杯子被打破了。} Le verbe \zh{打} seul ne peut pas exprimer le résultat « cassé » ; il faut le complément de résultat \zh{破}.
\item Faute : \zh{门是被开了。} → Correction : \zh{门被打开了。} \zh{是} ne s'emploie pas dans la \zh{被字句} simple (calque du français « a été ») ; le verbe doit porter un complément (\zh{打开} « ouvrir » avec résultat).
\item Faute : \zh{我的书被同学借。} → Correction : \zh{我的书被同学借走了。} Le verbe ne peut rester nu : il faut le complément de direction \zh{走} et la particule \zh{了}.
\end{enumerate}
\textbf{Points de vigilance} : trois pièges classiques des francophones — place de la négation, verbe nu sans complément, ajout parasite de \zh{是} calqué sur le français « a été ».
\end{corrige}

\begin{corrige}[Exercice 11 — Type Bac : choix multiple et expression écrite]
\textbf{Barème indicatif (sur 10 pts)} : A. choix multiple 2,5 pts (0,5 pt par réponse) ; B. expression écrite 7,5 pts (respect des consignes 3 pts : deux \zh{被} + une négation \zh{没有被} ; correction grammaticale 2,5 pts ; vocabulaire et ponctuation 2 pts).

\textbf{A. Choix multiple}
\begin{enumerate}
\item 我的词典 \textbf{被} 同桌借走了。 — passif neutre/écrit.
\item 他 \textbf{把} 门关上，就回家了。 — construction en \zh{把} : il agit \emph{sur} la porte (actif disposant).
\item 我的新鞋 \textbf{叫} 小狗咬破了。 — style familier : \zh{叫} remplace \zh{被} (accepté aussi : \zh{让}).
\item 妈妈 \textbf{让} 我去市场买西红柿。 — ici \zh{让} = « faire faire, demander à » (causatif), non passif.
\item 教室 \textbf{被} 我们打扫得干干净净。 — passif avec complément de manière \zh{得干干净净}.
\end{enumerate}
\textbf{B. Expression écrite — proposition de rédaction modèle}
\medskip
\noindent \textbf{Texte chinois :}\par
\zh{上星期一，我在雅温得遇到了倒霉的一天。早上，我的雨伞被风刮坏了，所以我的衣服被雨淋湿了。到了学校，我发现我的作业被弟弟画花了。幸好，我没有被老师批评，老师只说：「明天再交吧！」下午回家，我的钱也丢了。唉，真是倒霉的一天！}\par
\medskip
\noindent \textbf{Pinyin :}\par
\emph{Shàng ge xīngqīyī, wǒ zài Yǎwēndé yùdào le dǎoméi de yì tiān. Zǎoshang, wǒ de yǔsǎn bèi fēng guā huài le, suǒyǐ wǒ de yīfu bèi yǔ lín shī le. Dào le xuéxiào, wǒ fāxiàn wǒ de zuòyè bèi dìdi huà huā le. Xìnghǎo, wǒ méiyǒu bèi lǎoshī pīpíng, lǎoshī zhǐ shuō : « Míngtiān zài jiāo ba! » Xiàwǔ huí jiā, wǒ de qián yě diū le. Ai, zhēn shì dǎoméi de yì tiān!}\par
\medskip
\noindent \textbf{Traduction :}\par
« Lundi dernier, j'ai vécu une journée de malchance à Yaoundé. Le matin, mon parapluie a été abîmé par le vent, donc mes vêtements ont été trempés par la pluie. Arrivé à l'école, j'ai découvert que mes devoirs avaient été gribouillés par mon petit frère. Heureusement, je n'ai pas été critiqué par le professeur ; il a seulement dit : "Tu rendras demain !" L'après-midi, en rentrant, j'ai aussi perdu mon argent. Quelle journée de malchance ! »
\medskip
\noindent Ce modèle compte environ 95 caractères ; l'élève peut viser 60 à 80 caractères en supprimant un incident. Il contient \textbf{quatre} phrases en \zh{被} (\zh{被风刮坏了}, \zh{被雨淋湿了}, \zh{被弟弟画花了}, \zh{没有被老师批评}) dont une négation \zh{没有被}.

\textbf{Points de vigilance pour la rédaction} :
\begin{itemize}
\item Vérifier que chaque verbe de la \zh{被字句} porte un complément ou \zh{了} (\zh{刮坏了}, \zh{淋湿了}, \zh{画花了}).
\item Placer la négation \zh{没有} avant \zh{被} : \zh{我没有被老师批评}, jamais l'inverse.
\item Employer la ponctuation chinoise pleine chasse ：，。！ et les guillemets 「」 ou « 》 correctement.
\item Ne pas calquer le français « j'ai été » par \zh{我是} : le passif chinois n'utilise pas \zh{是}.
\end{itemize}
\end{corrige}

\newpage

\coursec{Fiche bilan}

\begin{popfigure}
\begin{tikzpicture}[mindmap, grow cyclic, every node/.style={concept, font=\small},
  root concept/.append style={concept color=popblue, font=\large\bfseries, text=white},
  level 1/.append style={level distance=3.4cm, sibling angle=51.4},
  level 1 concept/.append style={font=\small, text=white}]
\node[root concept, align=center]{被字句\\ \footnotesize phrase passive}
  child[concept color=poporange]{ node[align=center]{Schéma\\ \footnotesize 受事 + 被 + 施事 + V + 其他} }
  child[concept color=popgreen]{ node[align=center]{Emplois\\ \footnotesize subir, subir un effet} }
  child[concept color=poppurple]{ node[align=center]{Variantes\\ \footnotesize 叫 / 让 / 给} }
  child[concept color=poppink]{ node[align=center]{Négation\\ \footnotesize 没(有) + 被} }
  child[concept color=popgold]{ node[align=center]{Compléments\\ \footnotesize 了 / 过 / 得 + degré} }
  child[concept color=popblue!70]{ node[align=center]{Pièges\\ \footnotesize ordre, verbe seul interdit} }
  child[concept color=popgreen!70]{ node[align=center]{Français\\ \footnotesize passif ≠ 被 toujours} };
\end{tikzpicture}
\legende{Carte mentale de la leçon : la phrase en 被 (bèi zì jù)}
\end{popfigure}

\begin{bilan}
\textbf{1. Lexique clé (à connaître par cœur)}
\begin{center}
\begin{tabularx}{\linewidth}{|X|X|X|X|}
\hline
\rowcolor{popblueL} Caractères & Pinyin & Nature & Traduction \\
\hline
被 & bèi & préposition & par (marque du passif) \\
\hline
叫 & jiào & préposition & par (familier) \\
\hline
让 & ràng & préposition & par (familier) \\
\hline
给 & gěi & préposition / verbe & par (oral) ; donner \\
\hline
受事 / 施事 & shòushì / shīshì & n. & patient / agent \\
\hline
打破 & dǎpò & v. & casser, briser \\
\hline
偷走 & tōuzǒu & v. & voler (et emporter) \\
\hline
批评 & pīpíng & v. / n. & critiquer ; critique \\
\hline
表扬 & biǎoyáng & v. & féliciter, louer \\
\hline
耽误 & dānwu & v. & retarder, faire perdre (du temps) \\
\hline
\end{tabularx}
\end{center}

\textbf{2. Règles de grammaire à connaître par cœur}
\begin{center}
\begin{tabularx}{\linewidth}{|X|X|X|}
\hline
\rowcolor{popblueL} Structure & Exemple (caractères + pinyin) & Traduction \\
\hline
受事 + 被 + 施事 + V + 其他 & 杯子被弟弟打破了。Bēizi bèi dìdi dǎpò le. & Le verre a été cassé par le petit frère. \\
\hline
受事 + 被 + V + 其他 (agent omis) & 钱包被偷走了。Qiánbāo bèi tōuzǒu le. & Le portefeuille a été volé. \\
\hline
受事 + 没(有) + 被 + 施事 + V & 作业没有被老师收走。Zuòyè méiyǒu bèi lǎoshī shōuzǒu. & Les devoirs n'ont pas été ramassés par le professeur. \\
\hline
受事 + 叫/让/给 + 施事 + V + 其他 & 衣服叫雨淋湿了。Yīfu jiào yǔ línshī le. & Les vêtements ont été mouillés par la pluie. \\
\hline
被 + V + 得 + complément de degré & 他被老师批评得很厉害。Tā bèi lǎoshī pīpíng de hěn lìhai. & Il a été sévèrement critiqué par le professeur. \\
\hline
\end{tabularx}
\end{center}

\begin{itemize}
\item Le verbe de la phrase en 被 ne peut \textbf{jamais} rester seul : il doit être suivi de 了, 过, d'un complément de résultat (打破, 偷走), de direction ou de degré.
\item La négation se place \textbf{avant} 被 : 没有被 / 没被 (jamais 被没有).
\item L'agent peut être omis quand il est inconnu ou sans importance : 自行车被偷了。
\item 被 marque souvent un événement \textbf{subi} (nuance défavorable), mais l'usage moderne accepte aussi le neutre ou le favorable : 他被选为班长。(Il a été élu délégué.)
\item Variantes familières : 叫 et 让 exigent un agent exprimé ; 给 est très oral.
\end{itemize}

\textbf{3. Méthodes express}
\begin{enumerate}
\item Identifier le \cle{patient} (celui qui subit) : il devient le sujet, placé avant 被.
\item Placer \cle{被} puis l'agent (ou rien si l'agent est omis).
\item Choisir un verbe \textbf{accompagné} : V + 了 / V-résultat / V + 得 + degré.
\item Vérifier la place de la négation (avant 被) et des compléments de temps (avant 被 aussi : 昨天被…).
\item Pour transformer actif → passif : inverser patient et agent, insérer 被, garder le reste inchangé.
\end{enumerate}
\end{bilan}

\begin{aretenir}[Checklist avant le Bac]
\begin{itemize}
\item[$\square$] Je sais écrire et prononcer 被 (bèi) avec le bon ton.
\item[$\square$] Je connais le schéma : 受事 + 被 + 施事 + V + 其他.
\item[$\square$] Je sais que le verbe ne peut pas être seul après 被 (il faut 了, 过 ou un complément).
\item[$\square$] Je place correctement la négation : 没(有) avant 被.
\item[$\square$] Je sais omettre l'agent quand il est inconnu : 门被打开了。
\item[$\square$] Je connais les variantes familières 叫, 让, 给 et leurs contraintes.
\item[$\square$] Je sais transformer une phrase active en phrase passive en 被, et inversement.
\item[$\square$] Je sais que le passif français ne se traduit pas toujours par 被 (parfois 受到, 获得, ou une tournure active).
\item[$\square$] Je place les compléments de temps et de lieu avant 被 : 昨天在学校被…
\item[$\square$] Je sais utiliser 被 avec un complément de degré : 被…批评得很厉害.
\item[$\square$] Je vérifie l'ordre des mots et les tons dans mes productions écrites.
\item[$\square$] Je sais repérer la nuance « subir » de 被 et choisir la tournure adaptée au contexte.
\end{itemize}
\end{aretenir}

\findecours

\end{document}
